% Dangers du courant électrique et sécurité — PCT 3ème — Cours signé OnBuch+
% Programme officiel MINESEC. Compiler avec : tectonic N23-dangers-courant-securite.tex
\newcommand{\DOCMATIERE}{PCT}
\newcommand{\DOCNIVEAU}{3ème}
\newcommand{\DOCMODULE}{Module 4 : Technologie et mécanique}
\newcommand{\DOCLECON}{N23}
\newcommand{\DOCTITRE}{Dangers du courant électrique et sécurité}
% OnBuch+ — préambule « cours signés » ENRICHI (compilé avec tectonic / XeLaTeX)
% Système visuel « Pop » d'OnBuch+ : fond crème (jamais blanc), bordures encre
% épaisses, ombres « dures » décalées, titres Archivo Black en capitales.
% Version MATHÉMATIQUES : graphes (pgfplots/tikz), boîtes pédagogiques
% (définition, propriété, méthode, attention, le savais-tu, exercice résolu,
% exercice, TP, bilan), page de garde et signature OnBuch+ ancrée en bas.
%
% Macros attendues dans main.tex AVANT \input{preamble} :
%   \DOCMATIERE  \DOCNIVEAU  \DOCMODULE  \DOCLECON (numéro)  \DOCTITRE
\documentclass[11pt]{article}

\usepackage{fontspec}
\usepackage[french]{babel}
\usepackage[a4paper,margin=2cm,top=3.4cm,bottom=2.8cm,headheight=1.4cm,footskip=1.3cm]{geometry}
\usepackage{amsmath,amssymb,amsfonts}
\usepackage{mathtools} % \xmapsto, \coloneqq... souvent utilisés par les modèles
\usepackage{cancel} % simplifications barrées (bilans de forces)
% \abs{z} (module, valeur absolue) et \norm{u} (norme), utilisés par les modèles
\providecommand{\abs}{}\let\abs\relax\DeclarePairedDelimiter{\abs}{\lvert}{\rvert}
\providecommand{\norm}{}\let\norm\relax\DeclarePairedDelimiter{\norm}{\lVert}{\rVert}
\usepackage[table]{xcolor} % [table] : fournit aussi \rowcolors (alternance de lignes)
\usepackage{pagecolor}
\usepackage{tikz}
\usetikzlibrary{arrows.meta,positioning,calc,shapes.geometric,shapes.misc,shapes.arrows,decorations.pathmorphing,decorations.pathreplacing,decorations.markings,patterns,shadows,mindmap,trees,fit,backgrounds,matrix,intersections,angles,quotes}
\usepackage{pgfplots}
\usepackage[version=4]{mhchem} % équations nucléaires \ce{^{235}_{92}U}
\usepackage[european,siunitx]{circuitikz} % schémas électriques
\pgfplotsset{compat=1.18}
\usepgfplotslibrary{fillbetween,groupplots} % aires entre courbes (calcul intégral)
\usepackage{enumitem}
\usepackage{fancyhdr}
% [most] charge toutes les bibliothèques tcolorbox (listings, minted, raster,
% documentation...) et gonfle inutilement la mémoire TeX sur un document long
% (~300 boîtes) : on ne charge que ce qui est réellement utilisé.
\usepackage[skins,breakable]{tcolorbox}
\usepackage{graphicx}
\usepackage{colortbl}
\usepackage{booktabs}
\usepackage{tabularx}
\usepackage{multirow}
\usepackage{diagbox} % case d'en-tête diagonale des tableaux à double entrée
% Colonnes extensibles centrée / gauche / droite (C, L, R), souvent utilisées par les modèles
\newcolumntype{C}{>{\centering\arraybackslash}X}
\newcolumntype{L}{>{\raggedright\arraybackslash}X}
\newcolumntype{R}{>{\raggedleft\arraybackslash}X}
\usepackage{array}
\usepackage{multicol}
\usepackage{siunitx}
% parse-numbers=false : accepte aussi \SI{2,5 \times 10^{-3}}{...}
\sisetup{locale=FR,output-decimal-marker={,},group-minimum-digits=4,parse-numbers=false}
\DeclareSIUnit\ohm{\ensuremath{\symup{\Omega}}} % U+2126 absent de la police
\DeclareSIUnit\division{div} % réglages d'oscilloscope (ms/div, V/div)
\DeclareSIUnit\tr{tr} % tours (vitesse de rotation, ex. \SI{1500}{\tr\per\minute})
\DeclareSIUnit\molar{M} % concentration molaire (ex. \SI{3}{\molar}, \SI{1}{\micro\molar})
\DeclareSIUnit\an{an} % année (le modèle confond parfois avec \year, un compteur TeX) — ex. \SI{5}{\kilo\meter\per\an}
\DeclareSIUnit\FCFA{FCFA} % franc CFA (ex. \SI{500}{\FCFA\per\kilo\gram})
\DeclareSIUnit\degreeBrix{\textdegree Brix} % teneur en sucre d'un sirop/jus (réfractométrie)
\DeclareSIUnit\month{mois} % le modèle utilise parfois \month, un compteur TeX, comme unité de durée
\DeclareSIUnit\microliter{\micro\liter} % le modèle écrit parfois \microliter en un seul token
\DeclareSIUnit\million{millions} % dénombrement (ex. \SI{15}{\million\per\mL} spermatozoïdes)
\usepackage{qrcode}
% \degree : commande fréquemment utilisée par les modèles pour un angle ou une
% température en dehors de \SI{}{} (non fournie par les paquets chargés).
\providecommand{\degree}{^{\circ}}
% \qed (fin de démonstration), utilisé par les modèles sans amsthm
\providecommand{\qed}{\hfill$\square$}
% \barre : double barre des valeurs interdites dans un tableau de variations (tkz-tab)
\providecommand{\barre}{\ensuremath{\|}}
% environnement proof (amsthm non chargé) utilisé par les modèles
\ifdefined\proof\else\newenvironment{proof}[1][Démonstration]{\par\noindent\textit{#1.}\ }{\qed\par}\fi
\usepackage{lastpage}
\usepackage{hyperref}
\hypersetup{hidelinks}

% ============================================================
% Palette « Pop » OnBuch+ — mêmes hex que la classe P (plus_ui.dart)
% ============================================================
\definecolor{poporange}{HTML}{F59321}
\definecolor{popdark}{HTML}{E07A0C}
\definecolor{popcream}{HTML}{FFF6E8}
\definecolor{popink}{HTML}{1C1714}
\definecolor{poppurple}{HTML}{7A5AE0}
\definecolor{popgreen}{HTML}{2E9E6A}
\definecolor{poppink}{HTML}{E0457B}
\definecolor{popblue}{HTML}{2F7DE1}
\definecolor{popgold}{HTML}{C98A12}
\definecolor{popmuted}{HTML}{8A7F76}
\definecolor{popline}{HTML}{EADFD2}
\definecolor{popgray}{HTML}{8A7F76}
\colorlet{popgrey}{popgray}
% Alias tolérants (le modèle nomme parfois les couleurs différemment)
\colorlet{popwhite}{white}
\colorlet{popblack}{popink}
\colorlet{popred}{poppink}
\colorlet{popyellow}{popgold}
% Teintes claires (fonds de tableaux/encadrés) — le modèle nomme parfois
% les couleurs avec un suffixe L (ex. popblueL) sans les avoir définies.
\colorlet{poporangeL}{poporange!20}
\colorlet{popdarkL}{popdark!20}
\colorlet{poppurpleL}{poppurple!20}
\colorlet{popgreenL}{popgreen!20}
\colorlet{poppinkL}{poppink!20}
\colorlet{popblueL}{popblue!20}
\colorlet{popgoldL}{popgold!20}
\colorlet{popredL}{popred!20}
\colorlet{popgrayL}{popgray!20}
% Teintes claires pour fonds de tableaux / zones
\colorlet{popblueL}{popblue!10!white}
\colorlet{poporangeL}{poporange!14!white}
\colorlet{popgreenL}{popgreen!12!white}
\colorlet{poppurpleL}{poppurple!12!white}
\colorlet{poppinkL}{poppink!12!white}
\colorlet{popgoldL}{popgold!14!white}

\pagecolor{popcream}

% ============================================================
% Typographie
% ============================================================
\defaultfontfeatures{Ligatures=TeX}
\setmainfont{PlusJakartaSans-Medium.ttf}[Path=../../../fonts/,BoldFont=PlusJakartaSans-Bold.ttf]
% Maths en sans-serif (Fira Math) avec chiffres et lettres droites de Plus
% Jakarta, pour que les formules se fondent dans le texte.
\usepackage[math-style=ISO,bold-style=ISO]{unicode-math}
\setmathfont{FiraMath-Regular.otf}[Path=../../../fonts/]
\setmathfont{PlusJakartaSans-Medium.ttf}[Path=../../../fonts/,range={up/num,up/{Latin,latin},\mathup}]
\usepackage{icomma} % virgule décimale sans espace en mode math
% Caractères mathématiques Unicode absents des polices de texte (Plus Jakarta,
% Archivo Black) : rendus par leur équivalent mathématique, en texte comme en
% formule — sinon ils s'affichent en carré vide (ex. « Arithmétique dans ℤ »).
\usepackage{newunicodechar}
\newunicodechar{ℝ}{\ensuremath{\mathbb{R}}}
\newunicodechar{ℤ}{\ensuremath{\mathbb{Z}}}
\newunicodechar{ℂ}{\ensuremath{\mathbb{C}}}
\newunicodechar{ℕ}{\ensuremath{\mathbb{N}}}
\newunicodechar{ℚ}{\ensuremath{\mathbb{Q}}}
\newunicodechar{𝔻}{\ensuremath{\mathbb{D}}}
\newunicodechar{α}{\ensuremath{\alpha}}
\newunicodechar{β}{\ensuremath{\beta}}
\newunicodechar{γ}{\ensuremath{\gamma}}
\newunicodechar{δ}{\ensuremath{\delta}}
\newunicodechar{ε}{\ensuremath{\varepsilon}}
\newunicodechar{θ}{\ensuremath{\theta}}
\newunicodechar{λ}{\ensuremath{\lambda}}
\newunicodechar{μ}{\ensuremath{\mu}}
\newunicodechar{σ}{\ensuremath{\sigma}}
\newunicodechar{τ}{\ensuremath{\tau}}
\newunicodechar{φ}{\ensuremath{\varphi}}
\newunicodechar{ψ}{\ensuremath{\psi}}
\newunicodechar{ω}{\ensuremath{\omega}}
\newunicodechar{Σ}{\ensuremath{\Sigma}}
% majuscules produites par \MakeUppercase dans les titres (α-aminés -> Α)
\newunicodechar{Α}{\ensuremath{\alpha}}
\newunicodechar{Β}{\ensuremath{\beta}}
\newunicodechar{Γ}{\ensuremath{\Gamma}}
\newunicodechar{Δ}{\ensuremath{\Delta}}
\newunicodechar{Ε}{\ensuremath{\varepsilon}}
\newunicodechar{Θ}{\ensuremath{\Theta}}
\newunicodechar{Λ}{\ensuremath{\Lambda}}
\newunicodechar{Μ}{\ensuremath{\mu}}
\newunicodechar{Τ}{\ensuremath{\tau}}
\newunicodechar{Φ}{\ensuremath{\Phi}}
\newunicodechar{Ψ}{\ensuremath{\Psi}}
\newunicodechar{Ω}{\ensuremath{\Omega}}
\newunicodechar{↦}{\ensuremath{\mapsto}}
\newunicodechar{∈}{\ensuremath{\in}}
\newunicodechar{∉}{\ensuremath{\notin}}
\newunicodechar{∧}{\ensuremath{\wedge}}
\newunicodechar{⇔}{\ensuremath{\Leftrightarrow}}
\newunicodechar{⟺}{\ensuremath{\Longleftrightarrow}}
\newunicodechar{⇒}{\ensuremath{\Rightarrow}}
\newunicodechar{⟹}{\ensuremath{\Longrightarrow}}
\newunicodechar{∘}{\ensuremath{\circ}}
\newunicodechar{∪}{\ensuremath{\cup}}
\newunicodechar{∩}{\ensuremath{\cap}}
\newunicodechar{≡}{\ensuremath{\equiv}}
\newunicodechar{⊂}{\ensuremath{\subset}}
\newunicodechar{⊥}{\ensuremath{\perp}}
\newunicodechar{∃}{\ensuremath{\exists}}
\newunicodechar{∀}{\ensuremath{\forall}}
\newunicodechar{∛}{\ensuremath{\sqrt[3]{\,}}}
\newunicodechar{∜}{\ensuremath{\sqrt[4]{\,}}}
\newunicodechar{⃗}{\ensuremath{{}^{\to}}}
\newunicodechar{ⁿ}{\ifmmode^{n}\else\textsuperscript{\ensuremath{n}}\fi}
\newunicodechar{ˣ}{\ifmmode^{x}\else\textsuperscript{\ensuremath{x}}\fi}
\newunicodechar{ᵇ}{\ifmmode^{b}\else\textsuperscript{\ensuremath{b}}\fi}
\newunicodechar{ʳ}{\ifmmode^{r}\else\textsuperscript{\ensuremath{r}}\fi}
\newunicodechar{ᵘ}{\ifmmode^{u}\else\textsuperscript{\ensuremath{u}}\fi}
\newunicodechar{ʸ}{\ifmmode^{y}\else\textsuperscript{\ensuremath{y}}\fi}
\newunicodechar{ᵃ}{\ifmmode^{a}\else\textsuperscript{\ensuremath{a}}\fi}
\newunicodechar{ᵉ}{\ifmmode^{e}\else\textsuperscript{\ensuremath{e}}\fi}
\newunicodechar{ᵏ}{\ifmmode^{k}\else\textsuperscript{\ensuremath{k}}\fi}
\newunicodechar{ᵖ}{\ifmmode^{p}\else\textsuperscript{\ensuremath{p}}\fi}
\newunicodechar{ᴺ}{\ifmmode^{N}\else\textsuperscript{\ensuremath{N}}\fi}
\newunicodechar{ᶜ}{\ifmmode^{c}\else\textsuperscript{\ensuremath{c}}\fi}
\newunicodechar{ʰ}{\ifmmode^{h}\else\textsuperscript{\ensuremath{h}}\fi}
\newunicodechar{ᵠ}{\ifmmode^{q}\else\textsuperscript{\ensuremath{q}}\fi}
\newunicodechar{ₙ}{\ifmmode_{n}\else\textsubscript{\ensuremath{n}}\fi}
\newunicodechar{ₐ}{\ifmmode_{a}\else\textsubscript{\ensuremath{a}}\fi}
\newunicodechar{ᵢ}{\ifmmode_{i}\else\textsubscript{\ensuremath{i}}\fi}
\newunicodechar{ᵣ}{\ifmmode_{r}\else\textsubscript{\ensuremath{r}}\fi}
\newunicodechar{ₖ}{\ifmmode_{k}\else\textsubscript{\ensuremath{k}}\fi}
\newunicodechar{ᵦ}{\ifmmode_{\beta}\else\textsubscript{\ensuremath{\beta}}\fi}
\newunicodechar{ₓ}{\ifmmode_{x}\else\textsubscript{\ensuremath{x}}\fi}
\newfontfamily\popdisplay{ArchivoBlack-Regular.ttf}[Path=../../../fonts/]
\newfontfamily\poplogo{SpaceGrotesk-Bold.ttf}[Path=../../../fonts/]
\newfontfamily\popsemi{PlusJakartaSans-SemiBold.ttf}[Path=../../../fonts/]
\newfontfamily\popxbold{PlusJakartaSans-ExtraBold.ttf}[Path=../../../fonts/]
\newfontfamily\popxboldls{PlusJakartaSans-ExtraBold.ttf}[Path=../../../fonts/,LetterSpace=3.0]

\setlist[enumerate]{leftmargin=1.7em,itemsep=5pt,topsep=4pt}
\setlist[itemize]{leftmargin=1.7em,itemsep=4pt,topsep=4pt,label={\color{poporange}\textbullet}}
\setlength{\parindent}{0pt}
\setlength{\parskip}{5pt}
\renewcommand{\arraystretch}{1.25}

% pgfplots : style Pop par défaut
\pgfplotsset{
  every axis/.append style={axis line style={popink,line width=1pt},tick style={popink},
    grid style={popline,line width=0.5pt},label style={font=\small},tick label style={font=\footnotesize}},
  popcurve/.style={poporange,line width=1.8pt,no markers,smooth},
}
% Même style disponible dans un \draw TikZ simple (hors pgfplots)
\tikzset{popcurve/.style={poporange,line width=1.8pt}}
\tikzset{popgrid/.style={popline,very thin}}
\colorlet{popcurve}{poporange} % le nom du style est parfois utilisé comme couleur

% ============================================================
% Pill / squiggle
% ============================================================
\newcommand{\poppill}[3]{%
  \tikz[baseline=-0.55ex]\node[draw=popink,line width=1.2pt,rounded corners=2.6pt,%
    fill=#2,inner xsep=6.5pt,inner ysep=3pt]{%
    \popxboldls\fontsize{7}{7}\selectfont\color{#3}\MakeUppercase{#1}};%
}
\newcommand{\popsquiggle}[1]{%
  \tikz[baseline=-0.4ex]{
    \draw[#1,line width=2.6pt,line cap=round]
      plot [smooth] coordinates {(0,0.09) (0.34,0.01) (0.68,0.21) (1.02,0.01) (1.36,0.10)};
  }%
}
% Mot-clé surligné (marqueur orange sous le texte)
\newcommand{\cle}[1]{{\popxbold\color{popdark}#1}}

% ============================================================
% En-tête / pied
% ============================================================
\pagestyle{fancy}
\fancyhf{}
\renewcommand{\headrulewidth}{0pt}
\renewcommand{\footrulewidth}{0pt}
\lhead{\raisebox{-0.15ex}{\poplogo\fontsize{13}{13}\selectfont\color{popink}OnBuch\color{poporange}+}}
\rhead{\poppill{\DOCMATIERE\ \textbullet\ \DOCNIVEAU\ \textbullet\ Leçon \DOCLECON}{white}{popink}}
\lfoot{\popxbold\fontsize{7.6}{7.6}\selectfont\color{popmuted}\MakeUppercase{Cours signé OnBuch+}\ \textbullet\ {\popsemi\DOCTITRE}}
\rfoot{\tikz[baseline=-0.6ex]\node[rounded corners=8.5pt,fill=popink,minimum height=17pt,minimum width=17pt,inner xsep=5pt,inner sep=0]{\popxbold\fontsize{8}{8}\selectfont\color{popcream}\thepage\,/\,\pageref*{LastPage}};}

% ============================================================
% Titres
% ============================================================
\newcommand{\chaptitre}[2]{%
  \begin{tcolorbox}[enhanced,colback=poporange,colframe=popink,boxrule=2.6pt,arc=11pt,
      drop shadow=popink,left=15pt,right=15pt,top=12pt,bottom=11pt,before skip=3pt,after skip=18pt]
    \poppill{#2}{popink}{popcream}\par\vspace{9pt}
    {\raggedright\hyphenpenalty=10000\exhyphenpenalty=10000\popdisplay\fontsize{22}{24}\selectfont\color{popcream}\MakeUppercase{#1}\par}
  \end{tcolorbox}
}
% Section numérotée : pastille orange avec le numéro + titre Archivo
\newcounter{popsec}
\newcommand{\coursec}[1]{%
  \stepcounter{popsec}\setcounter{popsub}{0}%
  \par\vspace{16pt}\noindent
  \begin{minipage}{\linewidth}
  \tikz[baseline=-0.7ex]\node[draw=popink,line width=1.6pt,fill=poporange,rounded corners=4pt,minimum size=22pt,inner sep=0]{\popdisplay\fontsize{12}{12}\selectfont\color{popcream}\thepopsec};%
  \hspace{8pt}{\popdisplay\fontsize{15}{17}\selectfont\color{popink}\MakeUppercase{#1}}\par
  \vspace{4pt}\noindent{\color{poporange}\rule{36pt}{3.6pt}}
  \end{minipage}\par\nopagebreak\vspace{8pt}
}
\newcounter{popsub}
\newcommand{\courssub}[1]{%
  \stepcounter{popsub}%
  \par\vspace{9pt}\noindent{\popxbold\fontsize{11.5}{13}\selectfont\color{popink}{\color{poporange}\thepopsec.\thepopsub}\hspace{6pt}#1}\par\nopagebreak\vspace{4pt}
}

% ============================================================
% Boîtes pédagogiques — même recette : fond blanc, bordure épaisse colorée,
% ombre dure encre, badge plein. Titre optionnel [#1].
% ============================================================
\tcbset{popbase/.style={breakable,enhanced,colback=white,boxrule=2.3pt,arc=9pt,
  shadow={3pt}{-3pt}{0pt}{popink},left=14pt,right=14pt,top=11pt,bottom=11pt,before skip=11pt,after skip=15pt,
  fontupper=\popsemi\fontsize{10.6}{14.5}\selectfont\color{popink},
  fontlower=\popsemi\fontsize{10.6}{14.5}\selectfont\color{popink}}}
% Coche dessinée (le glyphe ✓ n'existe pas dans les polices du document)
\newcommand{\popcheck}[1]{\tikz[baseline=-0.5ex]\draw[#1,line width=1.6pt,line cap=round,line join=round](-0.13,0.0)--(-0.04,-0.09)--(0.13,0.1);}
\providecommand{\coche}{\popcheck{popgreen}}
\providecommand{\resultat}[1]{\par\smallskip\noindent\fcolorbox{popgreen}{popgreenL}{\parbox{\dimexpr\linewidth-2\fboxsep-2\fboxrule\relax}{\textbf{Résultat.} #1}}\par\smallskip}
\newcommand{\popboxhead}[3]{\poppill{#1}{#2}{white}\ifx\relax#3\relax\else\hspace{6pt}{\popxbold\fontsize{10.6}{12}\selectfont\color{popink}#3}\fi\par\vspace{7pt}}

\newtcolorbox{aretenir}[1][]{popbase,colframe=popgreen,before upper={\popboxhead{À retenir}{popgreen}{#1}}}
\newtcolorbox{exemplebox}[1][]{popbase,colframe=poporange,before upper={\popboxhead{Exemple}{poporange}{#1}}}
\newtcolorbox{definition}[1][]{popbase,colframe=popblue,before upper={\popboxhead{Définition}{popblue}{#1}}}
\newtcolorbox{propriete}[1][]{popbase,colframe=poppurple,before upper={\popboxhead{Propriété}{poppurple}{#1}}}
\newtcolorbox{methode}[1][]{popbase,colframe=popgold,before upper={\popboxhead{Méthode}{popgold}{#1}}}
\newtcolorbox{attention}[1][]{popbase,colframe=poppink,before upper={\popboxhead{Attention}{poppink}{#1}}}
\newtcolorbox{experience}[1][]{popbase,colframe=popblue,colback=popblueL,before upper={\popboxhead{Expérience}{popblue}{#1}}}
% « Le savais-tu ? » : carte violette pleine (contexte, culture, Cameroun)
\newtcolorbox{savaistu}[1][]{popbase,colback=poppurple,colframe=popink,
  fontupper=\popsemi\fontsize{10.4}{14.2}\selectfont\color{white},
  before upper={\poppill{Le savais-tu\,?}{white}{poppurple}\ifx\relax#1\relax\else\hspace{6pt}{\popxbold\color{white}#1}\fi\par\vspace{7pt}}}
% Exercice résolu : énoncé en haut, \tcblower puis la solution sur fond crème
\newcounter{popexo}\newcounter{popexr}
\newtcolorbox{exoresolu}[1][]{popbase,colframe=poporange,colbacklower=popcream,
  segmentation style={popink,line width=1.2pt,dashed},
  before upper={\stepcounter{popexr}\popboxhead{Exercice résolu \thepopexr}{poporange}{#1}},
  before lower={\poppill{Solution}{popink}{popcream}\par\vspace{6pt}}}
% Exercice (non résolu en place) — la correction est dans la partie Corrigés
\newtcolorbox{exercice}[1][]{popbase,colframe=popink,
  before upper={\stepcounter{popexo}\popboxhead{Exercice \thepopexo}{popink}{#1}}}
% Corrigé
\newtcolorbox{corrige}[1][]{popbase,colframe=popgreen,colback=popgreenL,
  before upper={\popboxhead{Corrigé}{popgreen}{#1}}}
% Objectifs / prérequis / bilan
\newtcolorbox{objectifs}{popbase,colframe=popink,colback=poporangeL,
  before upper={\popboxhead{Objectifs de la leçon}{popink}{}}}
\newtcolorbox{prerequis}{popbase,colframe=popink,colback=popblueL,
  before upper={\popboxhead{Prérequis}{popblue}{}}}
\newtcolorbox{bilan}{popbase,colframe=popink,colback=white,boxrule=2.8pt,
  before upper={\popboxhead{Fiche bilan}{poporange}{L'essentiel à connaître pour l'examen}}}
% Figure encadrée avec légende
\newtcolorbox{popfigure}[1][]{enhanced,breakable=false,colback=white,colframe=popline,boxrule=1.4pt,arc=8pt,
  left=10pt,right=10pt,top=10pt,bottom=8pt,before skip=10pt,after skip=12pt,halign=center,
  lower separated=false,halign lower=center,
  fontlower=\popsemi\fontsize{9}{11.5}\selectfont\color{popmuted},
  #1}
\newcommand{\legende}[1]{\par\vspace{6pt}{\popsemi\fontsize{9}{11.5}\selectfont\color{popmuted}#1}}

% ============================================================
% Signature OnBuch+
% ============================================================
\newcommand{\poplogoinline}[1]{{\poplogo\fontsize{#1}{#1}\selectfont\color{popink}OnBuch\color{poporange}+}}
\newcommand{\signatureonbuch}{%
  \begin{tcolorbox}[enhanced,colback=popink,colframe=popink,boxrule=2.2pt,arc=10pt,
      drop shadow=poporange,left=14pt,right=14pt,top=10pt,bottom=10pt,before skip=0pt,after skip=0pt]
    \tikz[baseline=-0.7ex]\node[circle,fill=popgreen,draw=popcream,line width=1.3pt,minimum size=22pt,inner sep=0]{\popcheck{white}};%
    \hspace{9pt}\begin{minipage}[c]{0.62\linewidth}
      {\popxbold\fontsize{12}{13}\selectfont\color{popcream}Signé\ \textbullet\ {\poplogo OnBuch\color{poporange}+}}\par\vspace{2pt}
      {\raggedright\popsemi\fontsize{8}{10.5}\selectfont\color{popline}\MakeUppercase{Équipe pédagogique OnBuch+}\\\MakeUppercase{Conforme au programme officiel MINESEC}\par}
    \end{minipage}\hfill
    {\poplogo\fontsize{22}{22}\selectfont\color{popcream}OnBuch\color{poporange}+}
  \end{tcolorbox}%
}

% ============================================================
% Page de garde
% #1 titre · #2 matière · #3 niveau · #4 module · #5 n° leçon · #6 durée ·
% #7 description / objectifs (texte ou itemize)
% ============================================================
\newcommand{\pagegarde}[7]{%
  \thispagestyle{empty}
  \newgeometry{a4paper,margin=1.7cm,top=1.6cm,bottom=1.4cm}%
  \begin{tikzpicture}[remember picture,overlay]
    \fill[poporange!18] ([xshift=-3.2cm,yshift=-3.6cm]current page.north east) circle (4.6cm);
    \fill[poppurple!14] ([xshift=2.4cm,yshift=7.6cm]current page.south west) circle (3.8cm);
  \end{tikzpicture}%
  {\poplogo\fontsize{30}{30}\selectfont\color{popink}OnBuch\color{poporange}+}\hfill
  \raisebox{0.6ex}{\poppill{Cours signé \textbullet\ Premium}{popink}{popcream}}\par
  \vspace{1pt}\popsquiggle{poppurple}\par
  \vspace{8pt}
  \begin{tcolorbox}[enhanced,colback=poporange,colframe=popink,boxrule=2.8pt,arc=13pt,
      drop shadow=popink,left=18pt,right=18pt,top=12pt,bottom=13pt,after skip=10pt]
    \poppill{#2 \textbullet\ #3}{popink}{popcream}\hspace{5pt}\poppill{#4}{white}{popink}\par\vspace{8pt}
    {\popdisplay\fontsize{14}{14}\selectfont\color{popink}LEÇON #5}\par\vspace{4pt}
    {\raggedright\hyphenpenalty=10000\exhyphenpenalty=10000\popdisplay\fontsize{25}{27}\selectfont\color{popcream}\MakeUppercase{#1}\par}
    \vspace{8pt}
    {\popxbold\fontsize{9.5}{9.5}\selectfont\color{popink}\MakeUppercase{Durée conseillée : #6}}
  \end{tcolorbox}
  \begin{tcolorbox}[enhanced,colback=white,colframe=popink,boxrule=2.2pt,arc=10pt,
      drop shadow=popink,left=16pt,right=16pt,top=10pt,bottom=10pt,after skip=8pt]
    \poppill{Dans cette leçon}{poppurple}{white}\par\vspace{5pt}
    \popsemi\fontsize{9.3}{12.6}\selectfont\color{popink}#7
  \end{tcolorbox}
  \noindent\begin{minipage}[t]{0.487\linewidth}
    \begin{tcolorbox}[enhanced,colback=popgreen,colframe=popink,boxrule=2.2pt,arc=10pt,
        drop shadow=popink,left=11pt,right=11pt,top=10pt,bottom=10pt,height=3.1cm,valign=center]
      \tcbox[colback=white,colframe=popink,boxrule=1.3pt,arc=5pt,boxsep=0pt,nobeforeafter,
        left=3pt,right=3pt,top=3pt,bottom=3pt,valign=center]{\qrcode[height=1.9cm]{https://play.google.com/store/apps/details?id=com.luvvix.onbuch&hl=fr}}%
      \hspace{8pt}\begin{minipage}[c]{0.5\linewidth}
        {\popdisplay\fontsize{11}{12.5}\selectfont\color{white}\MakeUppercase{Télécharge OnBuch}}\par\vspace{4pt}
        {\raggedright\popsemi\fontsize{8.4}{11}\selectfont\color{white}Gratuit \textbullet\ Google Play\\Tuteur IA, annales, concours, résultats\par}
      \end{minipage}
    \end{tcolorbox}
  \end{minipage}\hfill
  \begin{minipage}[t]{0.487\linewidth}
    \begin{tcolorbox}[enhanced,colback=poppurple,colframe=popink,boxrule=2.2pt,arc=10pt,
        drop shadow=popink,left=11pt,right=11pt,top=10pt,bottom=10pt,height=3.1cm,valign=center]
      \tikz[baseline=-0.7ex]\node[circle,fill=white,draw=popink,line width=1.3pt,minimum size=1.6cm,inner sep=0]{\popdisplay\fontsize{24}{24}\selectfont\color{poppurple}+};%
      \hspace{8pt}\begin{minipage}[c]{0.6\linewidth}
        {\popdisplay\fontsize{11}{12.5}\selectfont\color{white}\MakeUppercase{Passe à OnBuch+}}\par\vspace{4pt}
        {\raggedright\popsemi\fontsize{8.4}{11}\selectfont\color{white}Tous les cours signés, Tuteur IA illimité, fiches PDF — onglet OnBuch+ de l'app\par}
      \end{minipage}
    \end{tcolorbox}
  \end{minipage}
  \vspace{8pt}
  \popcontenu
  \vfill
  \signatureonbuch
  \restoregeometry
  \newpage
}

% Rangée « Ce cours contient » (page de garde)
\newcommand{\poptile}[3]{% #1 couleur #2 gros texte #3 légende
  \begin{tcolorbox}[enhanced,colback=#1,colframe=popink,boxrule=2pt,arc=9pt,shadow={2.5pt}{-2.5pt}{0pt}{popink},
      left=6pt,right=6pt,top=9pt,bottom=9pt,halign=center,valign=center,height=1.75cm,width=\linewidth,nobeforeafter]
    {\popdisplay\fontsize{12.5}{13}\selectfont\color{white}\MakeUppercase{#2}}\par\vspace{3pt}
    {\centering\popsemi\fontsize{7.8}{9.5}\selectfont\color{white}#3\par}
  \end{tcolorbox}}
\newcommand{\popcontenu}{%
  \par\noindent{\popxboldls\fontsize{7.5}{7.5}\selectfont\color{popmuted}CE COURS SIGNÉ CONTIENT}\par\vspace{7pt}
  \noindent\begin{minipage}[t]{0.235\linewidth}\poptile{popblue}{Cours}{complet, illustré, pas à pas}\end{minipage}\hfill
  \begin{minipage}[t]{0.235\linewidth}\poptile{poporange}{Méthodes}{et exercices résolus}\end{minipage}\hfill
  \begin{minipage}[t]{0.235\linewidth}\poptile{poppink}{Exercices}{type examen + corrigés}\end{minipage}\hfill
  \begin{minipage}[t]{0.235\linewidth}\poptile{popgreen}{Fiche bilan}{l'essentiel à retenir}\end{minipage}\par}

% Sommaire
\newtcolorbox{sommaire}{enhanced,colback=white,colframe=popink,boxrule=2.2pt,arc=10pt,
  drop shadow=popink,left=18pt,right=18pt,top=13pt,bottom=13pt,before skip=6pt,after skip=20pt,
  before upper={\poppill{Sommaire}{poppurple}{white}\par\vspace{9pt}\popsemi\fontsize{10.6}{14.5}\selectfont\color{popink}}}

% Signature de fin de cours — poussée tout en bas de la dernière page
\newcommand{\findecours}{%
  \par\vfill\nopagebreak
  \begin{center}\popsquiggle{poporange}\end{center}\vspace{4pt}
  \signatureonbuch
}
\AtBeginDocument{\def\llbracket{\mathopen{[\mkern-3.5mu[}}\def\rrbracket{\mathclose{]\mkern-3.5mu]}}} % intervalles d'entiers (glyphe ⟦ absent de la police)
% Glyphes absents de Fira Math : redéfinis à partir de glyphes disponibles
\AtBeginDocument{%
  \def\setminus{\mathbin{\backslash}}\let\smallsetminus\setminus
  \def\vdots{\mathord{\vbox{\baselineskip3.2pt\lineskiplimit0pt\kern4pt\hbox{.}\hbox{.}\hbox{.}}}}%
  \def\ddots{\mathinner{\mkern1mu\raise6.5pt\hbox{.}\mkern2mu\raise3.6pt\hbox{.}\mkern2mu\raise0.7pt\hbox{.}\mkern1mu}}%
  \def\triangle{\mathord{\scalebox{0.9}{$\symup{\Delta}$}}}%
  \def\nabla{\mathord{\raisebox{\depth}{\scalebox{1}[-1]{$\symup{\Delta}$}}}}%
  \def\checkmark{\popcheck{popdark}}%
  \def\star{\mathbin{\ast}}\def\bigstar{\mathord{\ast}}%
  \def\diamond{\mathbin{\scalebox{0.6}{$\Diamond$}}}%
  \def\bigcup{\mathop{\mathchoice{\raisebox{-0.3em}{\scalebox{1.7}{$\cup$}}}{\scalebox{1.25}{$\cup$}}{\cup}{\cup}}\displaylimits}%
  \def\bigcap{\mathop{\mathchoice{\raisebox{-0.3em}{\scalebox{1.7}{$\cap$}}}{\scalebox{1.25}{$\cap$}}{\cap}{\cap}}\displaylimits}%
  \def\lesseqqgtr{\mathrel{\lessgtr}}\def\bigoplus{\mathop{\oplus}\displaylimits}%
}
\providecommand{\ssi}{si et seulement si} % abréviation fréquente
\AtBeginDocument{\providecommand{\wideparen}{\overparen}} % arc de cercle

%% FIGFIX-BEGIN v4 — correctifs globaux des figures (docs/onbuch-generation/tools/figfix)
\usepackage{adjustbox}\usepackage{etoolbox}
% toute figure TikZ/pgfplots plus large que la ligne est réduite à la largeur disponible
\BeforeBeginEnvironment{tikzpicture}{\begin{adjustbox}{max width=\linewidth}}
\AfterEndEnvironment{tikzpicture}{\end{adjustbox}}
% légendes pgfplots lisibles par-dessus les courbes
\pgfplotsset{every axis/.append style={legend style={fill=white,fill opacity=0.92,text opacity=1,draw=black!15,inner sep=3pt}}}
% étiquettes d'arêtes (syntaxe edge["..."]) avec fond blanc
\tikzset{every edge quotes/.append style={fill=white,inner sep=1.2pt}}
%% FIGFIX-END
\begin{document}

\pagegarde{Dangers du courant électrique et sécurité}{PCT}{3ème}{Module 4 : Technologie et mécanique}{N23}{2 séances (1 h chacune)}{Cette leçon étudie les dangers du courant électrique pour le corps humain (électrisation, électrocution) et les risques pour les installations (surcharge, court-circuit). Elle présente les dispositifs de sécurité (fusible, disjoncteur, prise de terre, différentiel) et les règles de comportement à adopter face au courant domestique.
\vspace{4pt}
\begin{multicols}{2}\small
\begin{itemize}
\item À la fin de cette leçon, je sais définir l'électrisation et l'électrocution et les distinguer.
\item À la fin de cette leçon, je sais expliquer, à partir de résultats expérimentaux, comment le courant traverse le corps humain et pourquoi il est dangereux.
\item À la fin de cette leçon, je sais identifier les facteurs qui aggravent le danger : intensité, durée, trajet du courant, humidité.
\item À la fin de cette leçon, je sais décrire les risques pour une installation : surcharge et court-circuit.
\item À la fin de cette leçon, je sais citer et expliquer le rôle des dispositifs de sécurité : fusible, disjoncteur, prise de terre, interrupteur différentiel.
\item À la fin de cette leçon, je sais interpréter le schéma simple d'une installation domestique protégée.
\end{itemize}\end{multicols}}

\begin{sommaire}
\begin{multicols}{2}
\begin{enumerate}
\item Le corps humain et le courant électrique
\item Les dangers du courant : effets sur l'organisme
\item Les risques pour les installations : surcharge et court-circuit
\item Les dispositifs de protection : fusible et disjoncteur
\item La prise de terre et l'interrupteur différentiel
\item Règles de sécurité et conduite à tenir
\item Activité d'intégration
\item Méthodes et exercices résolus
\item Exercices
\item Corrigés des exercices
\item Fiche bilan
\end{enumerate}
\end{multicols}
\end{sommaire}

\begin{objectifs}
\begin{itemize}
\item À la fin de cette leçon, je sais définir l'électrisation et l'électrocution et les distinguer.
\item À la fin de cette leçon, je sais expliquer, à partir de résultats expérimentaux, comment le courant traverse le corps humain et pourquoi il est dangereux.
\item À la fin de cette leçon, je sais identifier les facteurs qui aggravent le danger : intensité, durée, trajet du courant, humidité.
\item À la fin de cette leçon, je sais décrire les risques pour une installation : surcharge et court-circuit.
\item À la fin de cette leçon, je sais citer et expliquer le rôle des dispositifs de sécurité : fusible, disjoncteur, prise de terre, interrupteur différentiel.
\item À la fin de cette leçon, je sais interpréter le schéma simple d'une installation domestique protégée.
\item À la fin de cette leçon, je sais appliquer les consignes de sécurité dans des situations concrètes de la vie courante.
\item À la fin de cette leçon, je sais réaliser un calcul simple d'intensité (I = P/U) pour vérifier le calibre d'un fusible.
\end{itemize}
\end{objectifs}

\begin{prerequis}
\begin{itemize}
\item Notions de 4ème : circuit électrique simple, tension et intensité du courant
\item Loi d'Ohm (U = R × I) vue en 4ème
\item Puissance électrique P = U × I (leçon sur l'énergie électrique, 3ème)
\item Conducteurs et isolants (4ème)
\item Installation électrique domestique : phase, neutre (leçons précédentes du module)
\end{itemize}
\end{prerequis}

\newpage

\coursec{Le corps humain et le courant électrique}

\textbf{Situation d'accroche.} À Yaoundé, pendant la saison des pluies, un jeune élève rentre de l'école trempé. Sans s'essuyer les mains, il ouvre le frigo de la maison pour prendre de l'eau fraîche : aussitôt, il ressent une violente secousse dans le bras et lâche la poignée. Il vient d'être \emph{électrisé}. Chaque année au Cameroun, des accidents domestiques liés au courant électrique provoquent des brûlures, des arrêts cardiaques et des incendies. Pourtant, des dispositifs simples — fusible, disjoncteur, prise de terre, interrupteur différentiel — permettent de protéger les personnes et les installations. \textbf{Problématique :} pourquoi le courant électrique est-il dangereux pour le corps humain ? Comment s'en protéger efficacement à la maison ? Pour répondre à la première question, commençons par comprendre comment le courant interagit avec notre corps.

\courssub{Le corps humain est un conducteur électrique}

\textbf{Intuition.} Tu as peut-être déjà entendu : « Ne touche jamais une prise avec les mains mouillées ! ». Pourquoi les mains \emph{mouillées} ? Parce que l'eau joue un rôle central : notre corps est constitué d'environ \cle{60 à 70 \% d'eau}, et cette eau n'est pas pure. Elle contient des \cle{sels minéraux dissous} (sodium, potassium, calcium, chlorures...) sous forme d'\cle{ions}, c'est-à-dire de particules chargées électriquement. Or un courant électrique, c'est précisément un déplacement de porteurs de charges. Les ions de nos liquides corporels (sang, sueur, lymphe) peuvent donc se déplacer et transporter le courant : \textbf{le corps humain conduit le courant électrique}.

\begin{propriete}[Le corps humain, conducteur ionique]
Le corps humain est un \cle{conducteur électrique} : les liquides de l'organisme (sang, sueur, lymphe) contiennent de l'eau et des sels minéraux dissous sous forme d'\cle{ions}, qui permettent le passage du courant. La résistance électrique du corps dépend fortement de l'état de la peau : elle est élevée quand la peau est sèche, et chute brutalement quand la peau est mouillée ou transpirante.
\end{propriete}

La peau sèche est la seule barrière un peu protectrice : c'est une mauvaise conductrice. Mais dès qu'elle est humide (transpiration, pluie, baignoire), les ions de la sueur la rendent conductrice et le courant pénètre facilement dans les tissus, qui, eux, conduisent très bien.

\begin{experience}[Mesure de la résistance du corps humain]
\textbf{Matériel :} un ohmmètre (multimètre en position $\Omega$), deux électrodes métalliques tenues une dans chaque main, un peu d'eau salée.\\
\textbf{Protocole :}
\begin{enumerate}
\item Saisir une électrode dans chaque main, mains bien sèches, et lire la résistance affichée.
\item Recommencer après avoir humidifié les paumes avec de l'eau salée.
\item Comparer les deux mesures.
\end{enumerate}
\textbf{Observations :} avec la peau sèche, on mesure une résistance de l'ordre de \SI{100000}{\ohm} ; avec la peau mouillée, la résistance tombe en dessous de \SI{1000}{\ohm}.\\
\textbf{Interprétation :} l'eau salée comble les pores de la peau et crée des chemins conducteurs. La résistance du corps est donc divisée par plus de 100 quand la peau passe de sèche à mouillée.\\
\textbf{Conclusion :} le corps humain conduit le courant, et il le conduit d'autant mieux que la peau est humide. C'est pourquoi l'élève de notre situation d'accroche, les mains mouillées par la pluie, a ressenti une secousse si violente.
\end{experience}

\begin{attention}[Les conditions qui font chuter la résistance du corps]
La résistance du corps humain n'est pas une constante : elle diminue avec l'\cle{humidité} de la peau, la \cle{surface de contact}, la \cle{pression} du contact et la \cle{durée} du contact. Ne jamais écrire au BEPC que « le corps a une résistance de \SI{100000}{\ohm} » sans préciser l'état de la peau : c'est une erreur fréquente.
\end{attention}

\courssub{Électrisation et électrocution}

\begin{definition}[Électrisation]
L'\cle{électrisation} est le passage d'un courant électrique à travers le corps humain, provoquant des lésions : brûlures, tétanisation des muscles (contraction incontrôlable), troubles respiratoires, voire arrêt cardiaque.
\end{definition}

\begin{definition}[Électrocution]
L'\cle{électrocution} est une électrisation \textbf{mortelle} : le courant traversant le corps a provoqué la mort de la victime, le plus souvent par fibrillation cardiaque (le cœur bat de façon désordonnée et ne pompe plus le sang).
\end{definition}

Retenons la nuance : toute électrocution est une électrisation, mais toute électrisation n'est pas une électrocution. L'élève de Yaoundé a été électrisé ; heureusement, il n'a pas été électrocuté.

\courssub{C'est l'intensité qui est dangereuse, pas seulement la tension}

\textbf{Intuition.} On entend souvent : « Attention, c'est du 220 V, ça tue ! ». Cette phrase cache une confusion. La \cle{tension} $U$ est la \emph{cause} : c'est elle qui « pousse » les charges. Mais ce qui agit directement sur les organes — muscles, nerfs, cœur — c'est l'\cle{intensité} $I$ du courant qui traverse réellement le corps. Le lien entre les deux est donné par la \textbf{loi d'Ohm} :
\[ I = \frac{U}{R} \]
où $R$ est la résistance du corps dans les conditions du contact. À tension égale, plus la résistance du corps est faible (peau mouillée), plus l'intensité qui le traverse est grande, et plus le danger est grave.

\begin{attention}[Erreur fréquente : « c'est la tension qui tue »]
Non ! Ce n'est pas la tension qui tue, c'est \textbf{l'intensité du courant qui traverse le corps}. La tension n'en est que la cause, par l'intermédiaire de la loi d'Ohm $I = U/R$. Une clôture électrique de brousse peut délivrer plusieurs milliers de volts mais une intensité très faible et très brève : elle donne une secousse désagréable mais rarement mortelle. À l'inverse, le \SI{220}{V} du réseau ENEO peut maintenir dans le corps une intensité élevée pendant longtemps : il est mortel. Au BEPC, rédige toujours : « la tension provoque le passage d'une intensité dangereuse dans le corps ».
\end{attention}

Les médecins ont mesuré les \cle{seuils physiologiques} de l'action du courant alternatif (\SI{50}{Hz}, celui du réseau ENEO) sur l'organisme :

\begin{center}
\begin{tabularx}{\linewidth}{|l|X|}\hline
\rowcolor{popblueL} \textbf{Intensité traversant le corps} & \textbf{Effet sur l'organisme} \\ \hline
$I \approx \SI{0,5}{mA}$ & \cle{Seuil de perception} : léger picotement, on « sent » le courant. \\ \hline
$I \approx \SI{10}{mA}$ & \cle{Seuil de non lâcher} : tétanisation musculaire, les muscles se contractent, la victime ne peut plus lâcher le conducteur. \\ \hline
$I > \SI{30}{mA}$ & Risque de \cle{fibrillation cardiaque} : le cœur se met à trembler sans pomper le sang. \\ \hline
$I \approx \SI{100}{mA}$ & \cle{Seuil mortel} : arrêt cardiaque, brûlures internes graves. \\ \hline
\end{tabularx}
\end{center}

\begin{aretenir}[Les quatre seuils à connaître]
Perception : $\approx \SI{0,5}{mA}$ ; tétanisation (non lâcher) : $\approx \SI{10}{mA}$ ; fibrillation cardiaque possible : au-delà de $\approx \SI{30}{mA}$ ; seuil mortel : $\approx \SI{100}{mA}$. Ces valeurs sont des ordres de grandeur issus de mesures médicales : le danger croît avec l'intensité \textbf{et} avec la durée du contact.
\end{aretenir}

\begin{exemplebox}[Pourquoi le 220 V est-il mortel avec les mains mouillées ?]
Avec la peau mouillée, la résistance du corps peut descendre à $R = \SI{1000}{\ohm}$. Sous la tension domestique $U = \SI{220}{V}$, la loi d'Ohm donne :
\[ I = \frac{U}{R} = \frac{220}{1000} = \SI{0,22}{A} = \SI{220}{mA}. \]
Cette intensité est \textbf{largement supérieure au seuil mortel} (\SI{100}{mA}) : le contact avec un appareil défectueux, mains mouillées, peut tuer en quelques secondes. Avec la peau sèche ($R \approx \SI{100000}{\ohm}$), la même tension ne ferait passer que $I = 220/100000 = \SI{0,0022}{A} = \SI{2,2}{mA}$ : secousse douloureuse, mais rarement mortelle. Voilà pourquoi l'élève de notre accroche a eu très chaud... et pourquoi on ne touche jamais un appareil électrique les mains mouillées.
\end{exemplebox}

La courbe suivante résume la situation : elle montre l'intensité traversant le corps en fonction de la tension de contact, pour la peau sèche et pour la peau mouillée, avec les zones de danger.

\begin{popfigure}
\begin{center}
\begin{tikzpicture}
\begin{axis}[
width=0.95\linewidth, height=7.2cm,
xlabel={Tension de contact $U$ (V)},
ylabel={Intensité $I$ (mA)},
xmin=0, xmax=250, ymin=0, ymax=120,
xtick={0,50,100,150,200,250},
ytick={0.5,10,30,100},
yticklabels={{0,5},{10},{30},{100}},
grid=major, grid style={popline!40},
legend style={at={(0.03,0.97)}, anchor=north west, font=\small},
]
% Zones de danger
\addplot[draw=none, fill=popgreenL, opacity=0.5] coordinates {(0,0) (250,0) (250,0.5) (0,0.5)} \closedcycle;
\addplot[draw=none, fill=popgoldL, opacity=0.55] coordinates {(0,0.5) (250,0.5) (250,10) (0,10)} \closedcycle;
\addplot[draw=none, fill=poporangeL, opacity=0.55] coordinates {(0,10) (250,10) (250,30) (0,30)} \closedcycle;
\addplot[draw=none, fill=poppinkL, opacity=0.6] coordinates {(0,30) (250,30) (250,120) (0,120)} \closedcycle;
% Droites I = U/R (en mA)
\addplot[popblue, very thick, domain=0:250, samples=50] {1000*x/100000};
\addlegendentry{Peau sèche ($R \approx \SI{100000}{\ohm}$)}
\addplot[poporange, very thick, domain=0:120, samples=50] {1000*x/1000};
\addlegendentry{Peau mouillée ($R \approx \SI{1000}{\ohm}$)}
% Seuils
\addplot[popdark, dashed, domain=0:250] {0.5};
\addplot[popdark, dashed, domain=0:250] {10};
\addplot[popdark, dashed, domain=0:250] {30};
\addplot[popdark, dashed, domain=0:250] {100};
\node[font=\scriptsize, popdark, anchor=west] at (axis cs:5,3) {perception};
\node[font=\scriptsize, popdark, anchor=west] at (axis cs:5,18) {tétanisation};
\node[font=\scriptsize, popdark, anchor=west] at (axis cs:5,60) {fibrillation};
\node[font=\scriptsize, popdark, anchor=west] at (axis cs:5,110) {danger mortel};
% Point 220 V peau sèche
\addplot[only marks, mark=*, popblue, mark size=2pt] coordinates {(220,2.2)};
\node[font=\scriptsize, popblue, anchor=west] at (axis cs:120,5) {220 V, peau sèche : $I \approx 2,2$ mA};
% Annotation peau mouillée
\node[font=\scriptsize, poporange, anchor=south east, align=right] at (axis cs:245,118) {220 V, peau mouillée :\\ $I \approx 220$ mA (hors graphique, zone mortelle)};
\end{axis}
\end{tikzpicture}
\end{center}
\legende{Intensité traversant le corps humain en fonction de la tension de contact, d'après la loi d'Ohm $I = U/R$. Pour la peau mouillée (droite orange), l'intensité croît très vite : dès \SI{100}{V} elle atteint \SI{100}{mA}, le seuil mortel ; sous \SI{220}{V} elle vaut environ \SI{220}{mA}, très au-delà du seuil mortel. Pour la peau sèche (droite bleue), l'intensité reste faible même sous \SI{220}{V}. Les zones colorées indiquent les seuils physiologiques : perception, tétanisation, fibrillation, danger mortel.}
\end{popfigure}

\courssub{Le trajet du courant dans le corps : un facteur décisif}

Pour qu'un courant traverse le corps, il faut que celui-ci soit soumis à une \cle{différence de potentiel} entre deux points de contact : par exemple une main qui touche un fil dénudé et les pieds nus sur le sol mouillé. Le courant entre par un point et ressort par un autre. Or tous les trajets ne se valent pas :

\begin{itemize}
\item trajet \cle{main-main} : le courant traverse le thorax et passe \textbf{à travers le cœur} ;
\item trajet \cle{main-pied} : le courant descend le long du corps et traverse également la région du cœur ;
\item trajet doigt-doigt de la même main : le courant reste localisé, le cœur est épargné.
\end{itemize}

Les trajets main-main et main-pied sont les plus dangereux car le cœur, organe musculaire commandé par des signaux électriques naturels très faibles, est extrêmement sensible aux courants extérieurs.

\begin{popfigure}
\begin{center}
\begin{tikzpicture}[scale=0.9, font=\small]
% Corps stylisé
\draw[popdark, thick, fill=popcream] (0,2.6) circle (0.45); % tête
\draw[popdark, thick, fill=popcream, rounded corners=6pt] (-0.7,-0.4) rectangle (0.7,2.2); % tronc
% bras gauche levé vers phase
\draw[popdark, thick, fill=popcream, rounded corners=4pt] (-0.7,1.9) -- (-1.9,2.9) -- (-1.7,3.2) -- (-0.6,2.3) -- cycle;
% bras droit vers le bas
\draw[popdark, thick, fill=popcream, rounded corners=4pt] (0.7,1.9) -- (1.5,0.6) -- (1.2,0.4) -- (0.6,1.5) -- cycle;
% jambes
\draw[popdark, thick, fill=popcream, rounded corners=4pt] (-0.6,-0.4) rectangle (-0.1,-2.2);
\draw[popdark, thick, fill=popcream, rounded corners=4pt] (0.1,-0.4) rectangle (0.6,-2.2);
% coeur
\node[popink, font=\large] at (0.15,1.3) {$\heartsuit$};
\node[popink, font=\scriptsize, anchor=west] at (0.35,1.3) {c\oe ur};
% fil phase en haut à gauche
\draw[poporange, very thick] (-3.4,3.1) -- (-1.8,3.1);
\node[poporange, anchor=east, font=\scriptsize] at (-3.4,3.1) {fil sous tension (phase)};
% sol
\draw[popdark, thick] (-2.6,-2.2) -- (2.6,-2.2);
\foreach \x in {-2.4,-1.9,...,2.5} \draw[popdark] (\x,-2.2) -- (\x-0.25,-2.5);
\node[popdark, font=\scriptsize, anchor=north west] at (2.0,-2.5) {sol};
% flèche courant main-pied
\draw[popink, very thick, -{Stealth[length=3mm]}] (-1.75,2.95) .. controls (-1.0,2.2) and (-0.3,1.8) .. (0.1,1.35);
\draw[popink, very thick, -{Stealth[length=3mm]}] (0.15,1.2) .. controls (0.3,0.4) and (0.35,-0.8) .. (0.35,-2.1);
\node[popink, font=\scriptsize, anchor=west, align=left] at (0.6,-1.0) {trajet main-pied\\(traverse le c\oe ur)};
% étiquette électrisation
\node[popblue, font=\scriptsize, anchor=south] at (0,3.35) {Le corps relie la phase au sol : il est parcouru par un courant $I = U/R$.};
\end{tikzpicture}
\end{center}
\legende{Trajet du courant dans le corps humain. La victime touche un fil sous tension avec la main gauche pendant que ses pieds sont en contact avec le sol : le courant entre par la main, traverse le thorax (et le cœur, en rouge) et ressort par les pieds. Ce trajet main-pied, comme le trajet main-main, est l'un des plus dangereux car il fait passer le courant à travers le cœur.}
\end{popfigure}

\begin{propriete}[Facteurs aggravants d'une électrisation]
La gravité d'une électrisation dépend de quatre facteurs principaux :
\begin{enumerate}
\item l'\cle{intensité} $I$ qui traverse le corps (elle-même fixée par la tension de contact $U$ et la résistance $R$ du corps, via $I = U/R$) ;
\item la \cle{durée du contact} : plus le courant passe longtemps, plus les lésions sont graves (d'où le danger de la tétanisation, qui empêche de lâcher) ;
\item le \cle{trajet du courant} : les trajets main-main et main-pied traversent le cœur ;
\item l'\cle{état de la peau} : peau mouillée ou transpirante = résistance faible = intensité élevée.
\end{enumerate}
\end{propriete}

\begin{exoresolu}[Peau sèche ou peau mouillée : qui risque le plus ?]
\textbf{Énoncé.} Un technicien touche accidentellement un conducteur sous tension de \SI{220}{V}. a) Calculer l'intensité qui traverse son corps si sa peau est sèche ($R = \SI{100000}{\ohm}$). b) Même question si sa peau est mouillée ($R = \SI{1000}{\ohm}$). c) Dans chaque cas, dire à quel seuil physiologique on se situe et conclure.
\tcblower
\textbf{Solution.} On applique la loi d'Ohm $I = U/R$.
\begin{align*}
\text{a) Peau sèche : } I &= \frac{220}{100000} = \SI{0,0022}{A} = \SI{2,2}{mA}.\\
\text{b) Peau mouillée : } I &= \frac{220}{1000} = \SI{0,22}{A} = \SI{220}{mA}.
\end{align*}
c) Avec la peau sèche, $I = \SI{2,2}{mA}$ : on est au-dessus du seuil de perception (\SI{0,5}{mA}) mais en dessous du seuil de non lâcher (\SI{10}{mA}) : le technicien sent une secousse et peut se dégager. Avec la peau mouillée, $I = \SI{220}{mA}$, très au-dessus du seuil mortel (\SI{100}{mA}) : risque de fibrillation cardiaque et de mort. \textbf{Conclusion :} la même tension de \SI{220}{V} est supportable avec la peau sèche mais mortelle avec la peau mouillée, car c'est l'intensité traversant le corps qui détermine le danger.
\end{exoresolu}

\begin{savaistu}[Pourquoi les oiseaux ne sont-ils pas électrocutés sur les lignes haute tension ?]
À Yaoundé comme ailleurs, on voit des oiseaux posés sur les câbles électriques sans être incommodés, alors que ces lignes portent des milliers de volts ! Le secret : le courant ne traverse pratiquement pas leur corps. Leurs deux pattes touchent le \emph{même} fil, donc il n'existe \textbf{pas de différence de potentiel} significative entre elles : sans tension entre deux points du corps, pas de courant. En revanche, si l'oiseau touchait en même temps le câble et un poteau relié à la terre, son corps relierait deux points de potentiels très différents : il serait électrocuté instantanément. C'est exactement la situation de l'homme qui touche un fil dénudé les pieds au sol : son corps devient le pont entre la phase et la terre.
\end{savaistu}

\begin{aretenir}[L'essentiel de la section]
\begin{itemize}
\item Le corps humain est \cle{conducteur} grâce à l'eau salée (ions) de ses liquides ; sa résistance passe d'environ \SI{100000}{\ohm} (peau sèche) à moins de \SI{1000}{\ohm} (peau mouillée).
\item L'\cle{électrisation} est le passage d'un courant à travers le corps ; l'\cle{électrocution} est une électrisation mortelle.
\item C'est l'\cle{intensité} $I = U/R$ traversant le corps qui est dangereuse, pas la tension seule.
\item Seuils : perception $\approx \SI{0,5}{mA}$ ; non lâcher $\approx \SI{10}{mA}$ ; fibrillation au-delà de $\approx \SI{30}{mA}$ ; seuil mortel $\approx \SI{100}{mA}$.
\item Sous \SI{220}{V} avec la peau mouillée, $I \approx \SI{220}{mA}$ : la tension domestique est mortelle. Les trajets main-main et main-pied sont les plus graves car ils traversent le cœur.
\end{itemize}
\end{aretenir}

\begin{exercice}[Pour t'entraîner]
Un élève touche une prise défectueuse sous \SI{220}{V} avec les mains légèrement humides ; la résistance de son corps vaut alors $R = \SI{11000}{\ohm}$. 1) Calculer l'intensité qui traverse son corps. 2) À quel seuil physiologique correspond cette valeur ? Que risque concrètement l'élève ? 3) Expliquer pourquoi il ne faut jamais manipuler une installation électrique les pieds nus sur un carrelage mouillé.
\end{exercice}

\begin{corrige}[Exercice]
1) Loi d'Ohm : $I = U/R = 220/11000 = \SI{0,02}{A} = \SI{20}{mA}$.\\
2) \SI{20}{mA} est supérieur au seuil de non lâcher (\SI{10}{mA}) : les muscles se tétanisent, l'élève risque de ne plus pouvoir lâcher le conducteur, ce qui prolonge le contact et aggrave les lésions ; on se rapproche de la zone de fibrillation (\SI{30}{mA}).\\
3) Les pieds nus sur un sol mouillé offrent un contact de grande surface avec la terre et une peau humide : la résistance du corps chute fortement. En cas de contact avec une phase, le courant traverse le corps selon le trajet main-pied, qui passe par le cœur, avec une intensité pouvant dépasser le seuil mortel.
\end{corrige}

\coursec{Les dangers du courant : effets sur l'organisme}

Dans la section précédente, nous avons vu que le corps humain est un \cle{conducteur} d'électricité : lorsqu'une personne touche un fil dénudé ou un appareil défectueux, le courant peut traverser son corps, d'un point d'entrée (souvent la main) vers un point de sortie (souvent les pieds posés au sol). Nous avons aussi appris que la résistance du corps dépend fortement des conditions : une peau sèche résiste bien, une peau mouillée très peu. Reste maintenant la question essentielle : \emph{que fait exactement le courant à notre organisme ?} Pourquoi une simple « châtaigne » peut-elle être mortelle ? C'est ce que nous allons étudier ici, en distinguant quatre effets du courant sur le corps : l'\cle{effet thermique}, l'effet sur les \cle{muscles}, l'effet sur le \cle{cœur} et l'effet sur la \cle{respiration}.

\begin{attention}[Pas d'expérience en classe]
Les seuils de danger que tu vas découvrir dans cette section ne proviennent \textbf{pas} d'expériences réalisables au collège : on ne fait évidemment jamais passer de courant dans le corps d'un être humain pour mesurer ses effets. Ces valeurs sont \cle{admises} : elles résultent de données médicales (accidents étudiés par les médecins), d'études sur des modèles animaux et de \cle{normes internationales} de sécurité électrique. Retiens simplement les ordres de grandeur : ils peuvent sauver des vies.
\end{attention}

\courssub{L'effet thermique : les brûlures électriques}

\textbf{Intuition.} Tu as déjà constaté qu'un fil qui transporte un courant trop intense chauffe : c'est l'\cle{effet Joule}, le même phénomène qui fait chauffer la résistance d'une bouilloire ou le filament d'une lampe. Or les tissus du corps humain (peau, muscles, nerfs, vaisseaux) sont aussi des conducteurs imparfaits : ils possèdent une résistance électrique. Lorsqu'un courant les traverse, une partie de l'énergie électrique est donc convertie en \cle{chaleur} à l'intérieur du corps.

\begin{definition}[Brûlure électrique]
La \cle{brûlure électrique} est une destruction des tissus provoquée par l'échauffement dû au passage du courant (effet Joule). Elle peut atteindre la peau (brûlure externe, visible aux points d'entrée et de sortie du courant) mais aussi les muscles, les vaisseaux et les nerfs situés \emph{en profondeur}, le long du trajet du courant (brûlure interne).
\end{definition}

\textbf{Mécanisme.} L'énergie thermique dégagée dans un conducteur de résistance $R$ traversé par un courant d'intensité $I$ pendant une durée $t$ est donnée par la loi de Joule :
\[ E = R \cdot I^{2} \cdot t \]
Cette formule explique deux choses capitales :
\begin{itemize}
\item l'échauffement dépend du \emph{carré} de l'intensité : si l'intensité double, la chaleur dégagée est multipliée par \textbf{quatre} ;
\item plus le courant passe \emph{longtemps}, plus l'échauffement est important : c'est pourquoi la rapidité de la coupure du circuit est vitale.
\end{itemize}

\begin{attention}[Le piège de la brûlure « invisible »]
Une brûlure électrique peut être \textbf{profonde même si la peau semble peu atteinte}. La surface de la peau ne montre parfois qu'une petite marque, alors que les tissus internes (muscles, vaisseaux) sont gravement détruits le long du trajet du courant. \textbf{Toute personne ayant subi un choc électrique doit consulter un médecin}, même si elle se sent bien et même si aucune brûlure n'est visible. Des complications cardiaques peuvent aussi apparaître plusieurs heures après l'accident.
\end{attention}

\courssub{L'effet sur les muscles : la tétanisation}

\textbf{Observation.} Lors d'accidents domestiques, on constate souvent que la victime reste « collée » au fil ou à l'appareil en défaut, incapable de le lâcher. Ce n'est ni de la volonté, ni de la peur : c'est un phénomène physiologique précis.

\begin{definition}[Tétanisation]
La \cle{tétanisation} est la contraction brutale et prolongée des muscles provoquée par le passage du courant électrique. Les muscles de la main se contractent si fort que la victime \textbf{ne peut plus lâcher} le conducteur qu'elle tient : le contact se prolonge, ce qui aggrave l'accident.
\end{definition}

\textbf{Explication.} Nos muscles reçoivent leurs ordres du cerveau sous forme de petits signaux \emph{électriques} transportés par les nerfs. Un courant extérieur qui traverse le corps vient parasiter ces signaux : il impose aux muscles une contraction permanente, beaucoup plus forte qu'une contraction volontaire. Comme les muscles qui \emph{referment} les doigts sont plus puissants que ceux qui les ouvrent, la main se crispe sur le conducteur au lieu de le repousser.

\begin{aretenir}[Le seuil de « non lâcher »]
À partir d'environ \SI{10}{mA} (soit \num{0,01} A) traversant le corps, la tétanisation des muscles devient bloquante : c'est le \cle{seuil de non lâcher}. En dessous, la victime ressent des picotements ou des crampes mais peut encore se dégager. Au-dessus, elle reste prisonnière du circuit.
\end{aretenir}

\begin{exemplebox}[Pourquoi on ne touche jamais directement la victime]
Si tu touches une personne tétanisée sur un fil, le courant peut traverser \emph{ton} corps à travers le sien : tu deviens à ton tour victime. La première action est toujours de \textbf{couper le courant} (disjoncteur, interrupteur général), jamais de tirer la victime à mains nues.
\end{exemplebox}

\courssub{L'effet sur le cœur : la fibrillation ventriculaire}

\textbf{Intuition.} Le cœur est un muscle particulier : il bat environ 60 à 100 fois par minute grâce à un « chef d'orchestre » naturel qui envoie des signaux électriques réguliers. Que se passe-t-il si un courant extérieur vient brouiller ce chef d'orchestre ?

\begin{definition}[Fibrillation ventriculaire]
La \cle{fibrillation ventriculaire} est un trouble du rythme cardiaque dans lequel les fibres musculaires du cœur se contractent de façon \textbf{désordonnée} et inefficace : le cœur tremble au lieu de battre. Le sang n'est plus propulsé dans l'organisme ; sans intervention rapide, c'est l'\cle{arrêt cardiaque}, puis la mort en quelques minutes.
\end{definition}

\begin{aretenir}[Le seuil de fibrillation]
Au-delà d'environ \SI{30}{mA} traversant le thorax pendant une durée suffisante, le risque de \cle{fibrillation ventriculaire} devient réel. Attention : ce seuil n'est pas un chiffre magique, il dépend aussi de la \emph{durée} du passage du courant (voir la courbe temps--intensité plus bas). C'est précisément pour cette raison que les \cle{interrupteurs différentiels} de sécurité sont réglés à \SI{30}{mA} : ils coupent le circuit \emph{avant} que le courant n'atteigne des valeurs et des durées mortelles (nous les étudierons en section 5).
\end{aretenir}

\courssub{L'effet sur la respiration : l'asphyxie}

La respiration est assurée par des muscles : le diaphragme et les muscles entre les côtes se contractent pour gonfler les poumons, puis se relâchent pour les vider. Un courant qui traverse le thorax peut \textbf{bloquer ces muscles respiratoires} par tétanisation : la victime ne peut plus inspirer ni expirer. C'est l'\cle{asphyxie} : le corps n'est plus alimenté en dioxygène, et la perte de conscience survient rapidement. Cet effet s'ajoute souvent aux troubles cardiaques, ce qui explique la gravité des accidents où le courant traverse le thorax (par exemple d'une main à l'autre).

\courssub{Deux facteurs décisifs : l'intensité et la durée}

\textbf{Démarche d'investigation à partir d'un document.} Les médecins et les organismes de normalisation ont rassemblé les données de nombreux accidents électriques. En reportant sur un graphique l'\emph{intensité} du courant qui a traversé le corps (axe vertical) et la \emph{durée} du passage (axe horizontal), on peut tracer des zones de danger croissant. Étudions ce document ensemble.

\begin{popfigure}
\begin{tikzpicture}
\begin{loglogaxis}[
width=0.95\linewidth,
height=7.5cm,
xlabel={Durée du passage du courant $t$ (s)},
ylabel={Intensité $I$ (mA)},
xmin=0.001, xmax=10,
ymin=0.5, ymax=1000,
grid=both,
grid style={popline!40},
]
\addplot[draw=none,fill=popgreenL,domain=0.001:10,samples=2]{10} \closedcycle;
\addplot[draw=none,fill=poporangeL,domain=0.01:10,samples=50]{min(300, 30/sqrt(max(0,x)))} \closedcycle;
\addplot[draw=none,fill=popgreenL,domain=0.01:10,samples=2]{10} \closedcycle;
\addplot[draw=none,fill=poppinkL,domain=0.01:10,samples=50]{1000} \closedcycle;
\addplot[draw=none,fill=poporangeL,domain=0.01:10,samples=50]{min(300, 30/sqrt(max(0,x)))} \closedcycle;
\addplot[popcurve,domain=0.01:10,samples=100]{min(300, 30/sqrt(max(0,x)))};
\addplot[popblue,very thick,dashed] coordinates {(0.03,0.5) (0.03,1000)};
\node[popblue,font=\small,anchor=west,rotate=90] at (axis cs:0.03,2) {coupure du différentiel 30 mA};
\node[popdark,font=\small\bfseries] at (axis cs:0.6,3) {Zone sans danger};
\node[popdark,font=\small\bfseries] at (axis cs:0.15,40) {Tétanisation};
\node[popdark,font=\small\bfseries] at (axis cs:0.35,400) {Fibrillation};
\end{loglogaxis}
\end{tikzpicture}
\legende{Courbe temps--intensité (échelles logarithmiques) : plus l'intensité est grande, plus la durée tolérable est courte. En vert, la zone sans danger ; en orange, la zone de tétanisation ; en rose, la zone où la fibrillation ventriculaire peut survenir. La ligne bleue en pointillés montre le temps de coupure d'un interrupteur différentiel de \SI{30}{mA} (environ \SI{0,03}{s}) : il agit dans la zone sûre, avant que le danger ne devienne mortel.}
\end{popfigure}

\textbf{Interprétation du document.}
\begin{enumerate}
\item \emph{Observation} : pour une intensité donnée, il existe une durée maximale tolérable ; plus l'intensité est grande, plus cette durée est courte.
\item \emph{Lecture} : à \SI{10}{mA}, la tétanisation apparaît ; au-dessus de la courbe limite (zone rose), la fibrillation menace.
\item \emph{Conclusion} : le danger ne dépend pas seulement de l'intensité, mais du \textbf{couple intensité--durée}. Un dispositif qui coupe le circuit en une fraction de seconde fait passer la victime de la zone mortelle à la zone sûre.
\end{enumerate}

\begin{propriete}[Propriété admise : les deux facteurs du danger]
Le danger d'un choc électrique \cle{croît avec l'intensité} du courant qui traverse le corps \cle{et avec la durée} de son passage. Une coupure rapide du circuit (en moins de \SI{0,1}{s}) empêche le courant d'agir assez longtemps pour provoquer la fibrillation : \textbf{une coupure rapide sauve des vies}. C'est le principe de tous les dispositifs de protection modernes.
\end{propriete}

\begin{exemplebox}[Exemple chiffré : la durée change tout]
Un courant de \SI{50}{mA} qui traverse le corps pendant \SI{1}{s} peut provoquer une fibrillation ventriculaire mortelle. Le \emph{même} courant de \SI{50}{mA} pendant seulement \SI{0,05}{s} est généralement supporté sans séquelle : le cœur n'a pas été perturbé assez longtemps. C'est exactement ce que réalise l'interrupteur différentiel : il ne réduit pas le courant, il réduit la \emph{durée} de l'accident.
\end{exemplebox}

\begin{exoresolu}[Lire la courbe temps--intensité]
Énoncé. Lors d'un accident domestique, un courant de \SI{100}{mA} traverse le corps d'une personne. Deux cas sont possibles : (a) le contact dure \SI{1}{s} car personne ne coupe le circuit ; (b) un interrupteur différentiel coupe le circuit en \SI{0,03}{s}. En t'aidant de la courbe temps--intensité, indique dans quelle zone se situe chaque cas et conclus sur l'état probable de la victime.
\tcblower
Solution détaillée.
\begin{itemize}
\item Cas (a) : $I = \SI{100}{mA}$ et $t = \SI{1}{s}$. Sur le graphique, le point correspondant se trouve \emph{au-dessus} de la courbe limite : c'est la \textbf{zone de fibrillation}. La victime risque l'arrêt cardiaque.
\item Cas (b) : $I = \SI{100}{mA}$ et $t = \SI{0,03}{s}$. Le point se situe \emph{en dessous} de la courbe limite, près de la ligne bleue en pointillés : la durée est trop courte pour que la fibrillation s'installe. La victime subit une secousse et une tétanisation brève, mais sa vie n'est normalement pas en danger.
\end{itemize}
Conclusion : à intensité égale, c'est la \textbf{durée} du passage du courant qui fait la différence entre un accident grave et un accident sans conséquence. D'où l'importance vitale des dispositifs de coupure rapide.
\end{exoresolu}

\begin{savaistu}[Le choc électrique qui soigne : le défibrillateur]
Le courant électrique peut tuer... mais il peut aussi \emph{ressusciter} ! Le \cle{défibrillateur}, appareil que l'on trouve désormais dans les hôpitaux, les stades et certains lieux publics, envoie un \textbf{choc électrique contrôlé} (intensité et durée précisément réglées) à travers le thorax d'une personne en fibrillation. Ce choc « remet à zéro » toutes les fibres musculaires du cœur en même temps, ce qui permet au chef d'orchestre naturel de reprendre le commandement : le cœur se remet à battre régulièrement. C'est un bel exemple de la science transformant un danger en outil de sauvetage.
\end{savaistu}

\begin{attention}[Erreurs fréquentes au BEPC]
\begin{itemize}
\item Ne confonds pas \cle{électrisation} (la victime est traversée par le courant mais survit) et \cle{électrocution} (le choc électrique entraîne la mort).
\item Ne dis pas « c'est la tension qui tue » sans nuance : c'est \textbf{l'intensité traversant le corps} (qui dépend de la tension \emph{et} de la résistance du corps) \textbf{et la durée} qui déterminent le danger.
\item Retiens les seuils dans le bon ordre : picotements dès environ \SI{1}{mA}, \textbf{non lâcher vers \SI{10}{mA}}, \textbf{fibrillation au-delà de \SI{30}{mA}} (pour une durée suffisante).
\end{itemize}
\end{attention}

\begin{exercice}[À toi de jouer]
\begin{enumerate}
\item Cite les quatre effets du courant électrique sur l'organisme et associe à chacun sa conséquence principale.
\item Explique, à l'aide de la loi de Joule $E = R I^{2} t$, pourquoi un courant deux fois plus intense est beaucoup plus dangereux qu'on ne pourrait le croire.
\item Une victime reste « collée » à un fil dénudé. Nomme le phénomène, donne le seuil d'intensité correspondant et indique la conduite à tenir pour la secourir sans danger.
\end{enumerate}
\end{exercice}

\begin{corrige}[Exercice : À toi de jouer]
\begin{enumerate}
\item Effet thermique $\rightarrow$ brûlures internes et externes ; effet sur les muscles $\rightarrow$ tétanisation (impossibilité de lâcher) ; effet sur le cœur $\rightarrow$ fibrillation ventriculaire puis arrêt cardiaque ; effet sur la respiration $\rightarrow$ blocage des muscles respiratoires et asphyxie.
\item L'énergie thermique dégagée est proportionnelle au \emph{carré} de l'intensité : si $I$ double, $E$ est multipliée par $2^{2} = 4$. Un courant deux fois plus intense échauffe donc les tissus quatre fois plus, d'où des brûlures bien plus graves.
\item C'est la \cle{tétanisation}, qui apparaît au seuil de non lâcher, environ \SI{10}{mA}. Conduite à tenir : couper d'abord le courant (disjoncteur ou interrupteur général), ne jamais toucher la victime à mains nues, puis alerter les secours et consulter un médecin même si la victime semble indemne.
\end{enumerate}
\end{corrige}

\coursec{Les risques pour les installations : surcharge et court-circuit}

Dans la section précédente, nous avons vu que le courant électrique peut être dangereux pour le \cle{corps humain} : électrisation, électrocution. Mais le courant peut aussi devenir dangereux \emph{sans même toucher une personne}. Il suffit parfois d'une multiprise surchargée dans un salon de Douala, ou de deux fils dénudés qui se touchent derrière un réfrigérateur, pour qu'un incendie se déclare. Chaque année au Cameroun, de nombreuses maisons et boutiques partent en fumée à cause d'installations électriques défaillantes. Dans cette section, nous allons comprendre les deux grands accidents qui menacent les \cle{installations} : la \cle{surcharge} et le \cle{court-circuit}.

\courssub{La surcharge : trop d'appareils sur une même ligne}

\textbf{Une situation de la vie courante.}\par
C'est le soir à Yaoundé. Sur une seule multiprise branchée au mur, on trouve : la télévision, le décodeur, le ventilateur, le chargeur de téléphone, et on vient d'y ajouter la bouilloire électrique pour préparer le thé. Au bout de quelques minutes, le câble de la multiprise devient tiède, puis chaud. Une odeur de plastique chauffé se répand dans la pièce. Que se passe-t-il ?

\textbf{Observation et hypothèse.}\par
Le câble chauffe alors qu'aucun appareil n'est en panne. Hypothèse : le courant qui circule dans le câble est devenu \emph{trop grand} pour ce câble. Vérifions cette hypothèse par l'expérience, puis par le calcul.

\begin{experience}[Pourquoi les intensités s'additionnent-elles ?]
\textbf{Matériel :} un générateur, trois lampes identiques, quatre ampèremètres, des fils de connexion.

\textbf{Protocole :}
\begin{enumerate}
\item Brancher une seule lampe et mesurer l'intensité $I_1$ qui la traverse : on lit par exemple $I_1 = \SI{0,2}{A}$.
\item Brancher une deuxième lampe en dérivation : mesurer l'intensité totale $I$ fournie par le générateur. On lit $I = \SI{0,4}{A}$.
\item Ajouter une troisième lampe : on lit $I = \SI{0,6}{A}$.
\end{enumerate}

\textbf{Observation :} à chaque lampe ajoutée en dérivation, l'intensité totale dans le fil principal augmente : c'est la somme des intensités de chaque branche ($0,2 + 0,2 = 0,4$ puis $0,2 + 0,2 + 0,2 = 0,6$).

\textbf{Interprétation :} dans une installation domestique, tous les appareils sont branchés en dérivation (en parallèle) sous la même tension de \SI{220}{V}. L'intensité totale qui traverse le fil d'arrivée et la multiprise est donc :
\[ I_{\text{total}} = I_1 + I_2 + I_3 + \cdots \]
Plus on branche d'appareils, plus l'intensité totale est grande. Or chaque fil est prévu pour supporter une intensité maximale. Au-delà, il chauffe dangereusement.

\textbf{Conclusion :} dans un circuit en dérivation, les intensités s'additionnent dans la branche principale. C'est la \emph{loi des nœuds}, déjà rencontrée dans les leçons sur les circuits.
\end{experience}

\begin{definition}[Surcharge]
On dit qu'un circuit électrique est en \cle{surcharge} lorsque l'intensité totale du courant qui le traverse \textbf{dépasse l'intensité maximale admissible} pour ses conducteurs (fils, prises, multiprises). Cela se produit en général lorsqu'on branche \textbf{trop d'appareils} sur une même ligne ou sur une même multiprise.
\end{definition}

\begin{popfigure}
\begin{tikzpicture}[scale=0.92, every node/.style={font=\small}]
% Mur et prise
\fill[popcream] (-0.4,-0.4) rectangle (1.2,3.4);
\draw[popmuted] (-0.4,-0.4) rectangle (1.2,3.4);
\node[popmuted, rotate=90] at (-0.1,1.5) {\footnotesize mur};
% Prise murale
\draw[fill=white, thick, popdark] (1.2,1.2) rectangle (2.0,2.0);
\fill[popdark] (1.6,1.45) circle (0.05);
\fill[popdark] (1.6,1.75) circle (0.05);
\node[below, popdark] at (1.6,1.2) {\footnotesize prise};
% Cable vers multiprise
\draw[very thick, poporange] (2.0,1.6) -- (3.2,1.6);
\node[above, poporange] at (2.6,1.65) {\footnotesize $I_{\text{total}}$};
% Multiprise
\draw[fill=poporangeL, very thick, poporange, rounded corners=3pt] (3.2,0.6) rectangle (5.4,2.6);
\node[popdark] at (4.3,2.85) {\footnotesize \textbf{multiprise}};
\foreach \y in {1.0,1.6,2.2}{
  \draw[fill=white, thick, popdark] (3.5,\y-0.18) rectangle (4.1,\y+0.18);
  \fill[popdark] (3.8,\y-0.07) circle (0.03);
  \fill[popdark] (3.8,\y+0.07) circle (0.03);
}
% Appareils
\draw[thick, popblue] (4.1,2.2) -- (6.0,2.2) -- (6.0,3.0);
\draw[fill=popblueL, thick, popblue] (6.0,3.0) rectangle (7.8,3.9);
\node[popdark] at (6.9,3.45) {\footnotesize bouilloire};
\node[popblue, below] at (6.9,3.0) {\footnotesize $I_1 \approx \SI{9,1}{A}$};
\draw[thick, popgreen] (4.1,1.6) -- (6.4,1.6) -- (6.4,2.0);
\draw[fill=popgreenL, thick, popgreen] (6.4,2.0) rectangle (8.2,2.6);
\node[popdark] at (7.3,2.3) {\footnotesize fer \`a repasser};
\node[popgreen, below] at (7.3,2.0) {\footnotesize $I_2 \approx \SI{5,5}{A}$};
\draw[thick, poppurple] (4.1,1.0) -- (6.8,1.0) -- (6.8,1.4);
\draw[fill=poppurpleL, thick, poppurple] (6.8,1.4) rectangle (8.6,0.6);
\node[popdark] at (7.7,1.0) {\footnotesize four};
\node[poppurple, below] at (7.7,0.6) {\footnotesize $I_3 \approx \SI{9,1}{A}$};
% Somme
\draw[->, very thick, poppink] (2.6,0.9) -- (2.6,0.2);
\node[below, poppink, align=center] at (2.6,0.15) {\footnotesize $I_{\text{total}} = I_1+I_2+I_3$\\ \footnotesize $\approx \SI{23,7}{A} > \SI{16}{A}$ : danger !};
% Chaleur sur cable
\foreach \x in {2.3,2.6,2.9}{
  \draw[thick, poppink, decorate, decoration={snake, amplitude=0.6mm, segment length=2mm}] (\x,1.85) -- (\x,2.25);
}
\end{tikzpicture}
\legende{Multiprise surchargée : les intensités des appareils branchés en dérivation s'additionnent dans le câble d'alimentation. Si $I_{\text{total}}$ dépasse l'intensité maximale admissible (ici \SI{16}{A}), le câble s'échauffe par effet Joule (ondulations roses).}
\end{popfigure}

\textbf{Pourquoi la surcharge fait-elle chauffer les fils ?}\par
Tout conducteur traversé par un courant dégage de la chaleur : c'est l'\cle{effet Joule}. La puissance dissipée en chaleur dans un fil de résistance $R$ traversé par un courant d'intensité $I$ est :
\[ P_{\text{chaleur}} = R \cdot I^2 \]
Remarque bien le carré : si l'intensité \emph{double}, la chaleur dégagée est \emph{multipliée par quatre}. Un fil prévu pour \SI{16}{A} et traversé par \SI{24}{A} dégage donc plus du double de chaleur que prévu (car $24^2 = 576$ contre $16^2 = 256$). Conséquences, par ordre de gravité :
\begin{itemize}
\item l'isolant plastique qui entoure le fil ramollit puis \textbf{fond} ;
\item les fils dénudés peuvent alors se toucher (ce qui provoque un court-circuit, voir plus loin) ;
\item les matériaux proches (bois, tissu, papier) peuvent \textbf{s'enflammer} : c'est le départ d'\cle{incendie}.
\end{itemize}

\begin{attention}[Un fil chaud n'est jamais anodin]
Erreur fréquente : croire qu'un fil ou une multiprise qui chauffe « un peu » est sans danger. \textbf{Faux !} Tout échauffement anormal d'un câble, d'une prise ou d'une fiche est un signal d'alarme : il annonce une surcharge ou un mauvais contact, et donc un risque d'incendie. Si tu sens une odeur de plastique chaud ou si une fiche est tiède : débranche immédiatement des appareils et préviens un adulte. De même, une prise noircie ou un câble dont l'isolant craquelle doivent être remplacés sans attendre.
\end{attention}

\courssub{Le court-circuit : le contact direct entre phase et neutre}

\textbf{Une situation de la vie courante.}\par
Dans un atelier de quartier, le fil d'un fer à repasser est usé : à force d'être plié, l'isolant s'est fendu et les deux fils de cuivre intérieurs — la \cle{phase} et le \cle{neutre} — se touchent. Au moment où on branche la fiche, une étincelle jaillit, un « clac » sec retentit, et toute la maison est plongée dans le noir : le fusible a fondu. Que s'est-il passé ?

\begin{definition}[Court-circuit]
Un \cle{court-circuit} est une \textbf{liaison accidentelle de très faible résistance} entre deux points d'un circuit qui sont à des potentiels différents (par exemple la phase et le neutre). Le courant ne traverse alors plus aucun appareil : la résistance du circuit devient \textbf{quasi nulle} et l'intensité devient \textbf{très grande}.
\end{definition}

\textbf{Pourquoi l'intensité devient-elle si grande ?}\par
Appliquons la loi d'Ohm : $I = \dfrac{U}{R}$. En fonctionnement normal, le courant traverse un appareil (bouilloire, fer, lampe) dont la résistance $R$ est assez grande : l'intensité reste modérée. En cas de court-circuit, le courant ne traverse plus que les fils, dont la résistance est très faible, par exemple $R \approx \SI{0,5}{\ohm}$ :
\[ I = \frac{U}{R} = \frac{220}{0,5} = \SI{440}{A} \]
C'est presque \textbf{trente fois} le calibre d'une ligne protégée par un fusible de \SI{16}{A} ! Une telle intensité, même pendant une fraction de seconde, provoque :
\begin{itemize}
\item des \textbf{étincelles} et une projection de métal fondu au point de contact ;
\item un \textbf{échauffement brutal} des conducteurs (effet Joule, proportionnel à $I^2$) ;
\item la \textbf{destruction} des appareils branchés et de l'installation ;
\item un \textbf{incendie} si aucun dispositif de protection ne coupe le circuit à temps.
\end{itemize}

\begin{popfigure}
\begin{tikzpicture}[scale=0.9, every node/.style={font=\small}]
% ===== Circuit normal =====
\node[popdark, font=\small\bfseries] at (2.6,3.6) {Circuit normal};
% source
\draw[thick, popdark] (0,0.4) -- (0,2.8);
\draw[very thick, popblue] (0,2.8) -- (5.2,2.8);
\draw[very thick, popgreen] (0,0.4) -- (5.2,0.4);
\node[above, popblue] at (2.6,2.85) {\footnotesize phase};
\node[below, popgreen] at (2.6,0.35) {\footnotesize neutre};
% lampe (appareil)
\draw[fill=popgoldL, thick, popdark] (3.4,0.4) rectangle (4.6,2.8);
\node[popdark, align=center] at (4.0,1.6) {\footnotesize appareil\\ \footnotesize $R$ grand};
% courant
\draw[->, very thick, poporange] (0.7,2.8) -- (2.2,2.8);
\node[above, poporange] at (1.45,2.85) {\footnotesize $I$ normale};
\node[popdark, align=center] at (2.6,-0.5) {\footnotesize Le courant traverse l'appareil :\\ \footnotesize $I = \dfrac{U}{R}$ reste modérée.};

% ===== Court-circuit =====
\begin{scope}[xshift=8.2cm]
\node[popdark, font=\small\bfseries] at (2.6,3.6) {Court-circuit};
\draw[thick, popdark] (0,0.4) -- (0,2.8);
\draw[very thick, popblue] (0,2.8) -- (5.2,2.8);
\draw[very thick, popgreen] (0,0.4) -- (5.2,0.4);
\node[above, popblue] at (1.2,2.85) {\footnotesize phase};
\node[below, popgreen] at (1.2,0.35) {\footnotesize neutre};
% contact direct
\draw[very thick, poppink] (3.4,0.4) -- (3.4,2.8);
\fill[poppink] (3.4,1.6) circle (0.09);
% etincelle
\foreach \a in {20,70,120,200,250,300,340}{
  \draw[thick, popgold] (3.4,1.6) -- ++(\a:0.35);
}
\node[poppink, right, align=left] at (3.6,1.6) {\footnotesize contact direct\\ \footnotesize $R \approx 0$};
% courant enorme
\draw[->, line width=2.2pt, poppink] (0.7,2.8) -- (2.4,2.8);
\node[above, poppink] at (1.55,2.9) {\footnotesize $I$ \textbf{très grande}};
\node[popdark, align=center] at (2.6,-0.5) {\footnotesize Le courant évite l'appareil :\\ \footnotesize $I = \dfrac{U}{R}$ devient énorme.};
\end{scope}
\end{tikzpicture}
\legende{Comparaison : à gauche, circuit normal, le courant traverse l'appareil (résistance élevée). À droite, court-circuit : la phase et le neutre se touchent directement (fils dénudés), la résistance devient quasi nulle et l'intensité devient très grande (étincelles au point de contact).}
\end{popfigure}

\textbf{Les causes fréquentes de court-circuit au quotidien :}
\begin{itemize}
\item des \textbf{fils dénudés} dont l'isolant est usé, coupé ou rongé par les rongeurs, et « réparés » à la va-vite avec un simple bout de ruban ;
\item un \textbf{appareil défectueux} dont les fils internes se touchent ;
\item de l'\textbf{eau} qui pénètre dans une prise ou un appareil (l'eau courante conduit le courant) ;
\item des \textbf{installations vétustes} ou réalisées par des bricoleurs non qualifiés, avec des fils de fortune ;
\item des multiprises branchées \textbf{en cascade} (une multiprise sur une autre) qui concentrent les courants et échauffent les contacts jusqu'à fondre l'isolant.
\end{itemize}

\begin{attention}[Surcharge et court-circuit : ne pas confondre]
Au BEPC, les élèves confondent souvent ces deux accidents. Retiens bien :
\begin{itemize}
\item la \cle{surcharge} : \emph{trop d'appareils} branchés en même temps ; l'intensité dépasse \emph{progressivement} la limite admissible ; les fils chauffent \emph{lentement} ;
\item le \cle{court-circuit} : un \emph{contact direct} entre phase et neutre ; l'intensité devient \emph{brutalement} énorme ; étincelles et échauffement quasi instantané.
\end{itemize}
Dans les deux cas, le danger final est le même : l'\textbf{incendie}. Et dans les deux cas, c'est le fusible ou le disjoncteur qui sauve l'installation en coupant le courant.
\end{attention}

\begin{savaistu}[Incendies d'origine électrique au Cameroun]
Une grande partie des incendies de marchés et d'habitations au Cameroun est d'origine électrique : installations vétustes, branchements sauvages, multiprises surchargées, groupes électrogènes mal raccordés pendant les délestages. Les pompiers le répètent : la plupart de ces drames auraient pu être évités par une installation aux normes, un disjoncteur adapté et le remplacement des câbles abîmés. La sécurité électrique est d'abord une question de \textbf{prévention}.
\end{savaistu}

\courssub{Vérifier une installation par le calcul : la méthode}

Chaque appareil électrique porte une \textbf{puissance} $P$ (en watts, W) indiquée sur sa plaque signalétique. Sous la tension du réseau $U = \SI{220}{V}$, l'intensité qu'il consomme se calcule à partir de la relation $P = U \times I$, d'où :
\[ I = \frac{P}{U} \]
Une ligne protégée par un fusible de \SI{16}{A} peut donc alimenter une puissance maximale :
\[ P_{\max} = U \times I_{\max} = 220 \times 16 = \SI{3520}{W} \approx \SI{3500}{W} \]

\begin{methode}[Vérifier qu'une ligne n'est pas surchargée]
\begin{enumerate}
\item \textbf{Relever} la puissance de chaque appareil branché sur la ligne (plaque signalétique ou notice).
\item \textbf{Additionner} les puissances : $P_{\text{total}} = P_1 + P_2 + P_3 + \cdots$
\item \textbf{Calculer} l'intensité totale : $I = \dfrac{P_{\text{total}}}{U}$ avec $U = \SI{220}{V}$.
\item \textbf{Comparer} au calibre du fusible ou du disjoncteur de la ligne.
\item \textbf{Conclure} : si $I$ dépasse le calibre, il y a surcharge : débrancher un ou plusieurs appareils ou répartir sur une autre ligne.
\end{enumerate}
\end{methode}

\begin{exoresolu}[La multiprise du salon]
Sur une ligne de \SI{220}{V} protégée par un fusible de \SI{16}{A}, on branche sur une même multiprise : une bouilloire de \SI{2000}{W}, un fer à repasser de \SI{1200}{W} et un four de \SI{2000}{W}. Y a-t-il surcharge ? Que risque-t-on ?
\tcblower
\textbf{Étape 1 : puissance totale.}
\[ P_{\text{total}} = 2000 + 1200 + 2000 = \SI{5200}{W} \]
\textbf{Étape 2 : intensité totale.}
\[ I = \frac{P_{\text{total}}}{U} = \frac{5200}{220} \approx \SI{23,6}{A} \]
\textbf{Étape 3 : comparaison.}
\[ \SI{23,6}{A} > \SI{16}{A} \quad \text{(calibre du fusible)} \]
\textbf{Conclusion :} l'intensité demandée dépasse largement le calibre du fusible : il y a \textbf{surcharge}. Le fusible va fondre pour protéger l'installation ; s'il était absent ou surdimensionné, les câbles chaufferaient par effet Joule, l'isolant fondrait et un incendie pourrait se déclarer. Il faut débrancher au moins un appareil : par exemple, avec la bouilloire et le fer seuls, $P = \SI{3200}{W}$, soit $I = \dfrac{3200}{220} \approx \SI{14,5}{A} < \SI{16}{A}$ : la ligne est alors utilisée en sécurité.
\end{exoresolu}

\begin{exemplebox}[Et le court-circuit, chiffré]
Un fil dénudé crée un contact direct phase-neutre de résistance $R = \SI{0,5}{\ohm}$ sous \SI{220}{V}. Alors :
\[ I = \frac{U}{R} = \frac{220}{0,5} = \SI{440}{A} \]
Cette intensité gigantesque fait fondre le fusible de \SI{16}{A} presque instantanément : c'est précisément le rôle des dispositifs de protection, que nous étudierons dans la section suivante.
\end{exemplebox}

\begin{aretenir}[L'essentiel de la section]
\begin{itemize}
\item Une \cle{surcharge} se produit quand trop d'appareils branchés sur une même ligne font dépasser l'intensité maximale admissible : les intensités s'additionnent ($I_{\text{total}} = I_1 + I_2 + \cdots$).
\item Conséquence : échauffement des fils par \cle{effet Joule} ($P = R\,I^2$), fonte des isolants, risque d'\textbf{incendie}.
\item Un \cle{court-circuit} est un contact direct de très faible résistance entre phase et neutre : $I = \dfrac{U}{R}$ devient énorme, d'où étincelles et incendie.
\item Pour vérifier une ligne : additionner les puissances, calculer $I = \dfrac{P}{U}$, comparer au calibre du fusible.
\item Tout échauffement anormal d'un câble ou d'une prise est un signal de danger : on débranche et on fait réparer.
\end{itemize}
\end{aretenir}

\begin{exercice}[À toi de jouer]
Dans une chambre, une ligne de \SI{220}{V} est protégée par un fusible de \SI{10}{A}. On veut y brancher en même temps : un ordinateur de \SI{300}{W}, une télévision de \SI{150}{W}, un ventilateur de \SI{100}{W} et un radiateur électrique de \SI{1800}{W}.
\begin{enumerate}
\item Calculer la puissance totale demandée.
\item Calculer l'intensité totale et la comparer au calibre du fusible. Conclure.
\item Quelle puissance maximale cette ligne peut-elle supporter ?
\item Proposer une solution pour utiliser ces appareils en sécurité.
\end{enumerate}
\end{exercice}

\begin{corrige}[Exercice : la chambre]
\begin{enumerate}
\item $P_{\text{total}} = 300 + 150 + 100 + 1800 = \SI{2350}{W}$.
\item $I = \dfrac{2350}{220} \approx \SI{10,7}{A}$. Or $\SI{10,7}{A} > \SI{10}{A}$ : il y a \textbf{surcharge}, le fusible fondra.
\item $P_{\max} = U \times I_{\max} = 220 \times 10 = \SI{2200}{W}$.
\item Il faut retirer un appareil : sans le radiateur, $P = 300 + 150 + 100 = \SI{550}{W}$ et $I = \dfrac{550}{220} \approx \SI{2,5}{A}$, très en dessous du calibre. Le radiateur doit être branché sur une autre ligne plus puissante (par exemple une ligne \SI{16}{A} dédiée, qui supporte jusqu'à \SI{3520}{W}).
\end{enumerate}
\end{corrige}

Nous savons désormais reconnaître et prévenir les deux grands accidents des installations. Reste une question essentielle : quels sont les \textbf{dispositifs} capables de couper automatiquement le courant avant la catastrophe ? C'est l'objet de la section suivante : le fusible, le disjoncteur, la prise de terre et l'interrupteur différentiel.

\coursec{Les dispositifs de protection : fusible et disjoncteur}

Dans la section précédente, nous avons vu que la \cle{surcharge} et le \cle{court-circuit} font circuler dans les fils des intensités dangereusement élevées : les conducteurs s'échauffent par \cle{effet Joule}, l'isolant fond, et l'incendie peut se déclarer. Face à ces risques, l'installation électrique doit posséder des « gardiens » capables de couper le courant avant la catastrophe. Ces gardiens, ce sont le \cle{fusible} et le \cle{disjoncteur}. Dans cette section, nous allons comprendre comment ils fonctionnent, comment choisir leur calibre, et pourquoi il ne faut jamais les trafiquer.

\courssub{Le fusible : un fil qui se sacrifie pour protéger l'installation}

\textbf{Une idée simple : le maillon faible volontaire.}\par
Imagine une chaîne dont l'un des maillons serait volontairement plus fragile que les autres. Si l'on tire trop fort, c'est ce maillon qui casse le premier, et la chaîne s'ouvre avant que le reste ne soit abîmé. Le fusible fonctionne exactement ainsi : on insère dans le circuit un \cle{fil conducteur calibré}, plus fragile que les fils de l'installation, qui fondra le premier en cas de surintensité.

\begin{definition}[Le fusible]
Un \cle{fusible} est un élément de protection constitué d'un fil conducteur calibré (souvent en plomb ou en alliage à bas point de fusion), placé \textbf{en série} dans le circuit, sur le fil de phase. Lorsque l'intensité du courant dépasse une valeur limite appelée \cle{calibre}, le fil fond par \cle{effet Joule} : le circuit s'ouvre et le courant est coupé. Un fusible fondu doit être \textbf{remplacé} par un fusible neuf de même calibre.
\end{definition}

\textbf{Pourquoi le fil fond-il ?}\par
Tout conducteur traversé par un courant s'échauffe : c'est l'effet Joule. Plus l'intensité est grande, plus l'échauffement est important. Le fil du fusible est choisi fin et fait d'un métal qui fond à basse température. En fonctionnement normal, l'échauffement reste modéré. Mais en cas de surcharge ou de court-circuit, l'intensité dépasse le calibre : la température du fil atteint son point de fusion, le fil fond, et le circuit est interrompu. Le fusible « se sacrifie » pour sauver l'installation.

\begin{experience}[Fusion d'un fil fin en court-circuit (démonstration en classe, très basse tension de sécurité)]
\textbf{Matériel :} un générateur de très basse tension ($\SI{6}{V}$), un interrupteur, un support, un fil de fer très fin (quelques centimètres), des fils de connexion.\par
\textbf{Protocole :}
\begin{enumerate}
\item Réaliser un circuit comportant le générateur, l'interrupteur et le fil fin tendu entre deux bornes.
\item Fermer l'interrupteur de manière à créer un court-circuit aux bornes du générateur.
\item Observer attentivement le fil fin, puis ouvrir aussitôt l'interrupteur.
\end{enumerate}
\textbf{Observations :} le fil fin rougeoit puis fond presque instantanément ; le circuit s'ouvre et le courant cesse de circuler.\par
\textbf{Interprétation :} le court-circuit a produit une intensité très élevée. L'effet Joule a porté le fil fin à sa température de fusion. C'est exactement le mécanisme d'un fusible : le fil calibré fond et protège le reste du circuit.\par
\textbf{Conclusion :} un conducteur calibré placé en série dans un circuit le protège contre les surintensités en fondant le premier.\par
\textbf{Sécurité :} cette expérience est réalisée par le professeur, en très basse tension ($\SI{6}{V}$), jamais sur le secteur $\SI{220}{V}$.
\end{experience}

\begin{popfigure}
\begin{center}
\begin{tikzpicture}[scale=1.0]
% Corps de la cartouche
\draw[thick, fill=popcream] (0,-0.5) rectangle (6,0.5);
% Embouts métalliques
\draw[thick, fill=popmuted] (-0.8,-0.6) rectangle (0,0.6);
\draw[thick, fill=popmuted] (6,-0.6) rectangle (6.8,0.6);
% Fil calibré intact (fusible du haut)
\draw[very thick, popblue] (0,0) -- (6,0);
\draw[popblue, ->] (3,0.15) -- (3,0.9) node[above, popblue]{\small fil calibr\'e intact};
% Labels
\node[below] at (-0.4,-0.7) {\small embout};
\node[below] at (6.4,-0.7) {\small embout};
\node at (3,-1.1) {\small \textbf{Fusible en bon état : le courant passe}};
% Deuxième fusible : fondu
\begin{scope}[yshift=-3.2cm]
\draw[thick, fill=popcream] (0,-0.5) rectangle (6,0.5);
\draw[thick, fill=popmuted] (-0.8,-0.6) rectangle (0,0.6);
\draw[thick, fill=popmuted] (6,-0.6) rectangle (6.8,0.6);
% Fil rompu
\draw[very thick, poporange] (0,0) -- (2.4,0);
\draw[very thick, poporange] (3.6,0) -- (6,0);
% Gouttelettes de métal fondu
\fill[poporange] (2.55,0) circle (0.08);
\fill[poporange] (3.45,0) circle (0.08);
\draw[poporange, ->] (3,0.2) -- (3,0.9) node[above, poporange]{\small fil fondu : circuit ouvert};
% Croix de coupure
\draw[very thick, poppink] (2.75,-0.35) -- (3.25,0.35);
\draw[very thick, poppink] (2.75,0.35) -- (3.25,-0.35);
\node at (3,-1.1) {\small \textbf{Fusible fondu (surcharge ou court-circuit) : le courant ne passe plus}};
\end{scope}
\end{tikzpicture}
\end{center}
\legende{Schéma en coupe d'un fusible à cartouche. En haut : le fil calibré est intact, le courant circule. En bas : après une surintensité, le fil a fondu par effet Joule ; le circuit est ouvert, l'installation est protégée.}
\end{popfigure}

\courssub{Le calibre d'un fusible : une valeur à respecter}

\begin{definition}[Calibre]
Le \cle{calibre} d'un fusible (ou d'un disjoncteur) est l'intensité maximale qu'il laisse passer durablement sans réagir. Il s'exprime en ampères (A) : calibres courants $\SI{10}{A}$, $\SI{16}{A}$, $\SI{20}{A}$, $\SI{32}{A}$. Le calibre doit être \textbf{adapté à la section des fils} de la ligne et aux appareils qu'elle alimente.
\end{definition}

Chaque ligne de l'installation (prises, éclairage, cuisson) est protégée par un fusible ou un disjoncteur dont le calibre correspond aux fils utilisés : des fils fins ne supportent qu'une faible intensité, donc un petit calibre ; des fils de grosse section (ligne de la cuisinière ou du chauffe-eau) supportent un calibre plus grand.

\begin{methode}[Choisir le calibre d'une protection]
\begin{enumerate}
\item Calculer l'intensité totale demandée par les appareils de la ligne : $I = \dfrac{P}{U}$ (avec $P$ en watts et $U = \SI{220}{V}$).
\item Choisir le calibre \textbf{immédiatement supérieur} à cette intensité, tout en restant compatible avec la section des fils.
\item Vérifier qu'en fonctionnement normal, l'intensité reste inférieure au calibre, sinon le fusible fondra ou le disjoncteur déclenchera sans défaut réel.
\end{enumerate}
\end{methode}

\begin{exemplebox}[Choisir le bon calibre]
Un chauffe-eau de puissance $P = \SI{2200}{W}$ est branché sous la tension $U = \SI{220}{V}$. L'intensité qui le traverse en fonctionnement normal est :
\[ I = \frac{P}{U} = \frac{2200}{220} = \SI{10}{A}. \]
On choisit un fusible de calibre immédiatement supérieur : un fusible de $\SI{16}{A}$ convient. Un fusible de $\SI{10}{A}$ risquerait de fondre en fonctionnement normal ; un fusible de $\SI{32}{A}$ laisserait passer des surcharges dangereuses sans réagir.
\end{exemplebox}

\begin{attention}[Le fusible ne protège pas les personnes !]
Erreur fréquente : croire que le fusible protège les personnes contre l'électrisation. C'est \textbf{faux}. Le seuil mortel pour le corps humain est d'environ $\SI{0,1}{A}$, alors que les calibres des fusibles domestiques vont de $\SI{10}{A}$ à $\SI{32}{A}$ : un courant cent fois supérieur au seuil mortel peut traverser ton corps sans que le fusible ne fonde. Le fusible protège \textbf{l'installation et les appareils} contre les surcharges et les courts-circuits. Seul l'\cle{interrupteur différentiel} (étudié dans la section suivante) protège efficacement les personnes.
\end{attention}

\courssub{Le disjoncteur : une protection réarmable}

\textbf{Le problème du fusible :} une fois fondu, il est hors d'usage et il faut le remplacer. Le soir, dans une maison sans fusible de rechange, c'est la panne assurée ! C'est pourquoi les installations modernes utilisent des \cle{disjoncteurs}.

\begin{definition}[Le disjoncteur]
Un \cle{disjoncteur} est un appareil de protection qui \textbf{coupe automatiquement} le circuit lorsque l'intensité dépasse son calibre (surcharge ou court-circuit). Contrairement au fusible, il n'est pas détruit : après élimination de la cause du défaut, on le \textbf{réarme} simplement en relevant son levier. On parle de \cle{disjoncteur divisionnaire} lorsqu'il protège une seule ligne de l'installation, et de \cle{disjoncteur général} (ou disjoncteur de branchement, posé par ENEO) lorsqu'il protège toute l'installation.
\end{definition}

\textbf{Comment déclenche-t-il ?}\par
Le disjoncteur possède deux mécanismes de détection :
\begin{itemize}
\item un mécanisme \textbf{thermique} (une lame bimétallique qui se courbe en chauffant) qui réagit aux \textbf{surcharges} : l'intensité dépasse légèrement le calibre pendant un certain temps ;
\item un mécanisme \textbf{magnétique} (une bobine qui attire un noyau) qui réagit instantanément aux \textbf{courts-circuits}, où l'intensité devient brutalement très grande.
\end{itemize}
Dans les deux cas, un levier bascule, les contacts s'ouvrent : le circuit est coupé. On dit couramment que le disjoncteur « saute ».

\begin{popfigure}
\begin{center}
\begin{tikzpicture}[scale=0.95]
% Arrivée ENEO
\draw[very thick] (0,0) -- (1.2,0);
\node[left] at (0,0) {\small Arriv\'ee ENEO};
% Disjoncteur général
\draw[thick, fill=popblueL] (1.2,-0.5) rectangle (2.4,0.5);
\node at (1.8,0) {\small DG};
\node[below, align=center] at (1.8,-0.6) {\tiny disjoncteur\\[-2pt]\tiny g\'en\'eral};
% Barre de répartition
\draw[very thick] (2.4,0) -- (3.2,0);
\draw[very thick, popdark] (3.2,-2.2) -- (3.2,2.2);
% Ligne 1 : éclairage
\draw[very thick] (3.2,1.8) -- (4.2,1.8);
\draw[thick, fill=popgreenL] (4.2,1.5) rectangle (5.2,2.1);
\node at (4.7,1.8) {\tiny 10 A};
\draw[very thick] (5.2,1.8) -- (6.6,1.8);
\node[right] at (6.6,1.8) {\small Ligne \'eclairage};
% Ligne 2 : prises
\draw[very thick] (3.2,0) -- (4.2,0);
\draw[thick, fill=popgreenL] (4.2,-0.3) rectangle (5.2,0.3);
\node at (4.7,0) {\tiny 16 A};
\draw[very thick] (5.2,0) -- (6.6,0);
\node[right] at (6.6,0) {\small Ligne prises de courant};
% Ligne 3 : cuisson / chauffe-eau
\draw[very thick] (3.2,-1.8) -- (4.2,-1.8);
\draw[thick, fill=popgreenL] (4.2,-2.1) rectangle (5.2,-1.5);
\node at (4.7,-1.8) {\tiny 20 A};
\draw[very thick] (5.2,-1.8) -- (6.6,-1.8);
\node[right] at (6.6,-1.8) {\small Ligne cuisson / chauffe-eau};
% Cadre du tableau
\draw[thick, popmuted, dashed] (0.9,-2.6) rectangle (6.9,2.6);
\node[popmuted] at (3.9,2.85) {\small Tableau \'electrique de la maison};
\end{tikzpicture}
\end{center}
\legende{Schéma d'un tableau électrique domestique. Le disjoncteur général (DG) protège toute l'installation. Chaque ligne est protégée par un disjoncteur divisionnaire dont le calibre est adapté : $\SI{10}{A}$ pour l'éclairage, $\SI{16}{A}$ pour les prises, $\SI{20}{A}$ pour la cuisson.}
\end{popfigure}

\begin{aretenir}[Fusible ou disjoncteur : ce qu'il faut retenir]
\begin{itemize}
\item Les deux protègent l'\textbf{installation} contre les surcharges et les courts-circuits, en coupant le circuit quand l'intensité dépasse le calibre.
\item Le \textbf{fusible} fond : il doit être \textbf{remplacé} après chaque défaut.
\item Le \textbf{disjoncteur} déclenche : il est \textbf{réarmable} après réparation du défaut.
\item Aucun des deux ne protège directement les personnes : c'est le rôle du différentiel.
\end{itemize}
\end{aretenir}

\courssub{Une règle d'or : ne jamais trafiquer la protection}

Il arrive qu'un fusible fonde « trop souvent » parce que la ligne est surchargée. La mauvaise réaction consiste à mettre un fusible de calibre plus grand, ou pire, à remplacer le fusible par un fil de fer ou un clou : on dit qu'on « shunte » le fusible. C'est extrêmement dangereux.

\begin{attention}[Jamais de fil quelconque à la place d'un fusible]
Remplacer un fusible par un calibre plus grand ou par un fil quelconque (clou, fil de fer) \textbf{supprime toute protection} : en cas de surcharge ou de court-circuit, rien ne coupe plus le circuit, les fils de l'installation s'échauffent jusqu'à provoquer un \textbf{incendie}. De nombreux incendies domestiques au Cameroun ont cette cause. La bonne conduite : remplacer le fusible par un fusible \textbf{de même calibre}, et si le problème revient, réduire le nombre d'appareils branchés sur la ligne ou faire vérifier l'installation par un technicien.
\end{attention}

\begin{exoresolu}[Un fusible qui fond : que faire ?]
\textbf{Énoncé.} Dans la maison de la famille Ngo Bell, la ligne des prises est protégée par un fusible de calibre $\SI{16}{A}$. Un soir, on branche sur cette ligne une bouilloire de $\SI{2000}{W}$, un fer à repasser de $\SI{1200}{W}$ et un radiateur de $\SI{800}{W}$, sous une tension de $\SI{220}{V}$. Le fusible fond.
\begin{enumerate}
\item Calculer l'intensité totale demandée par les trois appareils.
\item Expliquer pourquoi le fusible a fondu.
\item Le père propose de mettre un fusible de $\SI{32}{A}$. Est-ce une bonne idée ? Proposer une conduite correcte.
\end{enumerate}
\tcblower
\textbf{Solution détaillée.}
\begin{enumerate}
\item La puissance totale est $P = 2000 + 1200 + 800 = \SI{4000}{W}$. L'intensité totale est :
\[ I = \frac{P}{U} = \frac{4000}{220} \approx \SI{18,2}{A}. \]
\item L'intensité demandée ($\SI{18,2}{A}$) dépasse le calibre du fusible ($\SI{16}{A}$) : c'est une \textbf{surcharge}. Par effet Joule, le fil calibré du fusible a fondu, ce qui a coupé le circuit et protégé l'installation.
\item Non, c'est une mauvaise idée : avec un fusible de $\SI{32}{A}$, la surcharge ne serait plus détectée et les fils de la ligne (prévus pour $\SI{16}{A}$) s'échaufferaient dangereusement, avec un risque d'incendie. La conduite correcte est de \textbf{remettre un fusible de $\SI{16}{A}$} et de \textbf{ne pas brancher les trois appareils en même temps} sur cette ligne (par exemple, débrancher le radiateur quand la bouilloire fonctionne).
\end{enumerate}
\end{exoresolu}

\begin{savaistu}[Le fusible, une invention née du télégraphe]
Le premier fusible a été imaginé au XIX\textsuperscript{e} siècle pour protéger les câbles télégraphiques contre la foudre. Les lignes télégraphiques, très longues et exposées aux orages, recevaient parfois des décharges électriques immenses qui détruisaient les appareils. Les ingénieurs eurent l'idée d'insérer dans la ligne un fil fin qui fondrait le premier, sacrifié pour sauver le matériel coûteux. Le principe n'a pas changé depuis : dans ton tableau électrique comme dans ta radio ou ton chargeur de téléphone, un petit fusible veille toujours. Au Cameroun, où la foudre frappe souvent les installations pendant la saison des pluies, ces protections restent indispensables.
\end{savaistu}

\begin{exercice}[À toi de jouer]
Une ligne de prises est protégée par un disjoncteur divisionnaire de calibre $\SI{16}{A}$, sous $\SI{220}{V}$.
\begin{enumerate}
\item Quelle est la puissance maximale totale que cette ligne peut alimenter ?
\item Peut-on y brancher simultanément un four de $\SI{2500}{W}$ et un fer à repasser de $\SI{1000}{W}$ ? Justifier par un calcul.
\item Le disjoncteur « saute ». Que faut-il faire avant de le réarmer ?
\end{enumerate}
\end{exercice}

\begin{corrige}[Exercice : À toi de jouer]
\begin{enumerate}
\item $P_{\max} = U \times I = 220 \times 16 = \SI{3520}{W}$.
\item Puissance totale : $2500 + 1000 = \SI{3500}{W}$, soit $I = \dfrac{3500}{220} \approx \SI{15,9}{A} < \SI{16}{A}$. Oui, c'est possible, mais on est très proche de la limite : tout appareil supplémentaire ferait déclencher le disjoncteur.
\item Avant de réarmer, il faut \textbf{identifier et supprimer la cause du déclenchement} : débrancher les appareils en trop (surcharge) ou faire réparer l'appareil ou le fil défectueux (court-circuit). Ensuite seulement, on relève le levier du disjoncteur.
\end{enumerate}
\end{corrige}

\begin{aretenir}[L'essentiel de la section]
\begin{itemize}
\item Le \cle{fusible} est un fil calibré qui \textbf{fond par effet Joule} quand l'intensité dépasse son calibre : il ouvre le circuit et doit être remplacé.
\item Le \cle{disjoncteur} coupe automatiquement le circuit en cas de surcharge ou de court-circuit ; il est \textbf{réarmable}.
\item Le \cle{calibre} (en ampères) doit être adapté à la section des fils et aux appareils de la ligne ; on le choisit grâce à la relation $I = \dfrac{P}{U}$.
\item Fusibles et disjoncteurs protègent \textbf{l'installation}, pas les personnes : la protection des personnes relève du différentiel, objet de la section suivante.
\item Règle d'or : \textbf{jamais} de fusible de calibre supérieur, jamais de fil quelconque à la place d'un fusible.
\end{itemize}
\end{aretenir}

\coursec{La prise de terre et l'interrupteur différentiel}

Dans la section précédente, nous avons vu que le fusible et le disjoncteur protègent les \emph{installations} contre les surcharges et les courts-circuits : ils coupent le circuit quand l'intensité devient trop grande. Mais ces dispositifs ont une limite : un courant de \SI{0,05}{A} suffit à tuer un être humain, alors qu'un fusible de \SI{16}{A} ne fondra jamais pour une si petite fuite ! Il faut donc des dispositifs spécialement conçus pour protéger les \emph{personnes}. Ce sont la prise de terre et l'interrupteur différentiel, que nous allons étudier maintenant.

\courssub{Le défaut d'isolement : quand la carcasse devient dangereuse}

\begin{experience}[Observation : un réfrigérateur qui « donne des coups »]
Dans une maison de quartier, chaque fois que Mama touche la porte métallique du réfrigérateur en étant pieds nus sur le carrelage humide, elle ressent un picotement désagréable. Pourtant, le frigo fonctionne normalement.
\begin{itemize}
\item \textbf{Observation :} la carcasse métallique de l'appareil semble porter une tension électrique.
\item \textbf{Hypothèse :} à l'intérieur de l'appareil, un fil sous tension touche accidentellement la carcasse.
\item \textbf{Vérification :} un électricien mesure une tension entre la carcasse et le sol ; il trouve un fil de phase dont l'isolant, usé, touche la tôle interne.
\item \textbf{Conclusion :} l'appareil présente un \cle{défaut d'isolement}.
\end{itemize}
\end{experience}

\begin{definition}[Défaut d'isolement]
Un \cle{défaut d'isolement} est une panne dans laquelle un conducteur sous tension (la \cle{phase}) touche accidentellement la \cle{carcasse métallique} d'un appareil (réfrigérateur, fer à repasser, machine à laver, bouilloire). La carcasse se trouve alors portée à la tension du réseau (\SI{220}{V}) : toute personne qui la touche risque l'électrisation, voire l'électrocution.
\end{definition}

Pourquoi la personne est-elle traversée par un courant ? Parce que le corps humain, posé sur le sol, relie la carcasse (à \SI{220}{V}) à la terre (à \SI{0}{V}). Le courant emprunte alors le chemin : phase $\rightarrow$ carcasse $\rightarrow$ corps $\rightarrow$ sol. C'est exactement le circuit de l'électrisation étudié dans les sections précédentes.

\courssub{La prise de terre : détourner le courant de fuite}

\begin{definition}[Prise de terre]
La \cle{prise de terre} est un dispositif qui relie les carcasses métalliques des appareils électriques au sol, par l'intermédiaire d'un \cle{conducteur de protection} (fil vert-jaune) connecté à une \cle{tige métallique} (piquet de terre) enfoncée profondément dans le sol. Son rôle est d'évacuer les courants de fuite vers la terre afin qu'ils ne traversent pas le corps humain.
\end{definition}

Le raisonnement physique est le suivant. Le courant électrique est « paresseux » : entre deux chemins possibles, il passe préférentiellement par celui qui offre la plus faible résistance. Or :
\begin{itemize}
\item la résistance du corps humain vaut environ \SI{2000}{\ohm} dans des conditions normales (peau légèrement humide, pieds nus) ;
\item la résistance du conducteur de terre (gros fil de cuivre + piquet) vaut moins de \SI{1}{\ohm} (souvent quelques dixièmes d'ohm).
\end{itemize}
En cas de défaut d'isolement, si la carcasse est reliée à la terre, le courant de fuite s'écoule presque entièrement par le fil de terre vers le sol. La personne qui touche la carcasse n'est traversée que par un courant négligeable : elle est protégée.

\begin{popfigure}
\begin{center}
\begin{tikzpicture}[scale=0.85, font=\small, every node/.style={transform shape}]
% ===== PANNEAU 1 : SANS TERRE =====
\begin{scope}
% sol
\fill[poppinkL] (-3.5,-0.5) rectangle (3.5,-1.0);
\draw[popdark, thick] (-3.5,-0.5) -- (3.5,-0.5);
\node[popdark] at (2.5,-0.75) {\footnotesize sol (terre)};
% carcasse (frigo)
\draw[fill=popmuted!40, thick] (-3.0,-0.5) rectangle (-1.0,2.5);
\node[align=center] at (-2.0,1.8) {\footnotesize carcasse\\ \footnotesize métallique};
% fil phase en défaut
\draw[popink, very thick] (-2.0,3.5) -- (-2.0,2.5);
\draw[popink, thick] (-2.0,3.0) -- (-1.3,2.5);
\node[popink] at (-2.0,3.7) {\footnotesize Phase 220 V};
\fill[popink] (-1.3,2.5) circle (0.08);
\node[popink, align=center] at (-0.3,2.8) {\footnotesize défaut\\ \footnotesize d'isolement};
% personne
\draw[thick] (1.2,2.0) circle (0.3);
\draw[thick] (1.2,1.7) -- (1.2,0.6);
\draw[thick] (1.2,1.5) -- (0.6,1.1);
\draw[thick] (0.6,1.1) -- (-0.8,1.4);
\draw[thick] (1.2,0.6) -- (0.8,-0.5);
\draw[thick] (1.2,0.6) -- (1.6,-0.5);
% courant à travers le corps
\draw[->, poppink, very thick] (-0.75,1.45) -- (0.4,1.15);
\draw[->, poppink, very thick] (0.7,1.0) -- (1.0,0.3);
\draw[->, poppink, very thick] (1.2,0.1) -- (1.2,-0.6);
\node[poppink, align=center] at (2.5,1.0) {\footnotesize courant à\\ \footnotesize travers le corps};
\node[font=\bfseries, popdark] at (0,-1.5) {SANS prise de terre : DANGER};
\end{scope}
% ===== PANNEAU 2 : AVEC TERRE =====
\begin{scope}[xshift=9.5cm]
\fill[popgreenL] (-3.5,-0.5) rectangle (3.5,-1.0);
\draw[popdark, thick] (-3.5,-0.5) -- (3.5,-0.5);
\node[popdark] at (2.5,-0.75) {\footnotesize sol (terre)};
\draw[fill=popmuted!40, thick] (-3.0,-0.5) rectangle (-1.0,2.5);
\node[align=center] at (-2.0,1.8) {\footnotesize carcasse\\ \footnotesize métallique};
\draw[popink, very thick] (-2.0,3.5) -- (-2.0,2.5);
\draw[popink, thick] (-2.0,3.0) -- (-1.3,2.5);
\node[popink] at (-2.0,3.7) {\footnotesize Phase 220 V};
\fill[popink] (-1.3,2.5) circle (0.08);
% conducteur de terre
\draw[popgreen, very thick] (-3.0,0.8) -- (-3.6,0.8) -- (-3.6,-0.95);
\draw[popgreen, thick] (-3.85,-0.95) -- (-3.35,-0.95);
\draw[popgreen, thick] (-3.8,-1.1) -- (-3.4,-1.1);
\draw[popgreen, thick] (-3.72,-1.25) -- (-3.48,-1.25);
\node[popgreen, align=center] at (-3.6,1.4) {\footnotesize fil de\\ \footnotesize terre (vert-jaune)};
% courant de fuite vers la terre
\draw[->, popgreen, very thick] (-2.95,0.75) -- (-3.55,0.75);
\draw[->, popgreen, very thick] (-3.6,0.6) -- (-3.6,-0.7);
% personne
\draw[thick] (1.2,2.0) circle (0.3);
\draw[thick] (1.2,1.7) -- (1.2,0.6);
\draw[thick] (1.2,1.5) -- (0.6,1.1);
\draw[thick] (0.6,1.1) -- (-0.8,1.4);
\draw[thick] (1.2,0.6) -- (0.8,-0.5);
\draw[thick] (1.2,0.6) -- (1.6,-0.5);
\node[popgreen, align=center] at (2.3,1.0) {\footnotesize personne\\ \footnotesize protégée};
\node[font=\bfseries, popdark] at (0,-1.5) {AVEC prise de terre : SÉCURITÉ};
\end{scope}
\end{tikzpicture}
\end{center}
\legende{Comparaison : sans prise de terre, le courant de fuite traverse le corps de la personne vers le sol ; avec la prise de terre, il s'écoule par le conducteur de protection (vert-jaune), dont la résistance est très inférieure à celle du corps.}
\end{popfigure}

\begin{methode}[Vérifier la présence d'une prise de terre chez soi]
\begin{enumerate}
\item Observer les prises murales : une prise avec terre possède une \textbf{broche métallique} (ou deux ergots latéraux) en plus des deux trous.
\item Vérifier que les appareils à carcasse métallique (frigo, machine à laver, fer) ont une fiche à \textbf{trois conducteurs}.
\item Repérer, près du compteur ou du tableau électrique, le fil vert-jaune descendant vers un piquet enfoncé dans le sol.
\item En cas de doute, faire appel à un électricien qualifié : ne jamais bricoler soi-même l'installation.
\end{enumerate}
\end{methode}

\courssub{L'interrupteur différentiel : le gardien des personnes}

La prise de terre détourne le courant de fuite, mais elle ne coupe pas le circuit. Il faut un deuxième dispositif, capable de \emph{détecter} la fuite et de \emph{couper} automatiquement le courant : c'est l'interrupteur différentiel.

\begin{definition}[Interrupteur différentiel]
L'\cle{interrupteur différentiel} est un dispositif de sécurité qui compare en permanence l'intensité du courant qui \emph{part} par le fil de phase et celle du courant qui \emph{revient} par le fil de neutre. Dès qu'il détecte une différence (une \cle{fuite de courant}) supérieure à sa \cle{sensibilité}, il coupe automatiquement le circuit. Pour la protection des personnes, on utilise un différentiel de sensibilité \SI{30}{mA}, appelé \cle{haute sensibilité}.
\end{definition}

\begin{experience}[Principe de fonctionnement du différentiel]
\textbf{Fonctionnement normal.} Dans un circuit sain, tout le courant qui part par la phase revient par le neutre : $I_{\text{phase}} = I_{\text{neutre}}$. Les deux fils passent à l'intérieur d'un anneau magnétique (tore) ; leurs effets magnétiques, égaux et opposés, s'annulent : rien ne se passe.

\textbf{En cas de fuite.} Si une personne touche une carcasse en défaut, une partie du courant s'échappe vers la terre (à travers le corps ou le fil de terre) et ne revient \emph{pas} par le neutre : $I_{\text{phase}} > I_{\text{neutre}}$. Le tore détecte ce déséquilibre, un électroaimant déclenche un mécanisme qui ouvre les contacts : le circuit est coupé.

\textbf{Conclusion :} le différentiel ne mesure pas le courant de fuite directement ; il détecte la \emph{différence} entre le courant aller et le courant retour, d'où son nom.
\end{experience}

\begin{popfigure}
\begin{center}
\begin{tikzpicture}[scale=1.0, font=\small]
% tore magnétique
\draw[fill=popgoldL, thick, popgold!80!black] (0,0) circle (1.5);
\draw[fill=white, thick, popgold!80!black] (0,0) circle (0.95);
\node[popgold!60!black] at (0,1.85) {\footnotesize tore magnétique};
% fil phase
\draw[popink, very thick] (-4.5,0.35) -- (-0.35,0.35);
\draw[popink, very thick] (0.35,0.35) -- (2.2,0.35);
\draw[popink, very thick] (2.2,0.35) -- (2.2,1.1) -- (3.4,1.1);
\node[popink] at (-4.5,0.65) {\footnotesize Phase};
\node[popink] at (0,0.62) {\footnotesize $I_{\text{phase}}$};
% fil neutre
\draw[popblue, very thick] (-4.5,-0.35) -- (-0.35,-0.35);
\draw[popblue, very thick] (0.35,-0.35) -- (2.2,-0.35);
\draw[popblue, very thick] (2.2,-0.35) -- (2.2,-1.1) -- (3.4,-1.1);
\node[popblue] at (-4.5,-0.65) {\footnotesize Neutre};
\node[popblue] at (0,-0.62) {\footnotesize $I_{\text{neutre}}$};
% récepteur
\draw[fill=popmuted!40, thick] (3.4,0.5) rectangle (4.9,-0.5);
\node at (4.15,0) {\footnotesize récepteur};
\draw[popink, very thick] (3.4,1.1) -- (4.0,1.1) -- (4.0,0.5);
\draw[popblue, very thick] (3.4,-1.1) -- (4.3,-1.1) -- (4.3,-0.5);
% fuite
\draw[->, poppink, very thick] (4.6,-0.5) -- (4.6,-1.5);
\node[poppink, align=left] at (5.6,-1.0) {\footnotesize fuite vers\\ \footnotesize la terre};
% bobine de détection
\draw[popgreen, thick] (0,-1.5) -- (0,-2.1);
\draw[popgreen, thick] (-0.3,-2.1) rectangle (0.3,-2.7);
\node[popgreen, align=center] at (0,-3.05) {\footnotesize détecteur de déséquilibre};
% contacts
\draw[thick] (-3.4,0.35) -- (-2.9,0.35);
\draw[thick] (-2.9,0.35) -- (-2.55,0.75);
\fill (-3.4,0.35) circle (0.06);
\fill (-2.55,0.35) circle (0.06);
\node[align=center] at (-2.95,1.15) {\footnotesize contacts (coupure)};
\draw[->, popgreen, dashed] (-0.3,-2.4) -- (-2.55,0.6);
\end{tikzpicture}
\end{center}
\legende{Principe de l'interrupteur différentiel : les fils de phase et de neutre traversent un tore magnétique. Si $I_{\text{phase}} \neq I_{\text{neutre}}$ (fuite de courant), le déséquilibre est détecté et le mécanisme ouvre les contacts en moins de \SI{0,1}{s}.}
\end{popfigure}

\begin{propriete}[Protection des personnes par le différentiel 30 mA (admise)]
Un interrupteur différentiel de sensibilité \SI{30}{mA} coupe le circuit en moins de \SI{0,1}{s} dès que le courant de fuite dépasse \SI{30}{mA}. Comme la fibrillation cardiaque n'apparaît qu'au-delà d'environ \SI{75}{mA} maintenus pendant plus d'une seconde, cette coupure quasi instantanée empêche la fibrillation cardiaque dans la plupart des cas : le différentiel \SI{30}{mA} protège efficacement les personnes.
\end{propriete}

\begin{exemplebox}[Le bouton « test »]
Tout interrupteur différentiel possède un \cle{bouton test} (souvent marqué « T »). En l'actionnant, on simule volontairement une petite fuite de courant : le différentiel doit \textbf{déclencher immédiatement}. Il est recommandé d'appuyer sur ce bouton \textbf{une fois par mois} pour vérifier le bon fonctionnement du dispositif. S'il ne déclenche pas, il faut le faire remplacer sans attendre.
\end{exemplebox}

\courssub{Terre et différentiel : une équipe inséparable}

\begin{aretenir}[La protection complète des personnes]
La protection des personnes exige l'\textbf{association des deux dispositifs} :
\begin{itemize}
\item la \cle{prise de terre} \textbf{canalise} le courant de fuite vers le sol ;
\item l'\cle{interrupteur différentiel} \SI{30}{mA} \textbf{détecte} cette fuite et \textbf{coupe} le circuit en moins de \SI{0,1}{s}.
\end{itemize}
L'un sans l'autre, la protection est incomplète.
\end{aretenir}

\begin{attention}[Pièges fréquents]
\begin{itemize}
\item \textbf{La prise de terre seule ne coupe pas le circuit.} Sans différentiel, une fuite faible (par exemple \SI{100}{mA}) peut durer indéfiniment : la carcasse reste dangereuse et l'installation peut s'échauffer.
\item \textbf{Le différentiel ne remplace pas la terre.} Si la carcasse n'est pas reliée à la terre, aucune fuite n'existe \emph{avant} que quelqu'un ne la touche : le différentiel ne déclenchera qu'au moment où le courant traversera déjà la personne. La terre permet de détecter le défaut \emph{avant} l'accident.
\item Ne confonds pas les rôles : le \textbf{disjoncteur} (ou fusible) protège l'\emph{installation} (surcharges, courts-circuits) ; le \textbf{différentiel} protège les \emph{personnes} (fuites de courant).
\end{itemize}
\end{attention}

\courssub{Règles de sécurité et conduite à tenir}

\begin{aretenir}[Règles d'or de sécurité électrique à la maison]
\begin{enumerate}
\item \textbf{Mains sèches, pieds secs :} ne jamais manipuler un appareil électrique, une prise ou un interrupteur avec les mains mouillées ou les pieds nus sur un sol humide.
\item \textbf{Tirer sur la fiche, pas sur le fil :} pour débrancher, saisir la fiche fermement ; tirer sur le fil abîme l'isolant et crée un risque de court-circuit.
\item \textbf{Eau et électricité ne font pas bon ménage :} éloigner tout appareil branché (radio, téléphone en charge, rallonge) de l'évier, de la baignoire, du seau d'eau.
\item \textbf{Ne pas surcharger les prises :} éviter les « pieuvres » (multiprises en cascade) qui font chauffer les fils et risquent l'incendie.
\item \textbf{Câbles abîmés = danger :} un fil dénudé, une fiche cassée, une rallonge écrasée sous un meuble doivent être remplacés immédiatement.
\item \textbf{En cas d'accident (électrisation) :} \textbf{ne pas toucher la victime} tant qu'elle est en contact avec le courant. 1) Couper le courant (disjoncteur général ou débrancher). 2) Appeler les secours (112 ou 118). 3) Si formation secourisme : pratiquer les gestes de premiers secours.
\end{enumerate}
\end{aretenir}

\begin{exoresolu}[Une machine à laver en défaut]
\textbf{Énoncé.} La carcasse métallique d'une machine à laver est accidentellement reliée à la phase (\SI{220}{V}). L'installation possède une prise de terre de résistance totale \SI{50}{\ohm} et un interrupteur différentiel de \SI{30}{mA}.
\begin{enumerate}
\item Calculer l'intensité du courant de fuite qui s'écoule vers la terre.
\item Le différentiel déclenche-t-il ? Justifier.
\item Que se passerait-il si la maison n'avait ni terre ni différentiel et qu'une personne touche la carcasse ?
\end{enumerate}
\tcblower
\textbf{Solution détaillée.}
\begin{enumerate}
\item D'après la loi d'Ohm, le courant de fuite vaut :
\[ I = \frac{U}{R} = \frac{220}{50} = \SI{4,4}{A}. \]
\item Oui : $4{,}4\ \text{A} = 4400\ \text{mA} \gg 30\ \text{mA}$. Le différentiel détecte cette fuite et coupe le circuit en moins de \SI{0,1}{s}. La machine est mise hors tension : plus aucun danger.
\item Sans terre ni différentiel, la carcasse reste à \SI{220}{V} sans que rien ne le signale. La personne qui la touche est traversée par un courant $I = \dfrac{220}{2000} \approx \SI{110}{mA}$ (corps de résistance \SI{2000}{\ohm}), valeur supérieure au seuil de fibrillation cardiaque (\SI{75}{mA}) : risque d'électrocution mortelle.
\end{enumerate}
\end{exoresolu}

\begin{savaistu}[Et chez nous, au Cameroun ?]
De nombreuses habitations anciennes, à Yaoundé, Douala ou dans les villages, ont été câblées \textbf{sans prise de terre} : les prises murales n'ont que deux trous, et les carcasses des appareils ne sont reliées à rien. C'est une des causes fréquentes d'accidents domestiques, surtout pendant la saison des pluies quand les sols sont humides. En attendant la mise aux normes de l'installation par un électricien, applique les réflexes de survie : ne jamais toucher un appareil électrique les \textbf{mains mouillées} ou \textbf{pieds nus}, débrancher par la \textbf{fiche} (jamais en tirant sur le fil), et éloigner tout appareil branché des points d'eau (cuisine, salle de bain).
\end{savaistu}

\begin{exercice}[Pour t'entraîner]
\begin{enumerate}
\item Définis la prise de terre et explique son rôle en cas de défaut d'isolement.
\item Explique, à l'aide des intensités $I_{\text{phase}}$ et $I_{\text{neutre}}$, comment l'interrupteur différentiel détecte une fuite de courant.
\item Pourquoi dit-on que la terre et le différentiel forment une « équipe » ? Que risque-t-on si l'un des deux manque ?
\item À quoi sert le bouton « test » du différentiel et à quelle fréquence faut-il l'utiliser ?
\item Cite trois règles de sécurité à respecter lors de l'utilisation d'une rallonge électrique à la maison.
\end{enumerate}
\end{exercice}

\begin{corrige}[Exercice « Pour t'entraîner »]
\begin{enumerate}
\item La prise de terre relie les carcasses métalliques des appareils au sol par un conducteur de protection (fil vert-jaune) et un piquet enfoncé dans la terre. En cas de défaut d'isolement, le courant de fuite s'écoule vers le sol par ce conducteur, de résistance très faible, au lieu de traverser le corps de la personne.
\item En fonctionnement normal, $I_{\text{phase}} = I_{\text{neutre}}$. Si une fuite existe (vers la terre ou le corps d'une personne), une partie du courant ne revient pas par le neutre : $I_{\text{phase}} > I_{\text{neutre}}$. Le tore magnétique détecte ce déséquilibre et déclenche l'ouverture des contacts.
\item La terre canalise la fuite vers le sol, ce qui crée un déséquilibre détectable ; le différentiel détecte ce déséquilibre et coupe le circuit en moins de \SI{0,1}{s}. Sans terre, le différentiel ne déclenche qu'au moment où la personne est déjà traversée par le courant ; sans différentiel, la fuite n'est jamais coupée et peut rester dangereuse.
\item Le bouton « test » simule une fuite de courant pour vérifier que le différentiel déclenche correctement. Il faut l'actionner environ une fois par mois.
\item Exemples : ne pas enrouler la rallonge (effet bobine/échauffement) ; ne pas la faire passer sous un tapis ou une porte (écrasement/usure) ; ne pas brancher d'appareils trop puissants (four, fer, chauffage) sur une rallonge fine ; vérifier l'état de l'isolant avant usage.
\end{enumerate}
\end{corrige}

\coursec{Règles de sécurité et conduite à tenir}

Dans la section précédente, nous avons découvert deux gardiens silencieux de nos maisons : la \cle{prise de terre}, qui offre au courant de fuite un chemin vers le sol, et l'\cle{interrupteur différentiel}, qui coupe le circuit dès qu'il détecte une fuite de courant supérieure à \SI{30}{mA}. Ces dispositifs sont formidables, mais ils ont une limite : ils ne peuvent pas tout. Ils n'empêchent pas un enfant d'introduire un clou dans une prise, ni une maman de brancher cinq appareils sur la même multiprise, ni un papa de réparer un fil dénudé sous tension. La meilleure protection, c'est encore le \cle{comportement} de chacun. Cette dernière section est peut-être la plus importante de toute l'année : elle peut sauver ta vie et celle de tes proches.

\courssub{Les règles de prévention à la maison}

\courssub{Eau et électricité : un couple dangereux}

Commençons par une question simple : pourquoi les notices des sèche-cheveux et des bouilloires répètent-elles sans cesse « ne pas utiliser près d'une baignoire » ? Parce que l'eau transforme notre corps en excellent conducteur.

\begin{attention}[L'eau courante n'est pas isolante]
L'eau \emph{chimiquement pure} conduit très mal le courant. Mais l'eau du robinet, de la forage ou de la pluie contient des \cle{sels minéraux} dissous (ions sodium, calcium, chlorure...). Ces ions la rendent \textbf{conductrice}. Une main mouillée diminue fortement la résistance du corps humain : le courant qui traverse le corps devient alors beaucoup plus intense, donc beaucoup plus dangereux, pour la même tension de \SI{220}{V}.
\end{attention}

Rappelons la loi d'Ohm appliquée au corps humain : $I = \dfrac{U}{R}$. Si la tension $U = \SI{220}{V}$ est fixe, plus la résistance $R$ du corps est faible, plus l'intensité $I$ est grande. Or une peau sèche présente une résistance élevée, tandis qu'une peau mouillée la fait chuter. C'est toute la différence entre une simple « châtaigne » et une électrocution.

\begin{aretenir}[Règles d'or avec l'eau]
\begin{itemize}
\item Ne jamais toucher un appareil électrique (interrupteur, prise, fer à repasser, réfrigérateur) avec les \cle{mains mouillées} ou les \cle{pieds nus} sur un sol humide.
\item Ne jamais utiliser ou brancher un appareil électrique dans la \cle{salle de bain} (sèche-cheveux, radio, chargeur de téléphone) à proximité de l'eau.
\item Se sécher soigneusement les mains avant toute manipulation électrique.
\end{itemize}
\end{aretenir}

\courssub{Multiprises, fils dénudés et bricolage interdit}

Dans beaucoup de maisons, on voit une multiprise sur laquelle sont branchés le téléviseur, le décodeur, le ventilateur, le fer à repasser et parfois... une deuxième multiprise ! C'est la \cle{surcharge} : la somme des intensités dépasse ce que le fil peut supporter, le fil chauffe, l'isolant fond, et l'incendie peut démarrer.

\begin{aretenir}[Règles d'or pour les installations]
\begin{itemize}
\item Ne pas \cle{surcharger} les multiprises : on additionne les puissances des appareils branchés et on ne dépasse pas la limite indiquée sur la multiprise.
\item Ne jamais réparer soi-même un \cle{fil dénudé} avec un simple ruban adhésif de fortune : faire appel à un \cle{électricien qualifié}.
\item Ne jamais introduire d'\cle{objet métallique} (fourchette, clou, trombone) dans une prise de courant.
\item Protéger les prises par des \cle{cache-prises} lorsqu'il y a de jeunes enfants à la maison.
\item Avant toute intervention sur l'installation (changer une douille, réparer une prise), \cle{couper le courant} au disjoncteur général.
\end{itemize}
\end{aretenir}

\begin{popfigure}
\begin{tikzpicture}[scale=0.9, every node/.style={font=\small}]
% --- Pictogramme 1 : mains mouillées + frigo barré ---
\begin{scope}[shift={(0,0)}]
\draw[thick, popblue, fill=popblueL] (-1.1,-0.9) rectangle (1.1,0.9);
\draw[thick, popdark] (-0.55,-0.5) rectangle (0.35,0.55); % frigo
\draw[thick, popdark] (0.35,-0.5) -- (0.35,0.55);
\draw[popdark] (0.28,0.35) -- (0.28,0.1);
\draw[thick, popdark, fill=popcream] (0.75,0.15) ellipse (0.22 and 0.3); % main
\draw[popblue, thick] (0.75,-0.55) -- (0.75,-0.25); % goutte tige
\draw[popblue, fill=popblue] (0.75,-0.62) circle (0.06);
\draw[popblue, fill=popblue] (0.55,-0.45) circle (0.05);
\draw[very thick, poppink] (-1.0,-0.8) -- (1.0,0.8);
\draw[very thick, poppink] (-1.0,0.8) -- (1.0,-0.8);
\node[below, font=\footnotesize\bfseries, align=center] at (0,-1.05) {Mains mouillées :\\ INTERDIT};
\end{scope}
% --- Pictogramme 2 : multiprise en cascade barrée ---
\begin{scope}[shift={(5.2,0)}]
\draw[thick, popblue, fill=popblueL] (-1.1,-0.9) rectangle (1.1,0.9);
\draw[thick, popdark, fill=white] (-0.9,0.35) rectangle (0.5,0.6); % multiprise 1
\draw[popdark] (-0.6,0.475) circle (0.05); \draw[popdark] (-0.2,0.475) circle (0.05); \draw[popdark] (0.2,0.475) circle (0.05);
\draw[thick, popdark] (0.5,0.475) .. controls (0.9,0.4) and (0.9,-0.1) .. (0.5,-0.2); % cable
\draw[thick, popdark, fill=white] (-0.9,-0.2) rectangle (0.5,-0.45); % multiprise 2
\draw[popdark] (-0.6,-0.325) circle (0.05); \draw[popdark] (-0.2,-0.325) circle (0.05); \draw[popdark] (0.2,-0.325) circle (0.05);
\draw[very thick, poppink] (-1.0,-0.8) -- (1.0,0.8);
\draw[very thick, poppink] (-1.0,0.8) -- (1.0,-0.8);
\node[below, font=\footnotesize\bfseries, align=center] at (0,-1.05) {Multiprises en cascade :\\ INTERDIT};
\end{scope}
% --- Pictogramme 3 : coupure disjoncteur validée ---
\begin{scope}[shift={(10.4,0)}]
\draw[thick, popgreen, fill=popgreenL] (-1.1,-0.9) rectangle (1.1,0.9);
\draw[thick, popdark, fill=white] (-0.35,-0.55) rectangle (0.35,0.55); % disjoncteur
\draw[thick, popdark, fill=popmuted] (-0.12,-0.05) rectangle (0.12,0.3); % levier baissé
\draw[thick, popdark] (0,0.3) -- (0,0.45);
\node[font=\footnotesize, popdark] at (0,-0.35) {OFF};
\draw[very thick, popgreen] (-0.75,-0.15) -- (-0.45,-0.5) -- (0.75,0.65); % coche
\node[below, font=\footnotesize\bfseries, align=center] at (0,-1.05) {Couper le disjoncteur\\ avant d'intervenir : OUI};
\end{scope}
\end{tikzpicture}
\legende{Pictogrammes de sécurité : 1) jamais d'appareil électrique avec les mains mouillées ; 2) jamais de multiprises branchées en cascade ; 3) toujours couper le disjoncteur général avant toute intervention.}
\end{popfigure}

\courssub{Conduite à tenir face à une victime d'électrisation}

Imaginons une scène malheureusement fréquente : pendant la saison des pluies à Douala, un voisin touche un fil dénudé de son groupe électrogène. Il crie, ses muscles se contractent, il ne peut pas lâcher le fil. Que faire ? Agir vite, oui — mais agir \emph{dans le bon ordre}.

\begin{attention}[L'erreur mortelle à ne jamais commettre]
L'erreur la plus fréquente — et la plus mortelle — est de \textbf{toucher la victime directement} pour la tirer en arrière alors qu'elle est encore en contact avec le conducteur. Le courant traverse alors le corps du secouriste : \textbf{le sauveteur s'électrise à son tour} et il y a deux victimes au lieu d'une. Tant que le courant n'est pas coupé, la victime est un conducteur comme un autre.
\end{attention}

\begin{methode}[Secourir une victime d'électrisation : les 4 étapes]
\begin{enumerate}
\item \textbf{Couper le courant} immédiatement : disjoncteur général, interrupteur, ou débrancher la prise si c'est possible sans danger.
\item \textbf{Ne jamais toucher la victime} tant qu'elle est en contact avec le conducteur. Si l'on ne peut pas couper le courant, écarter la victime avec un \cle{objet isolant et sec} : bâton de bois sec, manche en plastique, corde sèche.
\item \textbf{Alerter les secours} : au Cameroun, composer le \cle{119} (pompiers) ou le SAMU, en précisant le lieu exact et la nature de l'accident.
\item \textbf{Surveiller la respiration} de la victime, la placer en position latérale de sécurité si elle respire, et pratiquer les gestes de premiers secours (massage cardiaque) \emph{uniquement si l'on est formé}.
\end{enumerate}
\end{methode}

\begin{popfigure}
\begin{tikzpicture}[
node distance=0.55cm,
etape/.style={rectangle, rounded corners=3pt, draw=popblue, thick, fill=popblueL, text width=6.2cm, align=center, font=\small, inner sep=5pt},
alerte/.style={rectangle, rounded corners=3pt, draw=poporange, thick, fill=poporangeL, text width=6.2cm, align=center, font=\small, inner sep=5pt},
danger/.style={rectangle, rounded corners=3pt, draw=poppink, thick, fill=poppinkL, text width=6.2cm, align=center, font=\small, inner sep=5pt},
fleche/.style={-{Stealth[length=2.5mm]}, very thick, popdark}]
\node[etape, align=center] (e1){\textbf{1. COUPER le courant}\\ (disjoncteur général, prise)};
\node[danger, below=of e1, align=center] (e2){\textbf{2. NE PAS TOUCHER la victime}\\ tant qu'elle est en contact\\ (sinon : l'écarter avec un objet isolant sec)};
\node[alerte, below=of e2, align=center] (e3){\textbf{3. ALERTER les secours}\\ Pompiers / SAMU : \textbf{119}};
\node[etape, below=of e3, align=center] (e4){\textbf{4. SECOURIR et SURVEILLER}\\ respiration, position latérale de sécurité,\\ premiers secours si formé};
\draw[fleche] (e1) -- (e2);
\draw[fleche] (e2) -- (e3);
\draw[fleche] (e3) -- (e4);
\end{tikzpicture}
\legende{Organigramme de la conduite à tenir face à une victime d'électrisation : couper, écarter sans contact direct, alerter, secourir. L'ordre des étapes est vital.}
\end{popfigure}

\begin{exemplebox}[Le différentiel, un gardien plus rapide que le cœur]
Un interrupteur différentiel de sensibilité \SI{30}{mA} déclenche en moins de \SI{0,1}{s}. Or un battement de cœur dure environ \SI{0,8}{s}. Le différentiel coupe donc le circuit \textbf{huit fois plus vite qu'un battement de cœur} : c'est ce qui lui permet, dans la plupart des cas, d'interrompre le courant avant que la fibrillation cardiaque ne s'installe. Mais attention : il ne protège pas contre tout, et il ne dispense jamais des règles de prudence.
\end{exemplebox}

\courssub{Les dangers hors de la maison}

L'électricité ne menace pas seulement entre les quatre murs de la maison. Dehors, les installations d'\cle{ENEO} (lignes, poteaux, transformateurs) présentent des risques spécifiques.

\begin{aretenir}[Règles d'or à l'extérieur]
\begin{itemize}
\item Ne jamais s'approcher d'un \cle{câble électrique tombé au sol} : le sol autour du câble peut être sous tension, surtout s'il est mouillé. Le courant peut remonter par une jambe et redescendre par l'autre si l'on marche à grands pas : c'est le danger dit du « \cle{pas de pas} ». Il faut rester à distance (au moins \SI{10}{m}) et prévenir ENEO ou les pompiers.
\item Ne jamais jouer, grimper ou faire sécher du linge près des \cle{transformateurs} et des poteaux électriques.
\item Ne jamais faire voler un cerf-volant ou lancer un objet près des lignes électriques.
\end{itemize}
\end{aretenir}

\begin{savaistu}[Le 220 V tue plus que la haute tension]
On croit souvent que seule la haute tension des lignes ENEO est dangereuse. C'est faux : la grande majorité des accidents électriques mortels surviennent avec la tension « ordinaire » du secteur, \SI{220}{V}, celle des prises de la maison. Pourquoi ? Parce que c'est avec elle que nous sommes en contact chaque jour, souvent dans de mauvaises conditions (mains mouillées, installations vétustes, bricolages). La haute tension, elle, tue souvent... sans même qu'on la touche, par simple proximité. La leçon : \textbf{aucune tension domestique n'est anodine}.
\end{savaistu}

\courssub{Pour aller plus loin : la foudre, une électrisation naturelle (hors programme)}

\begin{savaistu}[La foudre et le paratonnerre — culture pour le Bac]
La \cle{foudre} est une gigantesque électrisation naturelle : une décharge de plusieurs millions de volts entre un nuage et le sol. Un coup de foudre direct est presque toujours mortel. Pour protéger les bâtiments, on installe un \cle{paratonnerre} : une tige métallique pointue placée au sommet du toit et reliée à une grosse tige enterrée dans le sol. Son rôle est exactement analogue à celui d'une \textbf{prise de terre géante} : il offre au courant de la foudre un chemin direct et facile vers le sol, en évitant qu'il ne traverse la maison... ou ses habitants. Pendant un orage, on évite donc de s'abriter sous un arbre isolé, de tenir un objet métallique ou de téléphoner avec un appareil branché.
\end{savaistu}

\courssub{Exercice résolu et bilan}

\begin{exoresolu}[Analyser une situation accidentelle]
\textbf{Énoncé.} Pendant la saison des pluies, Marlyse voit son petit frère toucher le fil dénudé d'une rallonge branchée ; l'enfant crie et ne peut pas lâcher le fil. Marlyse, paniquée, veut le tirer par le bras. 
\begin{enumerate}
\item Pourquoi ce geste est-il dangereux pour Marlyse ?
\item Décris, dans l'ordre, les trois premières actions que Marlyse doit accomplir.
\item L'installation de la maison est protégée par un différentiel de \SI{30}{mA} qui coupe en \SI{0,08}{s}. Pourquoi ce dispositif augmente-t-il les chances de survie de l'enfant ?
\end{enumerate}
\tcblower
\textbf{Solution détaillée.}
\begin{enumerate}
\item Tant que l'enfant est en contact avec le fil sous tension, son corps est traversé par le courant : il se comporte comme un conducteur. En le touchant, Marlyse fermerait un nouveau circuit à travers son propre corps vers le sol : elle s'électriserait à son tour.
\item Marlyse doit : \textbf{1)} couper le courant au disjoncteur général ; \textbf{2)} si la coupure est impossible, écarter l'enfant avec un objet isolant et sec (bâton de bois sec) sans jamais le toucher directement ; \textbf{3)} alerter les secours en composant le 119, puis surveiller la respiration de l'enfant.
\item Le différentiel détecte la fuite de courant vers la terre et coupe le circuit en \SI{0,08}{s}, soit environ dix fois plus vite qu'un battement de cœur. La durée de passage du courant dans le corps de l'enfant est ainsi très courte, ce qui réduit fortement le risque de fibrillation cardiaque et de brûlures graves.
\end{enumerate}
\end{exoresolu}

\begin{exercice}[À toi de jouer]
\begin{enumerate}
\item Cite trois règles de sécurité à respecter dans une salle de bain équipée d'une prise électrique.
\item Pourquoi est-il dangereux de brancher une multiprise sur une autre multiprise ?
\item Un câble ENEO est tombé sur la route devant chez toi après une pluie. Que fais-tu, et pourquoi ne faut-il pas s'en approcher ?
\end{enumerate}
\end{exercice}

\begin{corrige}[Exercice]
\begin{enumerate}
\item Ne pas utiliser d'appareil électrique près de l'eau ; ne jamais toucher la prise ou l'interrupteur avec les mains mouillées ou les pieds nus ; ne pas laisser un appareil branché (sèche-cheveux) à portée de la baignoire ou du lavabo.
\item Brancher des multiprises en cascade provoque une surcharge : l'intensité totale dépasse ce que le fil peut supporter, le fil s'échauffe, l'isolant peut fondre et provoquer un court-circuit ou un incendie.
\item Je reste à distance (au moins \SI{10}{m}), j'empêche les autres de s'approcher et j'alerte ENEO ou les pompiers (119). Le sol mouillé autour du câble peut être sous tension : en marchant, le courant traverserait mon corps d'une jambe à l'autre (danger du pas de pas).
\end{enumerate}
\end{corrige}

\begin{aretenir}[L'essentiel de la leçon]
\begin{itemize}
\item L'eau courante est conductrice : jamais d'électricité avec les mains mouillées, les pieds nus ou près d'un point d'eau.
\item Pas de surcharge des multiprises, pas de bricolage électrique : on coupe le disjoncteur général avant toute intervention et on fait appel à un électricien qualifié.
\item Face à une victime d'électrisation : \textbf{couper le courant, ne jamais toucher directement, alerter le 119, surveiller et secourir}.
\item Dehors : on ne s'approche jamais d'un câble tombé au sol ni d'un transformateur.
\item Le \SI{220}{V} domestique est la première cause des accidents électriques mortels : la prudence est la meilleure protection.
\end{itemize}
\end{aretenir}

\newpage

\coursec{Activité d'intégration}

\courssub{Objectifs de l'activité}

À la fin de cette activité d'intégration, tu seras capable de :

\begin{itemize}
\item mobiliser l'ensemble des savoirs de la leçon (effets du courant sur le corps humain, surcharge, court-circuit, fusible, disjoncteur, prise de terre, interrupteur différentiel) pour analyser une situation réelle ;
\item lire et exploiter un dossier documentaire : récit d'accident, description d'installation, tableau de données, plaques signalétiques ;
\item effectuer un calcul simple d'intensité à partir des puissances des appareils ($I = \dfrac{P}{U}$) et conclure sur une surcharge ;
\item identifier les défauts d'une installation électrique domestique et proposer des corrections conformes aux règles de sécurité ;
\item rédiger un court rapport argumenté, comme le ferait un technicien expert.
\end{itemize}

\begin{aretenir}[Compétence visée]
« Expertiser » une installation, c'est raisonner comme un professionnel : observer les faits, chercher les causes, vérifier par le calcul, puis proposer des solutions. C'est exactement la démarche d'investigation que tu as pratiquée toute l'année : \textbf{observation} $\rightarrow$ \textbf{hypothèse} $\rightarrow$ \textbf{vérification} $\rightarrow$ \textbf{conclusion}.
\end{aretenir}

\courssub{Situation}

\textbf{Un matin à Bonabéri (Douala)...}

Monsieur Essomba habite une maison à Bonabéri, quartier de Douala. Un samedi matin, alors qu'il touche la carcasse métallique de sa machine à laver en marche, il ressent des \cle{picotements} désagréables dans la main et le bras. Il retire vivement sa main. Sa fille Aïcha, élève en classe de 3ème, l'empêche de retoucher l'appareil et déclare : « Papa, c'est un problème d'électricité, il ne faut plus y toucher ! » La famille fait appel à un technicien, et toi, avec ton groupe, tu es chargé d'établir le \textbf{rapport d'expertise} de l'installation.

\textbf{Document 1 : description de l'installation}

L'habitation est alimentée par le réseau ENEO sous une tension $U = \SI{220}{V}$. L'installation présente les caractéristiques suivantes :

\begin{itemize}
\item la maison ne possède \textbf{pas de prise de terre} (aucun piquet métallique planté dans le sol, aucun conducteur vert-jaune) ;
\item l'installation ne comporte \textbf{pas d'interrupteur différentiel} ;
\item une seule ligne part du compteur vers la cuisine et la buanderie ; elle est protégée par un \textbf{fusible de calibre $\SI{32}{A}$} ;
\item sur cette ligne, une \textbf{multiprise} est branchée, et sur cette multiprise est branchée une \textbf{seconde multiprise} (branchement « en cascade ») qui alimente simultanément : le réfrigérateur, le fer à repasser et la bouilloire électrique ;
\item la machine à laver, dont la carcasse est métallique, est branchée sur une prise murale à deux broches, sans borne de terre ;
\item le fil de la bouilloire est dénudé sur quelques centimètres près de la fiche.
\end{itemize}

\textbf{Document 2 : plaques signalétiques des appareils}

\begin{center}
\begin{tabularx}{\linewidth}{|l|X|X|}
\hline
\rowcolor{popblueL} \textbf{Appareil} & \textbf{Tension} & \textbf{Puissance} \\
\hline
Réfrigérateur & $\SI{220}{V}$ & $\SI{150}{W}$ \\
\hline
Fer à repasser & $\SI{220}{V}$ & $\SI{1000}{W}$ \\
\hline
Bouilloire électrique & $\SI{220}{V}$ & $\SI{2000}{W}$ \\
\hline
Machine à laver & $\SI{220}{V}$ & $\SI{2200}{W}$ \\
\hline
\end{tabularx}
\end{center}

\textbf{Document 3 : seuils des effets du courant alternatif traversant le corps humain}

\begin{center}
\begin{tabularx}{\linewidth}{|l|X|}
\hline
\rowcolor{popblueL} \textbf{Intensité $I$ traversant le corps} & \textbf{Effet sur l'organisme} \\
\hline
$I < \SI{0,5}{mA}$ & Aucune sensation \\
\hline
$\SI{0,5}{mA} \leq I < \SI{10}{mA}$ & Picotements, fourmillements \\
\hline
$\SI{10}{mA} \leq I < \SI{30}{mA}$ & Contractions musculaires, difficulté à lâcher l'objet \\
\hline
$\SI{30}{mA} \leq I < \SI{500}{mA}$ & Tétanisation, risque de paralysie respiratoire \\
\hline
$I \geq \SI{500}{mA}$ & Fibrillation cardiaque, risque de mort (électrocution) \\
\hline
\end{tabularx}
\end{center}

\textbf{Document 4 : mesures effectuées par le technicien}

\begin{itemize}
\item le courant de fuite de la machine à laver (dû à un défaut d'isolement interne) traverse le corps de M. Essomba avec une intensité mesurée $I_{\text{fuite}} = \SI{8}{mA}$ ;
\item la résistance du corps humain, en conditions humides, est estimée à $R = \SI{1000}{\Omega}$.
\end{itemize}

\begin{popfigure}
\begin{center}
\begin{tikzpicture}[scale=0.9, every node/.style={font=\small}]
% Machine à laver
\draw[thick, fill=popblueL] (0,0) rectangle (2.2,2.6);
\draw[thick, fill=white] (1.1,1.5) circle (0.7);
\node at (1.1,2.9) {\textbf{Machine à laver}};
% Fil de fuite vers carcasse
\draw[thick, poporange] (1.1,1.5) -- (2.2,1.9);
\fill[poporange] (2.2,1.9) circle (0.07);
\node[poporange, right] at (2.3,2.3) {défaut d'isolement :};
\node[poporange, right] at (2.3,1.9) {carcasse sous tension};
% Personne
\draw[thick] (4.6,1.9) circle (0.3);
\draw[thick] (4.6,1.6) -- (4.6,0.6);
\draw[thick] (4.6,1.3) -- (4.1,0.9);
\draw[thick] (4.6,1.3) -- (2.2,1.9);
\draw[thick] (4.6,0.6) -- (4.2,0);
\draw[thick] (4.6,0.6) -- (5.0,0);
\node[below] at (4.6,-0.1) {M. Essomba};
% Sol
\draw[thick] (-0.5,0) -- (6.5,0);
\foreach \x in {-0.3,0.3,0.9,1.5,2.1,2.7,3.3,3.9,4.5,5.1,5.7,6.3} \draw (\x,0) -- (\x-0.25,-0.25);
\node at (6.0,-0.55) {sol};
% Courant de fuite
\draw[-{Stealth}, very thick, poppink] (2.4,1.75) -- (4.0,1.35) node[midway, above, sloped] {$I_{\text{fuite}} = \SI{8}{mA}$};
\draw[-{Stealth}, very thick, poppink] (4.4,0.5) -- (4.1,0.05);
% Absence de terre
\node[popdark, align=center] at (1.1,-0.9) {pas de prise de terre :\\ le courant traverse le corps};
\end{tikzpicture}
\end{center}
\legende{Schéma de l'accident : à cause du défaut d'isolement, la carcasse métallique est sous tension. En l'absence de prise de terre, le courant de fuite traverse le corps de M. Essomba pour rejoindre le sol.}
\end{popfigure}

\courssub{Consignes}

Par groupes de quatre élèves, répondez aux questions suivantes, puis rédigez le rapport d'expert.

\begin{enumerate}
\item \textbf{Comprendre l'accident.} En utilisant les documents 1, 3 et 4, expliquer pourquoi M. Essomba a ressenti des picotements en touchant la carcasse de sa machine à laver. Nommer précisément le phénomène dont il a été victime (électrisation ou électrocution ?). Justifier.
\item \textbf{Vérifier par le calcul.} Sachant que le courant de fuite traverse le corps de M. Essomba dont la résistance est $R = \SI{1000}{\Omega}$, calculer la tension de contact à laquelle il a été soumis (loi d'Ohm : $U = R \times I$). Commenter le résultat.
\item \textbf{Expliquer le rôle manquant de la terre.} Expliquer, en une phrase, ce qui se serait passé si la maison avait possédé une prise de terre reliée à la carcasse de la machine.
\item \textbf{Étudier la surcharge de la ligne.} Calculer l'intensité du courant absorbé par chacun des trois appareils branchés sur les multiprises en cascade (réfrigérateur, fer, bouilloire), puis l'intensité totale $I_{\text{totale}}$ dans la ligne. Comparer au calibre du fusible ($\SI{32}{A}$) et à l'intensité maximale admissible d'une multiprise courante ($\SI{16}{A}$). Conclure sur le risque encouru.
\item \textbf{Recenser les défauts.} Dresser la liste complète des défauts de l'installation (au moins cinq).
\item \textbf{Rédiger le rapport d'expert.} Rédiger un texte de 10 à 15 lignes comportant : le diagnostic (causes de l'accident), les corrections techniques à apporter (avec leurs noms précis : prise de terre, interrupteur différentiel de $\SI{30}{mA}$, calibres de protection adaptés, répartition des appareils) et trois règles de comportement à rappeler à la famille Essomba.
\end{enumerate}

\courssub{Ressources à mobiliser}

\begin{aretenir}[Boîte à outils de la leçon]
\begin{itemize}
\item \textbf{Effets du courant sur le corps} : l'\cle{électrisation} est le passage du courant dans le corps avec des lésions (brûlures, contractions) ; l'\cle{électrocution} est une électrisation mortelle. La gravité dépend de l'intensité $I$, de la durée et du trajet du courant.
\item \textbf{Loi d'Ohm} : $U = R \times I$, avec $U$ en volts (V), $R$ en ohms ($\Omega$), $I$ en ampères (A).
\item \textbf{Puissance et intensité} : $P = U \times I$, donc $I = \dfrac{P}{U}$, avec $P$ en watts (W), $U$ en volts (V) et $I$ en ampères (A).
\item \textbf{Surcharge} : l'intensité totale demandée dépasse ce que la ligne ou la multiprise peut supporter ; les fils s'échauffent (effet Joule) : risque d'incendie.
\item \textbf{Court-circuit} : contact direct entre la phase et le neutre ; l'intensité devient très grande brutalement.
\item \textbf{Fusible et disjoncteur} : coupent le circuit quand l'intensité dépasse leur calibre ; le calibre doit être adapté à la ligne protégée.
\item \textbf{Prise de terre} : conduit les courants de fuite vers le sol au lieu de les laisser traverser le corps.
\item \textbf{Interrupteur différentiel $\SI{30}{mA}$} : compare le courant qui part et celui qui revient ; dès qu'une fuite supérieure à $\SI{30}{mA}$ est détectée, il coupe le circuit en une fraction de seconde.
\end{itemize}
\end{aretenir}

\courssub{Démarche et corrigé commenté}

\begin{exoresolu}[Rapport d'expertise de l'installation Essomba]
Reprendre les six consignes de l'activité et produire le diagnostic complet de l'installation de la famille Essomba.
\tcblower
\textbf{Question 1 : cause de l'accident.}

La machine à laver présente un \textbf{défaut d'isolement} : un fil sous tension touche sa carcasse métallique, qui se trouve alors portée à une tension dangereuse. Comme la maison n'a \textbf{pas de prise de terre}, ce courant de fuite ne peut pas s'écouler vers le sol. En touchant la carcasse, M. Essomba a fermé le circuit : le courant de fuite $I_{\text{fuite}} = \SI{8}{mA}$ a traversé son corps vers la terre.

D'après le document 3, pour $\SI{0,5}{mA} \leq I < \SI{10}{mA}$, l'effet est des picotements et fourmillements : cela correspond exactement aux symptômes décrits. M. Essomba a été victime d'une \textbf{électrisation} (et non d'une électrocution, car l'accident n'a pas été mortel). Il a eu de la chance : avec un courant de fuite plus important, il aurait pu être tétanisé, incapable de lâcher la carcasse.

\textbf{Question 2 : tension de contact.}

D'après la loi d'Ohm, en convertissant d'abord les milliampères en ampères ($\SI{8}{mA} = \num{0,008}{A}$) :
\[
U = R \times I = \SI{1000}{\Omega} \times \num{0,008}{A} = \SI{8}{V}.
\]
La tension de contact est de $\SI{8}{V}$. Elle paraît faible, mais c'est l'\textbf{intensité} qui est dangereuse, pas seulement la tension : avec un corps mouillé (résistance plus faible), la même fuite produirait une intensité bien plus grande. C'est pourquoi on ne touche jamais un appareil électrique avec les mains humides.

\textbf{Question 3 : rôle de la prise de terre.}

Si la carcasse avait été reliée à une prise de terre, le courant de fuite se serait écoulé directement dans le sol par ce chemin de faible résistance, au lieu de traverser le corps de M. Essomba ; de plus, un interrupteur différentiel aurait détecté la fuite et coupé le circuit instantanément.

\textbf{Question 4 : étude de la surcharge.}

On applique $I = \dfrac{P}{U}$ avec $U = \SI{220}{V}$ :
\begin{align*}
I_{\text{frigo}} &= \frac{150}{220} \approx \SI{0,68}{A} \\
I_{\text{fer}} &= \frac{1000}{220} \approx \SI{4,5}{A} \\
I_{\text{bouilloire}} &= \frac{2000}{220} \approx \SI{9,1}{A}
\end{align*}
Intensité totale traversant la première multiprise :
\[
I_{\text{totale}} = 0,68 + 4,5 + 9,1 \approx \SI{14,3}{A}.
\]
Deux conclusions :
\begin{itemize}
\item $I_{\text{totale}} \approx \SI{14,3}{A} < \SI{32}{A}$ : le fusible de $\SI{32}{A}$ ne fond pas, il est \textbf{mal calibré} (trop élevé) : il laisse passer sans réagir un courant dangereux pour les multiprises ;
\item $I_{\text{totale}} \approx \SI{14,3}{A}$ est proche de la limite de $\SI{16}{A}$ des multiprises, et le branchement en cascade fait passer tout ce courant dans la \textbf{première} multiprise et son cordon : ceux-ci s'échauffent par effet Joule. Si l'on ajoute un appareil (par exemple un grille-pain), la limite de $\SI{16}{A}$ est dépassée : risque d'\textbf{échauffement, de fonte des isolants et d'incendie}. Le fil dénudé de la bouilloire ajoute un risque de \textbf{court-circuit} et de contact direct avec la phase.
\end{itemize}

\textbf{Question 5 : liste des défauts.}

\begin{enumerate}
\item Absence de prise de terre.
\item Absence d'interrupteur différentiel.
\item Fusible de calibre trop élevé ($\SI{32}{A}$) pour la ligne et les multiprises.
\item Multiprises branchées en cascade.
\item Concentration de plusieurs appareils puissants sur une seule prise.
\item Prise murale à deux broches sans borne de terre pour un appareil à carcasse métallique.
\item Fil de la bouilloire dénudé.
\end{enumerate}

\textbf{Question 6 : rapport d'expert (modèle de rédaction).}

\emph{« L'accident de M. Essomba est une électrisation causée par un défaut d'isolement de la machine à laver, aggravé par l'absence totale de prise de terre : le courant de fuite de $\SI{8}{mA}$ a traversé son corps. Nous préconisons : (1) l'installation d'une prise de terre (piquet métallique et conducteur vert-jaune relié aux carcasses des appareils) ; (2) la pose d'un interrupteur différentiel de $\SI{30}{mA}$ en tête de l'installation ; (3) le remplacement du fusible de $\SI{32}{A}$ par un calibre adapté à la ligne ($\SI{16}{A}$) ; (4) la suppression des multiprises en cascade et la répartition des appareils puissants (bouilloire, fer, machine à laver) sur des prises distinctes ; (5) la réparation du cordon dénudé de la bouilloire. Enfin, nous rappelons à la famille trois règles : ne jamais toucher un appareil électrique avec les mains mouillées ; ne jamais brancher plusieurs appareils puissants sur une même multiprise ; en cas d'accident, couper d'abord le courant au disjoncteur général avant de porter secours à la victime. »}
\end{exoresolu}

\begin{attention}[Pièges fréquents dans ce type d'exercice]
\begin{itemize}
\item Ne confonds pas \textbf{électrisation} (accident non mortel) et \textbf{électrocution} (accident mortel).
\item Dans le calcul $I = \dfrac{P}{U}$, vérifie que $P$ est en watts et $U$ en volts ; n'oublie pas de convertir les milliampères : $\SI{8}{mA} = \num{0,008}{A}$.
\item Un fusible de calibre \textbf{trop élevé} est un défaut : il ne protège plus la ligne.
\item Le différentiel $\SI{30}{mA}$ protège les \textbf{personnes} ; le fusible protège les \textbf{installations}. Les deux sont complémentaires, l'un ne remplace pas l'autre.
\end{itemize}
\end{attention}

\begin{exercice}[Pour vérifier que tu as compris]
Sur la même ligne de la famille Essomba, on branche maintenant, en plus du réfrigérateur, du fer et de la bouilloire, un grille-pain de puissance $P = \SI{800}{W}$.
\begin{enumerate}
\item Calculer la nouvelle intensité totale dans la première multiprise.
\item La limite de $\SI{16}{A}$ de la multiprise est-elle dépassée ? Conclure.
\item Le fusible de $\SI{32}{A}$ réagit-il ? Que faut-il faire ?
\end{enumerate}
\end{exercice}

\begin{corrige}[Exercice : Pour vérifier que tu as compris]
\begin{enumerate}
\item Intensité du grille-pain : $I_{\text{gp}} = \dfrac{800}{220} \approx \SI{3,6}{A}$. Nouvelle intensité totale : $I'_{\text{totale}} = 14,3 + 3,6 \approx \SI{17,9}{A}$.
\item $I'_{\text{totale}} \approx \SI{17,9}{A} > \SI{16}{A}$ : la limite de la multiprise est dépassée. Les fils et la multiprise s'échauffent par effet Joule : risque de fonte des isolants et d'incendie.
\item $I'_{\text{totale}} \approx \SI{17,9}{A} < \SI{32}{A}$ : le fusible ne fond pas, il ne protège pas la ligne. Il faut le remplacer par un fusible (ou un disjoncteur divisionnaire) de calibre $\SI{16}{A}$ et répartir les appareils sur plusieurs prises.
\end{enumerate}
\end{corrige}

\courssub{Bilan de l'activité}

Cette expertise a permis de mettre en évidence que :

\begin{itemize}
\item un accident électrique domestique est presque toujours la conséquence de \textbf{plusieurs défauts cumulés} : ici, le défaut d'isolement de la machine n'est devenu dangereux que parce que la prise de terre et le différentiel manquaient ;
\item le \textbf{calcul} est un outil du technicien : la relation $I = \dfrac{P}{U}$ a permis de démontrer objectivement le risque de surcharge, sans avoir à provoquer l'incendie pour le constater ;
\item la sécurité électrique repose sur une \textbf{chaîne de protections complémentaires} : prise de terre (écoulement des fuites), interrupteur différentiel $\SI{30}{mA}$ (protection des personnes), fusibles et disjoncteurs bien calibrés (protection des installations contre surcharges et courts-circuits) ;
\item les \textbf{comportements} comptent autant que le matériel : mains sèches, pas de multiprises en cascade, coupure du courant avant toute intervention ou tout secours.
\end{itemize}

Tu es désormais capable, comme Aïcha, de repérer les dangers d'une installation électrique et de conseiller ta famille : c'est une compétence utile pour le BEPC... et pour toute la vie.

\newpage

\coursec{Méthodes et exercices résolus}

\courssub{Méthode 1 : Identifier les dangers du courant pour le corps humain}

\begin{methode}[Distinguer électrisation et électrocution]
Pour analyser une situation d'accident domestique :
\begin{enumerate}
\item \textbf{Repérer le passage du courant} : le corps humain est conducteur (il contient de l'eau salée). Un courant le traverse-t-il ?
\item \textbf{Identifier les points d'entrée et de sortie} : main--main, main--pieds, etc. Le trajet passe-t-il par le cœur ?
\item \textbf{Estimer l'intensité} $I$ qui traverse le corps : c'est elle qui détermine la gravité, pas la tension seule.
\item \textbf{Conclure avec le vocabulaire exact} :
\begin{itemize}
\item \cle{Électrisation} : le courant traverse le corps et provoque des lésions (brûlures, tétanisation, arrêt respiratoire), \emph{sans entraîner la mort}.
\item \cle{Électrocution} : l'accident électrique entraîne la \emph{mort}.
\end{itemize}
\end{enumerate}
\textbf{Réflexe} : retenir les seuils pour un courant alternatif de \SI{50}{Hz} : $I \approx \SI{1}{mA}$ (seuil de perception), $\SI{10}{mA}$ (tétanisation des muscles, on ne peut plus lâcher), $\SI{30}{mA}$ (risque de fibrillation cardiaque), au-delà de $\SI{1}{A}$ (brûlures graves, arrêt cardiaque).
\end{methode}

\begin{aretenir}[Les seuils de danger du courant alternatif (\SI{50}{Hz})]
\begin{center}
\begin{tabularx}{\linewidth}{|l|X|}\hline
\rowcolor{popblueL} \textbf{Intensité traversant le corps} & \textbf{Effet sur l'organisme} \\ \hline
$I \approx \SI{1}{mA}$ & Seuil de perception : picotements. \\ \hline
$I \approx \SI{10}{mA}$ & Tétanisation des muscles : on ne peut plus lâcher le conducteur. \\ \hline
$I \approx \SI{30}{mA}$ & Risque de fibrillation cardiaque (contractions désordonnées du cœur). \\ \hline
$I > \SI{1}{A}$ & Brûlures graves, arrêt cardiaque : danger de mort. \\ \hline
\end{tabularx}
\end{center}
C'est l'\cle{intensité} du courant qui traverse le corps (et sa durée) qui détermine la gravité de l'accident.
\end{aretenir}

\begin{exoresolu}[Accident domestique à Douala]
Un élève veut débrancher son fer à repasser. Les mains mouillées, il touche la fiche dénudée de la prise et reçoit une « secousse » violente qui le projette en arrière. Il s'en sort avec des brûlures légères aux doigts.
\begin{enumerate}
\item Définir les termes électrisation et électrocution.
\item Dire de quel accident il s'agit ici. Justifier.
\item Expliquer pourquoi les mains mouillées aggravent le danger.
\end{enumerate}
\tcblower
\textbf{1. Définitions.} L'\cle{électrisation} est l'ensemble des lésions provoquées par le passage d'un courant électrique à travers le corps humain, sans entraîner la mort : brûlures, tétanisation musculaire, troubles cardiaques. L'\cle{électrocution} est l'accident électrique qui entraîne la mort de la victime, le plus souvent par fibrillation cardiaque ou arrêt respiratoire.

\textbf{2. Nature de l'accident.} L'élève a été traversé par un courant (secousse, brûlures aux doigts) mais il a survécu : il s'agit donc d'une \textbf{électrisation}, et non d'une électrocution.

\textbf{3. Rôle de l'eau.} Le corps humain est conducteur grâce à l'eau salée qu'il contient. La peau sèche offre une résistance électrique relativement élevée ; la peau \emph{mouillée} a une résistance beaucoup plus faible. À tension égale ($U = \SI{220}{V}$ du réseau ENEO), la loi d'Ohm $I = \dfrac{U}{R}$ montre que si $R$ diminue, l'intensité $I$ qui traverse le corps \emph{augmente}. Le courant devient alors plus dangereux : il peut dépasser le seuil de tétanisation ($\SI{10}{mA}$) voire le seuil de fibrillation ($\SI{30}{mA}$).

\textbf{Conclusion :} l'élève a subi une électrisation, aggravée par l'humidité de ses mains qui a diminué la résistance de son corps et augmenté l'intensité du courant le traversant.
\end{exoresolu}

\courssub{Méthode 2 : Expliquer et prévenir la surcharge d'une installation}

\begin{methode}[Diagnostiquer une surcharge]
\begin{enumerate}
\item \textbf{Lister les appareils} branchés simultanément sur la même prise ou la même rallonge (multiprise).
\item \textbf{Additionner les puissances} : $P_{\text{totale}} = P_1 + P_2 + \cdots$
\item \textbf{Calculer l'intensité totale} appelée : $I = \dfrac{P_{\text{totale}}}{U}$ avec $U = \SI{220}{V}$.
\item \textbf{Comparer} $I$ à l'intensité maximale admissible $I_{\max}$ du circuit (calibre du fusible ou du disjoncteur).
\item \textbf{Conclure} : si $I > I_{\max}$, il y a \cle{surcharge} : les fils chauffent (effet Joule), l'isolant fond, risque d'incendie. Le dispositif de protection doit couper le circuit.
\end{enumerate}
\textbf{Réflexe} : la surcharge vient de \emph{trop d'appareils branchés en même temps}, jamais d'un appareil seul en bon état.
\end{methode}

\begin{exoresolu}[La rallonge de la boutique]
Dans une boutique de quartier, on branche sur une même multiprise un réfrigérateur de puissance $P_1 = \SI{150}{W}$, un fer à repasser $P_2 = \SI{1000}{W}$ et une bouilloire $P_3 = \SI{2000}{W}$. La tension du réseau est $U = \SI{220}{V}$. Le fusible protégeant ce circuit est calibré à $I_{\max} = \SI{10}{A}$.
\begin{enumerate}
\item Calculer l'intensité totale traversant la multiprise.
\item Y a-t-il surcharge ? Que risque-t-il de se passer ?
\item Proposer une solution.
\end{enumerate}
\tcblower
\textbf{1. Calcul de l'intensité totale.} La puissance totale est :
\[ P_{\text{totale}} = P_1 + P_2 + P_3 = 150 + 1000 + 2000 = \SI{3150}{W}. \]
La relation $P = U \cdot I$ donne :
\[ I = \frac{P_{\text{totale}}}{U} = \frac{3150}{220} \approx \SI{14,3}{A}. \]

\textbf{2. Comparaison.} On a $I \approx \SI{14,3}{A} > I_{\max} = \SI{10}{A}$ : il y a \textbf{surcharge}. Le fusible de \SI{10}{A} va fondre et couper le circuit. Si la protection est défaillante (fusible remplacé par un fil de fer, pratique dangereuse mais fréquente), les fils de la rallonge vont chauffer par effet Joule, l'isolant plastique peut fondre et provoquer un \textbf{incendie}.

\textbf{3. Solution.} Ne pas brancher simultanément les trois appareils : par exemple utiliser la bouilloire seule ($I = \dfrac{2000}{220} \approx \SI{9,1}{A} < \SI{10}{A}$, c'est acceptable), ou brancher le fer sur une autre prise alimentée par un autre circuit.

\textbf{Conclusion :} avec $I \approx \SI{14,3}{A} > \SI{10}{A}$, la multiprise est en surcharge ; il faut répartir les appareils sur des circuits différents ou limiter leur usage simultané.
\end{exoresolu}

\courssub{Méthode 3 : Expliquer le court-circuit et le rôle du fusible}

\begin{methode}[Raisonner sur un court-circuit]
\begin{enumerate}
\item \textbf{Définir} : il y a \cle{court-circuit} lorsque deux points d'un circuit (phase et neutre) sont reliés directement par un conducteur de résistance quasi nulle, sans passer par un appareil.
\item \textbf{Appliquer la loi d'Ohm} : $I = \dfrac{U}{R}$. Si $R \to 0$, alors $I$ devient \emph{très grande} (plusieurs dizaines d'ampères).
\item \textbf{Conséquence} : échauffement brutal des fils (effet Joule), étincelles, incendie.
\item \textbf{Rôle du fusible} : le \cle{fusible} contient un fil métallique fin qui \emph{fond} dès que $I$ dépasse son calibre ; il coupe le circuit et protège l'installation. Il est à usage unique : après fusion, on le \emph{remplace} par un fusible de même calibre.
\end{enumerate}
\textbf{Réflexe} : causes typiques de court-circuit : fils dénudés qui se touchent, isolant usé, eau renversée sur une prise, fiche mal câblée.
\end{methode}

\begin{popfigure}
\begin{center}
\begin{tikzpicture}[scale=1, every node/.style={font=\small}]
% Fils phase et neutre
\draw[very thick, popink] (0,1.2) -- (8,1.2) node[right, popink]{Phase};
\draw[very thick, popblue] (0,-1.2) -- (8,-1.2) node[right, popblue]{Neutre};
% Fusible sur la phase
\draw[fill=popgoldL, draw=popdark] (1.2,0.95) rectangle (2.4,1.45);
\node[popdark] at (1.8,1.75) {Fusible};
% Appareil (résistor)
\draw[fill=popblueL, draw=popdark] (5.4,-1.2) rectangle (6.6,1.2);
\node[popdark] at (6,0) {Appareil};
% Court-circuit : fil dénudé reliant phase et neutre
\draw[very thick, poporange, decorate, decoration={zigzag, segment length=4, amplitude=1.5}] (3.6,1.2) -- (3.6,-1.2);
\node[poporange, align=center] at (3.6,-1.9) {Court-circuit\\ ($R \approx 0$)};
% Flèche courant
\draw[-{Stealth}, very thick, popgreen] (2.6,1.2) -- (3.4,1.2);
\node[popgreen, above] at (3.0,1.25) {$I$ très grande};
\end{tikzpicture}
\end{center}
\legende{Schéma d'un court-circuit : un contact direct entre la phase et le neutre court-circuite l'appareil ; l'intensité devient très grande et le fusible fond pour couper le circuit.}
\end{popfigure}

\begin{exoresolu}[Le grille-pain défectueux]
Dans une cuisine, les fils dénudés du cordon d'un grille-pain se touchent accidentellement. Une étincelle jaillit et le fusible de \SI{16}{A} « saute » instantanément.
\begin{enumerate}
\item Expliquer pourquoi l'intensité devient très grande au moment du contact.
\item Expliquer le rôle joué par le fusible.
\item Pourquoi ne faut-il jamais remplacer le fusible fondu par un clou ou un fil de fer ?
\end{enumerate}
\tcblower
\textbf{1. Intensité du court-circuit.} Quand les fils de phase et de neutre se touchent directement, le courant ne traverse plus le résistor du grille-pain : la résistance du circuit devient quasi nulle ($R \approx 0$). D'après la loi d'Ohm $I = \dfrac{U}{R}$, avec $U = \SI{220}{V}$ et $R$ très faible, l'intensité $I$ devient \textbf{très grande} (elle peut dépasser \SI{100}{A}) : c'est le court-circuit.

\textbf{2. Rôle du fusible.} Le fusible de \SI{16}{A} contient un fil fin calibré. Dès que $I$ dépasse \SI{16}{A}, ce fil fond par effet Joule : le circuit est \emph{ouvert}, le courant cesse de circuler. Le fusible a donc \textbf{protégé l'installation} contre l'échauffement des fils et le risque d'incendie.

\textbf{3. Danger du bricolage.} Un clou ou un fil de fer ne fond pas à \SI{16}{A} : il laisserait passer un courant très intense sans couper le circuit. Les fils de l'installation chaufferaient alors jusqu'à faire fondre leur isolant et déclencher un incendie. Il faut toujours remplacer un fusible par un fusible \emph{neuf de même calibre}, après avoir réparé la cause du court-circuit.

\textbf{Conclusion :} le court-circuit produit une intensité énorme car $R \approx 0$ ; le fusible, en fondant, coupe le circuit et protège la maison. Le remplacer par un objet métallique quelconque supprime cette protection.
\end{exoresolu}

\courssub{Méthode 4 : Choisir le calibre d'un disjoncteur ou d'un fusible}

\begin{methode}[Calculer le calibre de protection]
\begin{enumerate}
\item \textbf{Calculer l'intensité nominale} du circuit : $I = \dfrac{P}{U}$ (ou somme des intensités des appareils prévus).
\item \textbf{Choisir le calibre normalisé juste supérieur} : \SI{10}{A}, \SI{16}{A}, \SI{20}{A}, \SI{32}{A}...
\item \textbf{Vérifier} que le calibre choisi est inférieur à l'intensité maximale que les fils peuvent supporter.
\item \textbf{Rédiger la conclusion} : « on choisit un disjoncteur (ou fusible) de calibre \SI{16}{A} ».
\end{enumerate}
\textbf{Réflexe} : un calibre \emph{trop grand} ne protège pas ; un calibre \emph{trop petit} fait sauter l'installation en permanence.
\end{methode}

\begin{exoresolu}[Protéger le circuit d'un chauffe-eau]
Un chauffe-eau porte les indications : $P = \SI{2200}{W}$, $U = \SI{220}{V}$. On veut le protéger par un disjoncteur. Calibres disponibles : \SI{6}{A}, \SI{10}{A}, \SI{16}{A}, \SI{20}{A}.
\begin{enumerate}
\item Calculer l'intensité absorbée par le chauffe-eau en fonctionnement normal.
\item Choisir le calibre du disjoncteur en justifiant.
\end{enumerate}
\tcblower
\textbf{1. Intensité absorbée.} De la relation $P = U \cdot I$, on tire :
\[ I = \frac{P}{U} = \frac{2200}{220} = \SI{10}{A}. \]

\textbf{2. Choix du calibre.} Le disjoncteur doit laisser passer le courant normal de \SI{10}{A} sans déclencher, mais couper en cas de dépassement anormal.
\begin{itemize}
\item Un calibre de \SI{6}{A} est \emph{trop faible} : il déclencherait dès la mise en route ($\SI{10}{A} > \SI{6}{A}$).
\item Un calibre de \SI{10}{A} est juste à la limite : la moindre fluctuation du réseau le ferait déclencher.
\item Un calibre de \SI{16}{A} laisse passer le fonctionnement normal et protège efficacement contre les surintensités : c'est le bon choix.
\item Un calibre de \SI{20}{A} protégerait moins finement.
\end{itemize}
\textbf{Conclusion :} on choisit un disjoncteur de calibre \SI{16}{A}, immédiatement supérieur à l'intensité normale de \SI{10}{A}.
\end{exoresolu}

\courssub{Méthode 5 : Expliquer le rôle de la prise de terre et de l'interrupteur différentiel}

\begin{methode}[Raisonner sur un défaut d'isolement]
\begin{enumerate}
\item \textbf{Identifier le défaut} : un fil de phase touche la carcasse métallique d'un appareil (machine à laver, réfrigérateur). La carcasse se retrouve sous tension.
\item \textbf{Sans prise de terre} : la personne qui touche la carcasse ferme le circuit vers le sol à travers son corps $\Rightarrow$ électrisation, voire électrocution.
\item \textbf{Avec prise de terre} : le fil de terre relie la carcasse au sol ; le courant de défaut s'écoule vers la terre par ce fil (résistance très faible) au lieu de traverser le corps.
\item \textbf{Rôle du différentiel} : l'\cle{interrupteur différentiel} compare le courant entrant (phase) et le courant sortant (neutre). Si la différence dépasse son seuil (par exemple \SI{30}{mA}), c'est qu'une partie du courant « fuit » vers la terre : il coupe le circuit en une fraction de seconde, \emph{protégeant les personnes}.
\end{enumerate}
\textbf{Réflexe} : le fusible/disjoncteur protège les \emph{installations} ; le différentiel protège les \emph{personnes} ; la prise de terre rend le différentiel efficace. Les trois sont complémentaires.
\end{methode}

\begin{popfigure}
\begin{center}
\begin{tikzpicture}[scale=1, every node/.style={font=\small}]
% Machine à laver
\draw[fill=popblueL, draw=popdark, thick] (4.5,-0.5) rectangle (6.5,1.5);
\draw[fill=white, draw=popdark] (5.5,0.5) circle (0.55);
\node[popdark] at (5.5,1.85) {Carcasse métallique};
% Fils phase et neutre
\draw[very thick, popink] (0,1.0) -- (4.5,1.0) node[midway, above, popink]{Phase};
\draw[very thick, popblue] (0,0.0) -- (4.5,0.0) node[midway, above, popblue]{Neutre};
% Différentiel
\draw[fill=popgoldL, draw=popdark] (1.2,-0.35) rectangle (2.2,1.35);
\node[popdark, align=center] at (1.7,2.0) {Interrupteur\\ différentiel};
% Fil de terre
\draw[very thick, popgreen] (5.5,-0.5) -- (5.5,-1.8);
\draw[popgreen, thick] (4.9,-1.8) -- (6.1,-1.8);
\draw[popgreen, thick] (5.1,-2.0) -- (5.9,-2.0);
\draw[popgreen, thick] (5.3,-2.2) -- (5.7,-2.2);
\node[popgreen, right] at (6.1,-1.9) {Prise de terre};
% Courant de défaut
\draw[-{Stealth}, very thick, poporange] (5.5,-0.6) -- (5.5,-1.6);
\node[poporange, left, align=right] at (5.3,-1.1) {Courant\\ de défaut};
% Personne
\draw[fill=poppinkL, draw=popdark] (7.6,1.0) circle (0.3);
\draw[popdark, thick] (7.6,0.7) -- (7.6,-0.6);
\draw[popdark, thick] (7.6,0.3) -- (6.5,0.6);
\draw[popdark, thick] (7.6,-0.6) -- (7.3,-1.6);
\draw[popdark, thick] (7.6,-0.6) -- (7.9,-1.6);
\node[popdark, right] at (8.0,0.9) {Personne protégée};
\end{tikzpicture}
\end{center}
\legende{En cas de défaut d'isolement, le fil de terre évacue le courant de défaut vers le sol et l'interrupteur différentiel coupe le circuit : la personne qui touche la carcasse est protégée.}
\end{popfigure}

\begin{exoresolu}[La machine à laver qui « pique »]
Dans une maison de Yaoundé, Madame Ngo sent des picotements en touchant la carcasse métallique de sa machine à laver. L'installation est ancienne : la prise de terre n'est pas raccordée et il n'y a pas d'interrupteur différentiel.
\begin{enumerate}
\item Expliquer l'origine des picotements.
\item Expliquer comment une prise de terre correctement raccordée aurait évité ce danger.
\item Expliquer le rôle complémentaire d'un interrupteur différentiel de \SI{30}{mA}.
\end{enumerate}
\tcblower
\textbf{1. Origine des picotements.} Un fil de phase, dont l'isolant est usé, touche la carcasse métallique de la machine : celle-ci est portée à la tension du réseau (\SI{220}{V}). En la touchant, Madame Ngo relie la carcasse au sol à travers son corps : un courant de défaut la traverse, d'où les picotements. C'est le début d'une électrisation ; si ses mains étaient mouillées, l'accident pourrait devenir mortel.

\textbf{2. Rôle de la prise de terre.} Le fil de terre (vert-jaune) relie la carcasse métallique au sol par un conducteur de très faible résistance. En cas de défaut, le courant s'écoule préférentiellement vers la terre par ce fil plutôt qu'à travers le corps humain, dont la résistance est beaucoup plus grande. La carcasse reste à un potentiel proche de celui du sol : la personne ne reçoit plus de choc électrique.

\textbf{3. Rôle de l'interrupteur différentiel.} En fonctionnement normal, le courant qui arrive par la phase est égal à celui qui repart par le neutre. En cas de fuite vers la terre, ces deux courants diffèrent. L'interrupteur différentiel de \SI{30}{mA} détecte cette différence dès qu'elle dépasse \SI{30}{mA} — seuil en dessous de la fibrillation cardiaque — et coupe automatiquement le circuit en moins d'une fraction de seconde, avant que le courant ne devienne mortel.

\textbf{Conclusion :} la prise de terre évacue le courant de défaut et le différentiel coupe le circuit dès qu'une fuite de \SI{30}{mA} est détectée : ensemble, ils protègent efficacement les personnes. Madame Ngo doit faire raccorder la terre et installer un différentiel par un électricien qualifié.
\end{exoresolu}

\courssub{Méthode 6 : Appliquer les règles de sécurité et la conduite à tenir}

\begin{methode}[Réagir face à un accident électrique]
\textbf{Prévention (avant l'accident)} :
\begin{enumerate}
\item Ne jamais toucher un appareil électrique avec les mains mouillées ou les pieds nus.
\item Ne jamais introduire d'objet métallique dans une prise.
\item Ne pas surcharger les multiprises ; débrancher en tirant sur la fiche, pas sur le fil.
\item Faire vérifier l'installation (terre, différentiel, fusibles adaptés).
\end{enumerate}
\textbf{Secours (après l'accident), dans l'ordre} :
\begin{enumerate}
\item \textbf{Couper le courant} immédiatement (disjoncteur général) : ne \emph{jamais} toucher la victime tant qu'elle est en contact avec le circuit, sous peine d'être électrisé à son tour.
\item \textbf{Alerter les secours} (pompiers : 118 au Cameroun ; SAMU : 119).
\item \textbf{Secourir} : si la victime ne respire plus, commencer les gestes de premiers secours en attendant les secours.
\end{enumerate}
\textbf{Réflexe} : la première action est toujours « couper le courant », jamais « tirer la victime ».
\end{methode}

\begin{exoresolu}[Accident près du groupe électrogène]
Pendant une coupure ENEO, un père de famille bricole le raccordement de son groupe électrogène. Son fils de 15 ans le trouve inconscient, la main encore collée au fil dénudé.
\begin{enumerate}
\item Quelle doit être la toute première action du fils ? Justifier.
\item Pourquoi ne doit-il pas tirer son père par le bras immédiatement ?
\item Citer deux règles de prévention qui auraient pu éviter cet accident.
\end{enumerate}
\tcblower
\textbf{1. Première action.} Le fils doit d'abord \textbf{couper le courant} : arrêter le groupe électrogène ou actionner l'interrupteur général. Tant que le courant n'est pas coupé, la victime reste traversée par le courant (tétanisation : elle ne peut pas lâcher le fil) et le danger persiste.

\textbf{2. Pourquoi ne pas toucher la victime.} Le corps humain est conducteur. Si le fils touche son père alors que celui-ci est encore en contact avec le fil sous tension, le courant traversera \emph{aussi} le corps du fils : il y aurait deux victimes au lieu d'une. On ne dégage la victime qu'après coupure du courant, éventuellement en la poussant avec un objet isolant sec (manche en bois) si la coupure est impossible.

\textbf{3. Prévention.} Par exemple :
\begin{itemize}
\item ne jamais intervenir sur une installation sous tension : couper d'abord l'alimentation ;
\item ne pas utiliser de câbles dénudés ou mal isolés, et faire les raccordements avec du matériel conforme (dominos, prises, disjoncteur adapté).
\end{itemize}

\textbf{Conclusion :} face à un accident électrique, la conduite à tenir est : couper le courant, alerter les secours (118/119), puis secourir. Toucher la victime avant la coupure expose le sauveteur au même danger.
\end{exoresolu}

\begin{savaistu}[Le paratonnerre, ancêtre de la prise de terre]
Inventé par Benjamin Franklin en 1752, le \cle{paratonnerre} protège les bâtiments de la foudre : une pointe métallique placée sur le toit capte la décharge électrique et un gros câble la conduit directement \emph{vers la terre}, évitant qu'elle ne traverse la maison. C'est exactement le même principe que la prise de terre de nos installations : offrir au courant un chemin facile vers le sol. Au Cameroun, où les orages sont violents en saison des pluies, de nombreux édifices publics en sont équipés.
\end{savaistu}

\begin{attention}[Les 5 erreurs qui coûtent des points au BEPC]
\begin{enumerate}
\item \textbf{Confondre électrisation et électrocution.} L'électrocution est réservée au cas où l'accident entraîne la \emph{mort}. Une victime vivante est électrisée. Ce vocabulaire est systématiquement exigé dans la justification.
\item \textbf{Dire « la tension tue » sans nuancer.} C'est l'\emph{intensité} traversant le corps qui détermine la gravité (seuils en mA). La tension élevée est dangereuse parce qu'elle produit une intensité élevée à travers la résistance du corps ($I = \dfrac{U}{R}$).
\item \textbf{Confondre surcharge et court-circuit.} Surcharge : trop d'appareils branchés, $I$ dépasse modérément le calibre. Court-circuit : contact direct phase--neutre, $R \approx 0$, $I$ devient énorme. Les causes et les schémas sont différents.
\item \textbf{Oublier l'unité ou le calcul dans le choix d'un calibre.} Il faut écrire $I = \dfrac{P}{U}$, faire l'application numérique avec les unités, puis choisir le calibre \emph{normalisé juste supérieur} et conclure par une phrase. Un calibre choisi sans calcul justificatif ne rapporte pas tous les points.
\item \textbf{Attribuer au fusible la protection des personnes.} Le fusible et le disjoncteur protègent les \emph{installations} (contre surcharge et court-circuit). Seuls la prise de terre et l'interrupteur différentiel (\SI{30}{mA}) protègent les \emph{personnes}. De même, en secours : écrire « on touche la victime pour la dégager » avant « on coupe le courant » est une faute grave.
\end{enumerate}
\end{attention}

\newpage

\coursec{Exercices}

\courssub{Je vérifie mes connaissances}

\begin{exercice}[QCM : les bons réflexes]
Pour chaque question, choisis la (ou les) bonne(s) réponse(s).
\begin{enumerate}
\item Le passage du courant électrique dans le corps humain s'appelle :
\begin{enumerate}
\item l'électrocution ;
\item l'électrisation ;
\item l'électrolyse.
\end{enumerate}
\item Le rôle d'un fusible est de :
\begin{enumerate}
\item protéger les personnes contre les fuites de courant ;
\item fondre et couper le circuit quand l'intensité devient trop grande ;
\item augmenter la tension du réseau.
\end{enumerate}
\item Un interrupteur différentiel de \SI{30}{mA} protège :
\begin{enumerate}
\item les appareils contre les surtensions ;
\item les personnes contre les contacts directs et indirects ;
\item l'installation contre la foudre uniquement.
\end{enumerate}
\item Face à une victime encore en contact avec un fil électrique, il faut d'abord :
\begin{enumerate}
\item la tirer par le bras ;
\item couper le courant (disjoncteur général) ;
\item lui jeter de l'eau pour la réveiller.
\end{enumerate}
\end{enumerate}
\end{exercice}

\begin{exercice}[Vrai ou faux, en justifiant]
Réponds par vrai ou faux, puis justifie ta réponse en une phrase.
\begin{enumerate}
\item Un courant de faible intensité ne peut jamais être dangereux pour l'homme.
\item Le danger du courant augmente si le contact dure longtemps.
\item On peut remplacer un fusible fondu par un fil de fer, puisque le fil conduit aussi le courant.
\item La prise de terre permet d'évacuer vers le sol les courants de fuite d'un appareil.
\item Brancher plusieurs multiprises en cascade sur une seule prise est sans danger si les appareils sont éteints.
\end{enumerate}
\end{exercice}

\begin{exercice}[Définitions et vocabulaire]
\begin{enumerate}
\item Définis les termes suivants : \cle{électrisation}, \cle{électrocution}, \cle{court-circuit}, \cle{surcharge}.
\item Recopie et complète : le \ldots\ldots\ldots protège les installations contre les surintensités ; l'interrupteur \ldots\ldots\ldots protège les personnes en comparant le courant qui part et le courant qui revient ; la prise de \ldots\ldots\ldots évacue les courants de défaut vers le sol.
\end{enumerate}
\end{exercice}

\begin{exercice}[Calculs directs]
Un grille-pain de puissance $P = \SI{1100}{W}$ fonctionne sous la tension du réseau $U = \SI{220}{V}$.
\begin{enumerate}
\item Calcule l'intensité $I$ du courant qui le traverse en régime normal.
\item Parmi les calibres de fusible suivants : \SI{2}{A}, \SI{6}{A}, \SI{10}{A}, \SI{16}{A}, lequel choisis-tu pour protéger ce circuit ? Justifie.
\item Quelle énergie électrique, en wattheures, ce grille-pain consomme-t-il s'il fonctionne pendant \SI{6}{min} ?
\end{enumerate}
\end{exercice}

\courssub{Je m'entraîne}

\begin{exercice}[Le chauffe-eau de la maison]
Dans une maison de Yaoundé, un chauffe-eau porte les indications : $P = \SI{2200}{W}$ ; $U = \SI{220}{V}$.
\begin{enumerate}
\item Calcule l'intensité du courant absorbé par le chauffe-eau en fonctionnement normal.
\item On dispose de fusibles de calibres \SI{6}{A}, \SI{10}{A}, \SI{16}{A} et \SI{32}{A}. Choisis le calibre adapté en justifiant ton choix.
\item Explique ce qui se passerait si l'on installait un fusible de calibre \SI{6}{A}, puis un fusible de calibre \SI{32}{A}.
\end{enumerate}
\end{exercice}

\begin{exercice}[La ligne de la cuisine]
Une prise de la cuisine, protégée par un disjoncteur de calibre \SI{16}{A}, alimente sous \SI{220}{V} une bouilloire de \SI{2000}{W}, un fer à repasser de \SI{1000}{W} et un téléviseur de \SI{150}{W}, branchés en même temps sur une multiprise.
\begin{enumerate}
\item Calcule l'intensité traversant chaque appareil.
\item Calcule l'intensité totale fournie par la ligne.
\item Le disjoncteur déclenche-t-il ? Justifie par le calcul.
\item Propose deux solutions concrètes pour éviter ce problème.
\end{enumerate}
\end{exercice}

\begin{exercice}[Lecture d'une courbe de danger]
Le tableau ci-dessous donne les effets d'un courant alternatif traversant le corps humain, en fonction de l'intensité $I$ et de la durée $t$ de contact.
\begin{center}
\begin{tabularx}{\linewidth}{|l|X|X|}\hline
\rowcolor{popblueL}
\textbf{Zone} & \textbf{Conditions} & \textbf{Effet sur l'organisme} \\ \hline
1 & $I < \SI{0,5}{mA}$ & Aucune sensation \\ \hline
2 & $\SI{0,5}{mA} \le I < \SI{10}{mA}$ & Picotements, contractures musculaires \\ \hline
3 & $\SI{10}{mA} \le I < \SI{30}{mA}$, $t$ courte & Tétanisation, difficultés respiratoires \\ \hline
4 & $I \ge \SI{30}{mA}$ pendant plus de \SI{0,1}{s} & Fibrillation cardiaque possible, danger de mort \\ \hline
\end{tabularx}
\end{center}
\begin{enumerate}
\item Un courant de \SI{50}{mA} traverse le corps d'une personne pendant \SI{0,5}{s}. Cette personne est-elle en danger ? Justifie à l'aide du tableau.
\item Explique pourquoi un interrupteur différentiel de sensibilité \SI{30}{mA} protège efficacement les personnes.
\item Un différentiel de \SI{300}{mA} protège-t-il aussi bien les personnes ? Justifie.
\end{enumerate}
\end{exercice}

\begin{exercice}[Chasse aux erreurs dans une installation]
Sur le plan de l'installation électrique d'une chambre et d'une salle de bain, on relève les situations suivantes :
\begin{itemize}
\item le fusible de la ligne a été remplacé par un bout de fil de cuivre ;
\item la prise de la salle de bain n'est pas reliée à la terre ;
\item trois multiprises sont branchées les unes sur les autres ;
\item un radiateur électrique est posé juste à côté de la baignoire.
\end{itemize}
\begin{enumerate}
\item Pour chacune des quatre situations, explique le danger encouru.
\item Propose la correction à apporter dans chaque cas.
\item Rédige une consigne de sécurité générale que tu afficherais dans cette maison.
\end{enumerate}
\end{exercice}

\courssub{Je me prépare au Bac}

\begin{exercice}[Sujet type BEPC : l'atelier de menuiserie]
Dans un atelier de menuiserie à Douala, l'installation électrique est alimentée sous $U = \SI{220}{V}$. On y trouve une scie circulaire de puissance $P_1 = \SI{1500}{W}$, une perceuse de puissance $P_2 = \SI{800}{W}$ et quatre lampes de \SI{100}{W} chacune. La ligne est protégée par un disjoncteur de calibre \SI{20}{A}.
\begin{enumerate}
\item Définis les mots \cle{électrisation} et \cle{électrocution}.
\item Calcule l'intensité totale du courant lorsque tous les appareils fonctionnent en même temps.
\item Le disjoncteur de \SI{20}{A} est-il correctement choisi ? Justifie ta réponse par le calcul.
\item Un jour, l'ouvrier perce par mégarde le câble d'alimentation de la scie : les deux fils se touchent.
\begin{enumerate}
\item Comment appelle-t-on ce défaut ?
\item Décris ce qui se passe alors dans le circuit et explique le rôle du disjoncteur.
\end{enumerate}
\item La carcasse métallique de la scie est reliée à la prise de terre. Explique en quelques phrases comment cette liaison, associée à un interrupteur différentiel de \SI{30}{mA}, protège l'ouvrier en cas de défaut d'isolement.
\end{enumerate}
\end{exercice}

\begin{exercice}[Sujet type BEPC : accident domestique et secours]
Un soir à Bafoussam, pendant une coupure ENEO, un groupe électrogène est branché sur l'installation de la maison. Un enfant touche un fil dénudé d'une rallonge humide et reste « collé » au fil. L'intensité du courant qui traverse son corps est estimée à \SI{80}{mA} et le contact dure \SI{2}{s} avant que le père n'intervienne.
\begin{enumerate}
\item Cite trois facteurs qui aggravent le danger du courant électrique pour le corps humain, et indique lesquels sont présents dans cette situation.
\item À l'aide des seuils suivants : tétanisation à partir de \SI{10}{mA} ; fibrillation cardiaque possible au-delà de \SI{30}{mA} pendant plus de \SI{0,1}{s} ; explique pourquoi la vie de l'enfant est en danger.
\item Décris, dans l'ordre, les trois premières actions que le père doit effectuer pour secourir l'enfant sans se mettre lui-même en danger.
\item Cite deux dispositifs qui, installés sur le tableau électrique, auraient pu couper automatiquement le courant avant l'accident, et précise le rôle de chacun.
\item Rédige un court message de sensibilisation (4 à 5 lignes) sur les dangers des rallonges dénudées et de l'humidité, à afficher dans ton quartier.
\end{enumerate}
\end{exercice}

\begin{exercice}[Problème de synthèse : sécuriser une maison]
Mme Etoundi fait construire une maison à Bafoussam. L'installation prévue, alimentée sous \SI{220}{V}, comprend : un chauffe-eau de \SI{2200}{W}, une cuisinière électrique de \SI{3000}{W}, un réfrigérateur de \SI{200}{W}, dix lampes de \SI{60}{W} et divers petits appareils. L'électricien propose : un disjoncteur général de \SI{32}{A}, un interrupteur différentiel de \SI{30}{mA}, une prise de terre et des fusibles calibrés par circuit.
\begin{enumerate}
\item Calcule l'intensité totale si le chauffe-eau, la cuisinière, le réfrigérateur et toutes les lampes fonctionnent simultanément. Le disjoncteur général de \SI{32}{A} convient-il ? Justifie.
\item Choisis le calibre du fusible du circuit du chauffe-eau parmi \SI{10}{A}, \SI{16}{A} et \SI{20}{A}, en justifiant.
\item Explique la différence de fonction entre le disjoncteur (ou le fusible) et l'interrupteur différentiel : lequel protège les biens ? lequel protège les personnes ?
\item Le plomb de la prise de terre mesure une résistance de prise de terre $R = \SI{50}{\Omega}$. En cas de défaut, un courant de fuite $I_d = \SI{0,5}{A}$ s'écoule vers la terre. Calcule la tension de contact $U_c = R \cdot I_d$ à laquelle serait soumise une personne touchant la carcasse de l'appareil, et conclus sur le danger sachant qu'une tension supérieure à \SI{50}{V} est dangereuse en milieu sec.
\item Propose deux règles de sécurité supplémentaires à respecter dans la salle de bain de cette maison.
\end{enumerate}
\end{exercice}



\newpage

\coursec{Corrigés des exercices}

\begin{corrige}[Exercice 1 — QCM : les bons réflexes]
\begin{enumerate}
\item Réponse \textbf{b} : l'\cle{électrisation} est le passage du courant dans le corps ; l'électrocution est l'électrisation mortelle ; l'électrolyse est une transformation chimique.
\item Réponse \textbf{b} : le fusible fond par effet Joule quand l'intensité dépasse son calibre, ce qui ouvre le circuit.
\item Réponse \textbf{b} : le différentiel de \SI{30}{mA} détecte les fuites de courant et protège les personnes.
\item Réponse \textbf{b} : il faut d'abord \textbf{couper le courant} avant tout contact avec la victime, sinon le secouriste s'électrise à son tour.
\end{enumerate}
\end{corrige}

\begin{corrige}[Exercice 2 — Vrai ou faux, en justifiant]
\begin{enumerate}
\item \textbf{Faux.} Dès \SI{10}{mA}, le courant provoque une tétanisation des muscles ; le danger dépend de l'intensité \emph{et} de la durée.
\item \textbf{Vrai.} Plus la durée de contact est longue, plus les lésions (brûlures, fibrillation cardiaque) sont graves.
\item \textbf{Faux.} Le fil de fer ne fond pas au bon calibre : la protection n'existe plus et l'installation peut s'échauffer jusqu'à l'incendie.
\item \textbf{Vrai.} Le fil de terre offre au courant de défaut un chemin vers le sol, ce qui fait déclencher le différentiel.
\item \textbf{Faux.} Même éteints, les appareils branchés en cascade créent un risque de surcharge dès leur mise en marche et de mauvais contacts échauffant la prise.
\end{enumerate}
\end{corrige}

\begin{corrige}[Exercice 3 — Définitions et vocabulaire]
\begin{enumerate}
\item \cle{Électrisation} : passage du courant électrique à travers le corps humain. \cle{Électrocution} : électrisation qui entraîne la mort. \cle{Court-circuit} : contact direct entre les deux fils d'alimentation (phase et neutre) sans récepteur, provoquant une intensité très grande. \cle{Surcharge} : fonctionnement simultané de trop d'appareils sur une même ligne, l'intensité totale dépassant celle prévue.
\item Le \textbf{fusible} (ou le disjoncteur) protège les installations contre les surintensités ; l'interrupteur \textbf{différentiel} protège les personnes en comparant le courant qui part et le courant qui revient ; la prise de \textbf{terre} évacue les courants de défaut vers le sol.
\end{enumerate}
\end{corrige}

\begin{corrige}[Exercice 4 — Calculs directs]
\begin{enumerate}
\item On utilise $P = U \cdot I$, donc :
\[ I = \frac{P}{U} = \frac{1100}{220} = \SI{5}{A}. \]
\item On choisit le calibre \textbf{immédiatement supérieur} à l'intensité normale : \SI{6}{A}. Le calibre \SI{2}{A} fondrait en fonctionnement normal ; \SI{10}{A} et \SI{16}{A} protégeraient moins finement.
\item Durée : $t = \SI{6}{min} = \SI{0,1}{h}$. Énergie :
\[ E = P \cdot t = 1100 \times 0,1 = \SI{110}{Wh}. \]
\end{enumerate}
\end{corrige}

\begin{corrige}[Exercice 5 — Le chauffe-eau de la maison]
\begin{enumerate}
\item $I = \dfrac{P}{U} = \dfrac{2200}{220} = \SI{10}{A}$.
\item On choisit le fusible de \SI{16}{A} : c'est le calibre immédiatement supérieur à \SI{10}{A}. Un fusible de \SI{10}{A} serait utilisé exactement à sa limite et risquerait de fondre à la moindre pointe de courant.
\item Avec un fusible de \SI{6}{A} : le courant normal (\SI{10}{A}) dépasse le calibre, le fusible fond dès la mise en route ; le chauffe-eau ne peut pas fonctionner. Avec un fusible de \SI{32}{A} : le chauffe-eau fonctionne, mais en cas de défaut (court-circuit partiel, échauffement), le fusible ne fond pas assez tôt ; la ligne peut surchauffer et provoquer un incendie. La protection n'est plus assurée.
\end{enumerate}
\end{corrige}

\begin{corrige}[Exercice 6 — La ligne de la cuisine]
\begin{enumerate}
\item Bouilloire : $I_1 = \dfrac{2000}{220} \approx \SI{9,1}{A}$. Fer à repasser : $I_2 = \dfrac{1000}{220} \approx \SI{4,5}{A}$. Téléviseur : $I_3 = \dfrac{150}{220} \approx \SI{0,7}{A}$.
\item Les appareils sont branchés en dérivation, les intensités s'ajoutent :
\[ I = I_1 + I_2 + I_3 \approx 9,1 + 4,5 + 0,7 \approx \SI{14,3}{A}. \]
\item $I \approx \SI{14,3}{A} < \SI{16}{A}$ : le disjoncteur \textbf{ne déclenche pas}, mais la ligne fonctionne déjà à près de $90\,\%$ de sa limite. Il suffirait d'ajouter un appareil de plus de \SI{400}{W} environ pour dépasser le calibre et provoquer le déclenchement.
\item Solutions : brancher la bouilloire sur une autre prise (un autre circuit) ; ne pas faire fonctionner la bouilloire et le fer en même temps ; utiliser une ligne de calibre adapté posée par un électricien.
\end{enumerate}
\end{corrige}

\begin{corrige}[Exercice 7 — Lecture d'une courbe de danger]
\begin{enumerate}
\item Oui : $I = \SI{50}{mA} \ge \SI{30}{mA}$ et $t = \SI{0,5}{s} > \SI{0,1}{s}$ : on est en zone 4, avec risque de fibrillation cardiaque et danger de mort.
\item Le différentiel de \SI{30}{mA} coupe le circuit dès qu'une fuite atteint \SI{30}{mA}, en une fraction de seconde (moins de \SI{0,1}{s}) : on reste donc en dessous des conditions de la zone 4.
\item Non : un différentiel de \SI{300}{mA} ne déclenche qu'à partir de \SI{300}{mA}, valeur très supérieure au seuil de fibrillation (\SI{30}{mA}). Il protège l'installation contre les incendies dus aux fuites, mais pas les personnes.
\end{enumerate}
\end{corrige}

\begin{corrige}[Exercice 8 — Chasse aux erreurs dans une installation]
\begin{enumerate}
\item Dangers :
\begin{itemize}
\item Fil de cuivre à la place du fusible : plus aucune protection contre les surintensités, risque d'échauffement et d'incendie.
\item Prise sans terre : en cas de défaut d'isolement, la carcasse d'un appareil reste sous tension ; risque d'électrisation, aggravé par l'humidité.
\item Multiprises en cascade : surcharge de la prise et mauvais contacts qui chauffent ; risque d'incendie.
\item Radiateur près de la baignoire : risque de chute dans l'eau et d'électrocution, car l'eau rend le corps très conducteur.
\end{itemize}
\item Corrections : poser un fusible (ou disjoncteur) du bon calibre ; relier la prise au fil de terre ; supprimer les cascades de multiprises et répartir les appareils ; éloigner tout appareil électrique de la baignoire et ne jamais l'utiliser les pieds mouillés.
\item Consigne : « Dans cette maison, on ne touche jamais un appareil électrique avec les mains mouillées, on ne remplace jamais un fusible par un fil, et l'eau et l'électricité ne se rencontrent jamais. »
\end{enumerate}
\end{corrige}

\begin{corrige}[Exercice 9 — Sujet type BEPC : l'atelier de menuiserie]
\begin{enumerate}
\item \cle{Électrisation} : passage du courant électrique à travers le corps humain. \cle{Électrocution} : électrisation qui entraîne la mort de la victime.
\item Puissance totale :
\[ P = P_1 + P_2 + 4 \times 100 = 1500 + 800 + 400 = \SI{2700}{W}. \]
Intensité totale :
\[ I = \frac{P}{U} = \frac{2700}{220} \approx \SI{12,3}{A}. \]
\item Oui : $I \approx \SI{12,3}{A} < \SI{20}{A}$. Le calibre est supérieur à l'intensité normale, donc le disjoncteur ne déclenche pas en fonctionnement, tout en restant assez proche pour couper en cas de défaut.
\item \begin{enumerate}
\item Ce défaut est un \cle{court-circuit}.
\item La résistance du circuit devient très faible, donc l'intensité augmente brutalement et devient très grande (surintensité) : les fils chauffent fortement, avec risque d'incendie. Le disjoncteur détecte cette surintensité et \textbf{coupe automatiquement} le circuit, protégeant l'installation.
\end{enumerate}
\item Si un fil dénudé touche la carcasse, le courant de défaut s'écoule vers le sol par le fil de terre au lieu de traverser le corps de l'ouvrier. Le courant qui « part » par la phase ne « revient » plus entièrement par le neutre : le différentiel détecte cette différence dès \SI{30}{mA} et coupe le circuit en moins de \SI{0,1}{s}, avant que le courant ne devienne mortel.
\end{enumerate}
\textbf{Barème indicatif (sur 10)} : définitions (2 pts) ; calcul de $P$ et de $I$ avec formule posée et unités (2 pts) ; justification du calibre par comparaison (1 pt) ; nom du défaut et explication du rôle du disjoncteur (2 pts) ; explication terre + différentiel (3 pts).\par
\textbf{Points de vigilance} : écrire la formule $I = P/U$ avant l'application numérique ; ne pas oublier les unités (\si{A}, \si{W}) ; pour la question 5, citer explicitement la comparaison entre le courant qui part et le courant qui revient, condition de fonctionnement du différentiel.
\end{corrige}

\begin{corrige}[Exercice 10 — Sujet type BEPC : accident domestique et secours]
\begin{enumerate}
\item Trois facteurs aggravants : l'intensité du courant, la durée du contact, l'humidité (qui diminue la résistance du corps). Ici, les trois sont présents : $I = \SI{80}{mA}$ (élevée), $t = \SI{2}{s}$ (longue), rallonge humide.
\item $I = \SI{80}{mA}$ est très supérieur au seuil de tétanisation (\SI{10}{mA}) : les muscles se contractent, l'enfant ne peut plus lâcher le fil. De plus, $I > \SI{30}{mA}$ pendant $t = \SI{2}{s} \gg \SI{0,1}{s}$ : le cœur peut entrer en fibrillation, ce qui peut entraîner la mort.
\item Dans l'ordre : (1) \textbf{couper le courant} (arrêter le groupe électrogène ou ouvrir le disjoncteur général) ; (2) \textbf{dégager la victime} sans contact direct, avec un objet isolant sec (bâton en bois sec) si nécessaire ; (3) \textbf{alerter les secours} et, si l'enfant ne respire plus, commencer les gestes de premiers secours.
\item Un \textbf{interrupteur différentiel de \SI{30}{mA}} : il détecte la fuite de courant à travers le corps et coupe le circuit en moins de \SI{0,1}{s}. Un \textbf{disjoncteur (ou fusible) calibré} : il coupe le circuit en cas de surintensité (court-circuit ou surcharge). C'est le différentiel qui protège les personnes ; le disjoncteur protège l'installation.
\item Exemple de message : « Attention ! Une rallonge dénudée peut tuer. L'eau rend le corps très conducteur : ne touche jamais un fil ou une prise avec les mains mouillées. Remplace les câbles abîmés, garde les rallonges au sec et fais réparer ton installation par un électricien. Protégeons nos enfants : la prudence sauve des vies. »
\end{enumerate}
\textbf{Barème indicatif (sur 10)} : facteurs aggravants identifiés et reliés à la situation (2 pts) ; comparaison aux seuils \SI{10}{mA} et \SI{30}{mA} avec conclusion (2 pts) ; conduite à tenir dans l'ordre (3 pts) ; deux dispositifs et leur rôle (2 pts) ; message de sensibilisation clair (1 pt).\par
\textbf{Points de vigilance} : l'ordre des actions de secours est exigé (couper, dégager, alerter) ; ne jamais écrire qu'on touche la victime avant d'avoir coupé le courant ; bien distinguer le rôle du différentiel (personnes) de celui du disjoncteur (installation).
\end{corrige}

\begin{corrige}[Exercice 11 — Problème de synthèse : sécuriser une maison]
\begin{enumerate}
\item Puissance totale :
\[ P = 2200 + 3000 + 200 + 10 \times 60 = 2200 + 3000 + 200 + 600 = \SI{6000}{W}. \]
Intensité totale :
\[ I = \frac{P}{U} = \frac{6000}{220} \approx \SI{27,3}{A}. \]
$I \approx \SI{27,3}{A} < \SI{32}{A}$ : le disjoncteur général de \SI{32}{A} \textbf{convient}, avec une marge raisonnable pour les petits appareils.
\item Pour le chauffe-eau : $I = \dfrac{2200}{220} = \SI{10}{A}$. On choisit le fusible de \SI{16}{A}, calibre immédiatement supérieur : \SI{10}{A} serait à sa limite exacte et fondrait facilement ; \SI{20}{A} protégerait moins bien.
\item Le disjoncteur (ou le fusible) réagit aux \textbf{surintensités} (surcharges, courts-circuits) : il protège les \textbf{biens} (fils, appareils) contre l'échauffement et l'incendie. L'interrupteur différentiel réagit aux \textbf{fuites de courant} vers la terre, même faibles (\SI{30}{mA}) : il protège les \textbf{personnes} contre l'électrisation. Les deux sont complémentaires.
\item Tension de contact :
\[ U_c = R \cdot I_d = 50 \times 0,5 = \SI{25}{V}. \]
$U_c = \SI{25}{V} < \SI{50}{V}$ : en milieu sec, cette tension n'est pas dangereuse. De plus, le courant de fuite $I_d = \SI{0,5}{A} = \SI{500}{mA}$ dépasse largement \SI{30}{mA} : le différentiel coupe le circuit presque instantanément, ce qui supprime le danger. En milieu humide (salle de bain), le seuil de danger est plus bas, d'où l'importance du différentiel.
\item Règles pour la salle de bain : ne jamais utiliser d'appareil électrique (sèche-cheveux, radiateur) près de la baignoire ou les pieds dans l'eau ; ne pas toucher les interrupteurs avec les mains mouillées ; vérifier que les prises sont reliées à la terre et protégées par le différentiel de \SI{30}{mA}.
\end{enumerate}
\textbf{Barème indicatif (sur 10)} : calcul de la puissance totale et de l'intensité avec conclusion justifiée (3 pts) ; choix du calibre du chauffe-eau justifié (2 pts) ; distinction disjoncteur/différentiel (2 pts) ; calcul de $U_c$ et conclusion avec comparaison au seuil de \SI{50}{V} (2 pts) ; deux règles de sécurité (1 pt).\par
\textbf{Points de vigilance} : poser $P = U \cdot I$ et $U_c = R \cdot I_d$ avant tout calcul ; convertir correctement \SI{0,5}{A} en \SI{500}{mA} pour la comparaison avec le seuil du différentiel ; toute conclusion doit s'appuyer sur une comparaison chiffrée avec le calibre ou le seuil de danger.
\end{corrige}

\newpage

\coursec{Fiche bilan}

\begin{popfigure}
\centering
\begin{tikzpicture}[
  central/.style={circle, draw=popdark, fill=popblue, text=white, font=\bfseries\small, align=center, minimum size=2.6cm},
  branche/.style={rounded corners=4pt, draw=popdark, font=\small, align=center, inner sep=4pt, text width=3.1cm},
  lien/.style={thick, popdark}
]
\node[central, align=center] (C){Dangers du courant\\ et sécurité};
\node[branche, fill=poporangeL, align=center] (A) at (-4.6,2.6){\textbf{Corps humain}\\ conducteur ;\\ résistance $R_{c}$\\ (mille à cent mille ohms)};
\node[branche, fill=poppinkL, align=center] (B) at (0,3.4){\textbf{Effets du courant}\\ picotements, tétanisation,\\ brûlures, fibrillation\\ électrisation / électrocution};
\node[branche, fill=popgoldL, align=center] (C1) at (4.6,2.6){\textbf{Risques installation}\\ surcharge ($I$ trop grand)\\ court-circuit ($R$ nulle)};
\node[branche, fill=popgreenL, align=center] (D) at (-4.6,-2.6){\textbf{Fusible et disjoncteur}\\ coupent le circuit\\ si $I > I_{nominal}$};
\node[branche, fill=popblueL, align=center] (E) at (0,-3.4){\textbf{Prise de terre}\\ évacue le courant de fuite\\ vers le sol};
\node[branche, fill=poppurpleL, align=center] (F) at (4.6,-2.6){\textbf{Différentiel}\\ compare $I$ entrant et sortant\\ coupe en cas de fuite\\ (sensibilité \SI{30}{mA})};
\node[branche, fill=popcream, align=center] (G) at (-7.6,0){\textbf{Règles de sécurité}\\ mains sèches, pas de\\ surcharge de multiprise,\\ couper avant d'intervenir};
\draw[lien] (C) -- (A);
\draw[lien] (C) -- (B);
\draw[lien] (C) -- (C1);
\draw[lien] (C) -- (D);
\draw[lien] (C) -- (E);
\draw[lien] (C) -- (F);
\draw[lien] (C) -- (G);
\end{tikzpicture}
\legende{Carte mentale de la leçon : les dangers du courant électrique et les moyens de protection.}
\end{popfigure}

\begin{bilan}
\textbf{Définitions clés}
\begin{itemize}
  \item \cle{Électrisation} : passage du courant électrique à travers le corps humain, avec des lésions plus ou moins graves.
  \item \cle{Électrocution} : électrisation mortelle.
  \item \cle{Surcharge} : l'intensité demandée par les appareils branchés dépasse l'intensité maximale supportée par l'installation (fils qui chauffent).
  \item \cle{Court-circuit} : contact direct entre la phase et le neutre (résistance presque nulle) ; l'intensité devient très grande et peut provoquer un incendie.
  \item \cle{Fusible} : fil conducteur qui fond et coupe le circuit quand l'intensité dépasse sa valeur nominale.
  \item \cle{Disjoncteur} : appareil qui coupe automatiquement le circuit en cas de surintensité ; on peut le réarmer.
  \item \cle{Prise de terre} : conducteur reliant les carcasses métalliques des appareils au sol ; elle évacue les courants de fuite.
  \item \cle{Interrupteur différentiel} : dispositif qui détecte une fuite de courant (différence entre courant entrant et courant sortant) et coupe le circuit très vite.
\end{itemize}

\textbf{Grandeurs et seuils à connaître par cœur}
\begin{center}
\begin{tabularx}{\linewidth}{|l|X|}\hline
\rowcolor{popblueL} \textbf{Notion} & \textbf{Valeur ou relation} \\ \hline
Tension du réseau domestique & $U = \SI{220}{V}$ (courant alternatif) \\ \hline
Résistance du corps humain & de \SI{1000}{\ohm} (peau mouillée) à \SI{100000}{\ohm} (peau sèche) \\ \hline
Seuil de danger mortel & environ $I = \SI{30}{mA}$ à travers le corps \\ \hline
Sensibilité du différentiel & \SI{30}{mA} (protection des personnes) \\ \hline
Loi d'Ohm & $I = \dfrac{U}{R}$ : plus $R$ est faible, plus $I$ est grand \\ \hline
Puissance et surcharge & $P = U \cdot I$ : trop d'appareils $\Rightarrow$ $I$ trop grand \\ \hline
\end{tabularx}
\end{center}

\textbf{Méthodes express}
\begin{itemize}
  \item \emph{Expliquer un danger} : écrire $I = \dfrac{U}{R}$, identifier ce qui diminue $R$ (peau mouillée, court-circuit) ou augmente $I$, conclure sur l'effet (brûlure, incendie, fibrillation).
  \item \emph{Choisir un fusible} : calculer $I = \dfrac{P}{U}$ de l'appareil, puis prendre le calibre juste supérieur.
  \item \emph{Conduite à tenir} : couper le courant au disjoncteur général avant toute intervention ; ne jamais toucher une victime avant d'avoir coupé le courant ; appeler les secours.
\end{itemize}
\end{bilan}

\begin{aretenir}[Checklist avant le BEPC]
\begin{itemize}
  \item $\square$ Je sais définir électrisation et électrocution sans les confondre.
  \item $\square$ Je connais les effets du courant sur le corps : picotements, tétanisation, brûlures, fibrillation cardiaque.
  \item $\square$ Je sais que le corps humain est conducteur et que sa résistance diminue quand la peau est mouillée.
  \item $\square$ Je sais distinguer surcharge et court-circuit, et citer leurs conséquences (échauffement, incendie).
  \item $\square$ Je sais appliquer $I = \dfrac{U}{R}$ et $I = \dfrac{P}{U}$ avec les bonnes unités (A, V, $\Omega$, W).
  \item $\square$ Je sais expliquer le rôle du fusible (il fond) et du disjoncteur (il se déclenche et se réarme).
  \item $\square$ Je sais que le calibre du fusible doit être adapté à l'installation : ni trop faible, ni trop fort.
  \item $\square$ Je sais expliquer le rôle de la prise de terre : évacuer le courant de fuite vers le sol.
  \item $\square$ Je sais expliquer le rôle de l'interrupteur différentiel (\SI{30}{mA}) : protéger les personnes.
  \item $\square$ Je sais schématiser une installation simple avec phase, neutre, terre, fusible et différentiel.
  \item $\square$ Je connais les règles de sécurité : mains sèches, pas de multiprise surchargée, couper le courant avant d'intervenir.
  \item $\square$ Je sais la conduite à tenir face à une victime : couper le courant d'abord, ne jamais la toucher directement, alerter les secours.
\end{itemize}
\end{aretenir}

\findecours

\end{document}
